%% 本文件由 工程文件/scripts/assemble.py 生成，勿手改。
\bdchap{ch:oper}{Oper}

% ---- ch03_u01 ----
% ch03_u01 —— 原书 p.61–p.70（3.1 全部、3.2 全部、3.3.1）
\SourcePage{61}

\bdsec{3.1. 定义与基本性质}

\bdp{3.1.1} 设 $G$ 是 $\mathbb{C}$ 上的连通约化群，其上固定 Borel 子群 $B = B_G \subset G$。令 $N = [B,B]$，于是 $H = B/N$ 是嘉当群。记相应的李代数为 $\mathfrak{n} \subset \mathfrak{b} \subset \mathfrak{g}$，$\mathfrak{h} = \mathfrak{b}/\mathfrak{n}$。$\mathfrak{g}$ 带一个典范的降李代数滤链 $\mathfrak{g}^k$，使得 $\mathfrak{g}^0 = \mathfrak{b}$，$\mathfrak{g}^1 = \mathfrak{n}$；且对任意 $k > 0$，$\mathfrak{h} = \operatorname{gr}^0\mathfrak{g}$ 在 $\operatorname{gr}^k\mathfrak{g}$（相应地 $\operatorname{gr}^{-k}\mathfrak{g}$）上作用的权，都是 $k$ 个单正根（相应地单负根）之和。特别地 $\operatorname{gr}^{-1}\mathfrak{g} = \bigoplus \mathfrak{g}^{\alpha}$，其中 $\alpha$ 是单负根。令 $Z = Z_G = \operatorname{Center} G$。

\bdp{3.1.2} 设 $X$ 是任意光滑曲线（不必完备），$\mathfrak{F}_B$ 是其上的 $B$-丛。记 $\mathfrak{F}_G$ 为诱导出的 $G$-torsor，于是 $\mathfrak{F}_B \subset \mathfrak{F}_G$。有相应的扭李代数 $\mathfrak{b}_{\mathfrak{F}} := \mathfrak{b}_{\mathfrak{F}_B}$ 与 $\mathfrak{g}_{\mathfrak{F}} := \mathfrak{g}_{\mathfrak{F}_B} = \mathfrak{g}_{\mathfrak{F}_G}$，它们带李代数滤链 $\mathfrak{g}^k_{\mathfrak{F}}$。考虑联络层 $\operatorname{Conn}(\mathfrak{F}_B)$、$\operatorname{Conn}(\mathfrak{F}_G)$：它们分别是 $\mathfrak{b}_{\mathfrak{F}}\otimes\omega_X$- 与 $\mathfrak{g}_{\mathfrak{F}}\otimes\omega_X$-torsor（主齐性空间）。有显然的嵌入 $\operatorname{Conn}(\mathfrak{F}_B) \subset \operatorname{Conn}(\mathfrak{F}_G)$。它定义投影 $c: \operatorname{Conn}(\mathfrak{F}_G) \to (\mathfrak{g}/\mathfrak{b})_{\mathfrak{F}}\otimes\omega_X$，使得 $c^{-1}(0) = \operatorname{Conn}(\mathfrak{F}_B)$，且对任意 $\nabla \in \operatorname{Conn}(\mathfrak{F}_G)$ 与 $\nu \in \mathfrak{g}_{\mathfrak{F}}\otimes\omega_X$ 有 $c(\nabla+\nu) = c(\nabla)+\nu \bmod \mathfrak{b}_{\mathfrak{F}}\otimes\omega_X$。

\bdp{3.1.3} \emph{定义.} $X$ 上的 $G$-oper 是一对 $(\mathfrak{F}_B,\nabla)$，其中 $\nabla \in \Gamma(X,\operatorname{Conn}(\mathfrak{F}_G))$，且满足

\begin{enumerate}[leftmargin=2.2em]
\item[(1)] $c(\nabla) \in \operatorname{gr}^{-1}\mathfrak{g}_{\mathfrak{F}}\otimes\omega_X \subset (\mathfrak{g}/\mathfrak{b})_{\mathfrak{F}}\otimes\omega_X$；
\item[(2)] 对任意单负根 $\alpha$，$\alpha$-分量 $c(\nabla)^{\alpha} \in \Gamma(X,\mathfrak{g}^{\alpha}_{\mathfrak{F}}\otimes\omega_X)$ 在 $X$ 的任一点都不为零。
\end{enumerate}

若 $\mathfrak{g}$ 是半单李代数，则 $\mathfrak{g}$-oper 就是 $G_{ad}$-oper，其中 $G_{ad}$ 是对应于 $\mathfrak{g}$ 的伴随群。

我们通常把 $G$-oper 看作 $G$-局部系统 $(\mathfrak{F}_G,\nabla)$，再附加 oper 结构，即满足上述条件 (1)、(2) 的 $B$-旗 $\mathfrak{F}_B \subset \mathfrak{F}_G$。

\SourcePage{62}

$X$ 上的 $G$-oper 构成群胚 $\operatorname{Op}_G(X)$。当 $X'$ 在 $X$ 上平展时，群胚 $\operatorname{Op}_G(X')$ 构成 $X_{\text{\'et}}$ 上的群胚层 $\operatorname{Op}_G$。

\bdp{3.1.4} \emph{命题.} 设 $(\mathfrak{F}_B,\nabla)$ 是 $G$-oper。若 $X$ 连通，则 $\operatorname{Aut}(\mathfrak{F}_B,\nabla) = Z$。\hfill$\Box$

特别地，$\mathfrak{g}$-oper 没有对称性，即 $\operatorname{Op}_{\mathfrak{g}}(X)$ 是集合，$\operatorname{Op}_{\mathfrak{g}}$ 是集合层。

\bdp{3.1.5} \emph{命题.} 设 $X$ 完备、连通，亏格 $g > 1$。设 $(\mathfrak{F}_G,\nabla)$ 是 $X$ 上的 $G$-局部系统，且具有 oper 结构。则

\begin{enumerate}[label=(\roman*),leftmargin=2.2em]
\item $(\mathfrak{F}_G,\nabla)$ 上的 oper 结构唯一：相应的旗 $\mathfrak{F}_B \subset \mathfrak{F}_G$ 是 Harder--Narasimhan 旗；
\item $\operatorname{Aut}(\mathfrak{F}_G,\nabla) = Z$；
\item $(\mathfrak{F}_G,\nabla)$ 不能约化到非平凡抛物子群 $P \subset G$。\hfill$\Box$
\end{enumerate}

当然，ii) 由 i) 与 3.1.4 推出。

\bdp{3.1.6} \emph{例.} $GL_n$-oper 可以看作 $\mathcal{O}_X$-模 $\mathcal{E}$，带联络 $\nabla: \mathcal{E} \to \mathcal{E}\otimes\omega_X$ 与滤链 $\mathcal{E} = \mathcal{E}_n \supset \mathcal{E}_{n-1} \supset \cdots \supset \mathcal{E}_0 = 0$，使得

\begin{enumerate}[label=(\roman*),leftmargin=2.2em]
\item 层 $\operatorname{gr}_i\mathcal{E}$（$n \ge i \ge 1$）可逆；
\item $\nabla(\mathcal{E}_i) \subset \mathcal{E}_{i+1}\otimes\omega_X$，且对 $n-1 \ge i \ge 1$，由 $\nabla$ 诱导的态射 $\operatorname{gr}_i\mathcal{E} \to \operatorname{gr}_{i+1}\mathcal{E}\otimes\omega_X$ 是同构。
\end{enumerate}

$GL_n$-oper 可以如下构造。设 $\mathcal{A}$、$\mathcal{B}$ 是可逆 $\mathcal{O}_X$-模，$\partial: \mathcal{A} \to \mathcal{B}$ 是 $n$ 阶微分算子，其象征 $\sigma(\partial) \in \Gamma(X,\mathcal{B}\otimes\mathcal{A}^{\otimes(-1)}\otimes\Theta_X^{\otimes n})$ 无处为零。这里的 $\partial$ 是 $\mathcal{B}\otimes_{\mathcal{O}_X}\mathcal{D}_X\otimes_{\mathcal{O}_X}\mathcal{A}^{\otimes(-1)}$ 的截面；等价地，是 $\mathcal{O}$-线性映射 $\mathcal{B}^{\otimes(-1)} \to \mathcal{D}_X\otimes\mathcal{A}^{\otimes(-1)}$。设 $I \subset \mathcal{D}_X\otimes\mathcal{A}^{\otimes(-1)}$ 是该映射的像生成的 $\mathcal{D}_X$-子模。令 $\mathcal{E} := \mathcal{D}_X\otimes\mathcal{A}^{\otimes(-1)}/I$；记 $\mathcal{E}_i$ 为 $\mathcal{D}_X$ 上按算子阶数给出的通常滤链在 $\mathcal{E}$ 上诱导的滤链。\SourcePage{63}

于是 $\mathcal{E}$ 是 $\mathcal{D}_X$-模，即带联络的 $\mathcal{O}_X$-模；滤链 $\mathcal{E}_i$ 满足条件 (i)、(ii)。因此 $(\mathcal{E},\{\mathcal{E}_i\},\nabla)$ 是 $GL_n$-oper。这一构造给出 $GL_n$-oper 的群胚与上述数据 $\partial: \mathcal{A} \to \mathcal{B}$ 的群胚之间的等价。

逆函子 $\Phi$ 把 $(\mathcal{E},\{\mathcal{E}_i\},\nabla)$ 对应到如下微分算子：$\mathcal{A} \to \mathcal{B}$，其中 $\mathcal{A} := \mathcal{E}_1^{\otimes(-1)}$，$\mathcal{B} := \omega_X\otimes(\mathcal{E}/\mathcal{E}_{n-1})^{\otimes(-1)}$。把 $\mathcal{E}$ 看作 $\mathcal{D}_X$-模。设 $\mathcal{D}_X^{(k)} \subset \mathcal{D}_X$ 是阶数 $\le k$ 的算子构成的子层。则态射 $\mathcal{D}_X^{(n-1)}\otimes_{\mathcal{O}_X}\mathcal{E}_1 \to \mathcal{E}$ 是同构；因此复合 $\mathcal{D}_X^{(n)}\otimes_{\mathcal{O}_X}\mathcal{E}_1 \to \mathcal{E} \xrightarrow{\sim} \mathcal{D}_X^{(n-1)}\otimes_{\mathcal{O}_X}\mathcal{E}_1$ 给出正合列 $0 \to \mathcal{D}_X^{(n-1)}\otimes_{\mathcal{O}_X}\mathcal{E}_1 \to \mathcal{D}_X^{(n)}\otimes_{\mathcal{O}_X}\mathcal{E}_1 \to \omega_X^{\otimes(-n)}\otimes\mathcal{E}_1 \to 0$ 的分裂，即一个态射 $\omega_X^{\otimes(-n)}\otimes\mathcal{E}_1 \to \mathcal{D}_X^{(n)}\otimes_{\mathcal{O}_X}\mathcal{E}_1$。这等同于一个微分算子 $\partial: \mathcal{A} \to \mathcal{B}$（注意：同构 $\operatorname{gr}_i\mathcal{E} \xrightarrow{\sim} \operatorname{gr}_{i+1}\mathcal{E}\otimes\omega_X$ 诱导同构 $\mathcal{E}_1 \xrightarrow{\sim} (\mathcal{E}/\mathcal{E}_{n-1})\otimes\omega_X^{\otimes(n-1)}$，故 $\omega_X^{\otimes(-n)}\otimes\mathcal{E}_1 = \omega_X^{\otimes(-1)}\otimes(\mathcal{E}/\mathcal{E}_{n-1}) = \mathcal{B}^{\otimes(-1)}$）。

把上述函子 $\Phi$ 施于 $\mathrm{SL}_2$-oper，得到微分算子 $\partial: \mathcal{A} \to \omega_X\otimes\mathcal{A}^{\otimes(-1)}$。容易证明，这样得到 $\mathrm{SL}_2$-oper 的群胚与下述对象的群胚之间的等价：由可逆层 $\mathcal{A}$ 与 Sturm--Liouville 算子 $\partial: \mathcal{A} \to \omega_X\otimes\mathcal{A}^{\otimes(-1)}$ 构成的对 $(\mathcal{A},\partial)$，即象征 $\sigma(\partial)$ 无处为零的自伴二阶微分算子 $\partial$。注意 $\sigma(\partial)$ 诱导同构 $\omega_X^{\otimes 2}\otimes\mathcal{A} \xrightarrow{\sim} \omega_X\otimes\mathcal{A}^{\otimes(-1)}$，所以 $\mathcal{A}$ 自动是 $\omega_X^{\otimes(-1)}$ 的平方根。

若 $(\mathcal{A},\partial)$ 是 Sturm--Liouville 算子，$\mathcal{M}$ 是带同构 $\mathcal{M}^{\otimes 2} \xrightarrow{\sim} \mathcal{O}_X$ 的线丛，则 $\mathcal{M}$ 有典范联络；于是用 $\mathcal{M}$ 扭 $(\mathcal{A},\partial)$，得到 Sturm--Liouville 算子 $(\tilde{\mathcal{A}},\tilde{\partial})$，$\tilde{\mathcal{A}} = \mathcal{A}\otimes\mathcal{M}$。称 $(\mathcal{A},\partial)$ 与 $(\tilde{\mathcal{A}},\tilde{\partial})$ 等价。容易看出，自然映射 $\operatorname{Op}_{\mathrm{SL}_2}(X) \to \operatorname{Op}_{\mathfrak{sl}_2}(X)$ 把 $\operatorname{Op}_{\mathfrak{sl}_2}(X)$ 与 Sturm--Liouville 算子的等价类集合等同。

其它典型群的 oper 可用类似方式描述（局部情形已在 [DS85, 第 8 节] 中完成）。

\SourcePage{64}

\bdp{3.1.7} 把 $\mathfrak{sl}_2$-oper 与 Sturm--Liouville 算子的等价类等同（见 3.1.6），可见 $\operatorname{Op}_{\mathfrak{sl}_2}$ 是 $\omega_X^{\otimes 2}$-torsor：$\omega_X^{\otimes 2}$ 的截面 $\eta$ 把 Sturm--Liouville 算子 $\partial: \mathcal{A} \to \mathcal{A}\otimes\omega_X^{\otimes 2}$（$\mathcal{A}^{\otimes(-2)} = \omega_X$）映到 $\partial - \eta$。下面不用 Sturm--Liouville 算子，来描述 $\omega_X^{\otimes 2}$ 在 $\operatorname{Op}_{\mathfrak{sl}_2}$ 上的这个作用。

利用 $\mathfrak{sl}_2$ 上的双线性型 $\operatorname{Tr}(AB)$，把 $\mathfrak{n} \subset \mathfrak{sl}_2$ 与 $(\mathfrak{sl}_2/\mathfrak{b})^*$ 等同。若 $\mathfrak{F} = (\mathfrak{F}_B,\nabla)$ 是 $\mathfrak{sl}_2$-oper，则由 3.1.3，截面 $c(\nabla)$ 平凡化了层 $(\mathfrak{sl}_2/\mathfrak{b})_{\mathfrak{F}_B}\otimes\omega_X$。于是 $(\mathfrak{sl}_2/\mathfrak{b})_{\mathfrak{F}_B} = \omega_X^{\otimes(-1)}$，$\mathfrak{n}_{\mathfrak{F}} = \omega_X$，并有嵌入 $\omega_X^{\otimes 2} = \mathfrak{n}_{\mathfrak{F}_B}\otimes\omega_X \hookrightarrow (\mathfrak{sl}_2)_{\mathfrak{F}_B}\otimes\omega_X$。用 $\omega_X^{\otimes 2} \subset (\mathfrak{sl}_2)_{\mathfrak{F}_B}\otimes\omega_X$ 的截面 $\mu$ 平移 $\nabla$，得到新的 oper，记作 $\mathfrak{F}+\mu$。$\operatorname{Op}_{\mathfrak{sl}_2}$ 上的这个 $\omega_X^{\otimes 2}$-作用与上面引入的作用一致，因此它使 $\operatorname{Op}_{\mathfrak{sl}_2}$ 成为 $\omega_X^{\otimes 2}$-torsor。

\emph{注.} 众所周知，这个 torsor 平凡（即使 $H^1(X,\omega_X^{\otimes 2}) \ne 0$，即 $g \le 1$；$\mathbb{P}^1$ 上或椭圆曲线上确实存在 Sturm--Liouville 算子）。然而对曲线的族 $X$，这个 torsor 可能不平凡。

\bdp{3.1.8} 3.1.9 中要用如下记号。设 $B_0 \subset PSL_2$ 是上三角矩阵构成的群。令 $N_0 := [B_0,B_0]$，$\mathfrak{b}_0 := \operatorname{Lie} B_0$，$\mathfrak{n}_0 := \operatorname{Lie} N_0$。借助伴随作用 $B_0/N_0 \to \operatorname{Aut}\mathfrak{n}_0 = \mathbb{G}_m$，把 $B_0/N_0$ 与 $\mathbb{G}_m$ 等同。用矩阵 $e := \begin{pmatrix}0&1\\0&0\end{pmatrix}$ 与 $f := \begin{pmatrix}0&0\\1&0\end{pmatrix}$，把 $\mathfrak{n}_0$ 与 $\mathfrak{sl}_2/\mathfrak{b}_0$ 都等同于 $\mathbb{C}$。于是对 $\mathfrak{sl}_2$-oper $\mathfrak{F} = (\mathfrak{F}_{B_0},\nabla)$，3.1.7 中的同构 $(\mathfrak{sl}_2/\mathfrak{b}_0)_{\mathfrak{F}_{B_0}} \xrightarrow{\sim} \omega_X^{\otimes(-1)}$（或同构 $\mathfrak{n}_{\mathfrak{F}_{B_0}} \xrightarrow{\sim} \omega_X$）诱导同构：$\mathfrak{F}_{B_0}$ 沿 $B_0 \to B_0/N_0 = \mathbb{G}_m$ 的前推同构于 $\mathbb{G}_m$-torsor $\omega_X$。

\bdp{3.1.9} 对任意半单李代数 $\mathfrak{g}$，我们将给出 $\operatorname{Op}_{\mathfrak{g}}(X)$ 相当显式的描述。特别地，我们将在 $\operatorname{Op}_{\mathfrak{g}}(X)$ 上引入一个「典范的」仿射空间结构（$\mathfrak{g} = \mathfrak{sl}_2$ 时的情形已在 3.1.7 中引入）。设 $G$ 是对应于 $\mathfrak{g}$ 的伴随群，$B$ 是其 Borel 子群。沿用 3.1.8 的记号。固定主嵌入 $i: \mathfrak{sl}_2 \hookrightarrow \mathfrak{g}$，使 $i(\mathfrak{b}_0) \subset \mathfrak{b}$；有相应的嵌入 $i_G: PSL_2 \hookrightarrow G$，\SourcePage{65}$i_B: B_0 \hookrightarrow B$。

令 $V = V_{\mathfrak{g}} := \mathfrak{g}^{N_0}$。于是 $\mathfrak{n}_0 \subset V \subset \mathfrak{n}$。有 $\mathbb{G}_m = B_0/N_0$ 在 $V$ 上的伴随作用 $\operatorname{Ad}$。定义 $V$ 上新的 $\mathbb{G}_m$-作用 $a$：$a(t)v := t\operatorname{Ad}(t)v$，$v \in V$，$t \in \mathbb{G}_m$。

考虑向量丛 $V_{\omega_X}$，即 $V$ 关于 $\mathbb{G}_m$-作用 $a$ 的 $\omega_X$-扭（这里把 $\omega_X$ 看作 $X$ 上的 $\mathbb{G}_m$-torsor）。用 $\omega_X$ 扭嵌入 $\mathbb{C} \xrightarrow{\sim} \mathbb{C}e = \mathfrak{n}_0 \hookrightarrow V$，得到嵌入 $\omega_X^{\otimes 2} \hookrightarrow V_{\omega_X}$。

对任意 $\mathfrak{sl}_2$-oper $\mathfrak{F}_0 = (\mathfrak{F}_{B_0},\nabla_0)$，其 $i$-前推 $i\mathfrak{F}_0 = (\mathfrak{F}_B,\nabla)$ 是 $\mathfrak{g}$-oper。由 3.1.8 有典范同构 $V_{\omega_X} = V_{\mathfrak{F}_{B_0}}\otimes\omega_X$，从而有典范嵌入 $V_{\omega_X} \subset \mathfrak{b}_{\mathfrak{F}_0}\otimes\omega_X = \mathfrak{b}_{\mathfrak{F}_B}\otimes\omega_X$。用 $V_{\omega_X}$ 的截面 $\nu$ 平移 $\nabla$，得到新的 $\mathfrak{g}$-oper，记作 $i\mathfrak{F}_0+\nu$。

令 $\underline{\operatorname{Op}}_{\mathfrak{g}}$ 为 $\omega_X^{\otimes 2}$-torsor $\operatorname{Op}_{\mathfrak{sl}_2}$ 经嵌入 $\omega_X^{\otimes 2} \subset V_{\omega_X}$ 诱导出的 $V_{\omega_X}$-torsor。$\underline{\operatorname{Op}}_{\mathfrak{g}}$ 的截面是一对如上所述的 $(\mathfrak{F}_0,\nu)$；并约定：对 $\omega_X^{\otimes 2}$ 的截面 $\mu$，有 $(\mathfrak{F}_0+\mu,\nu) = (\mathfrak{F}_0,\mu+\nu)$。有典范映射
\begin{equation}\tag{43}
\underline{\operatorname{Op}}_{\mathfrak{g}} \longrightarrow \operatorname{Op}_{\mathfrak{g}},
\end{equation}
它把 $(\mathfrak{F}_0,\nu)$ 映到 $i\mathfrak{F}_0+\nu$。

\bdp{3.1.10} \emph{命题.} 映射（43）是双射。\hfill$\Box$

\emph{注.}
\begin{enumerate}[label=(\roman*),leftmargin=2.2em]
\item 尽管双射（43）是典范的，我们不确定它给出 $\operatorname{Op}_{\mathfrak{g}}$ 的合理描述。
\item 3.1.9 中的空间 $V = V_{\mathfrak{g}}$ 依赖于主嵌入 $i: \mathfrak{sl}_2 \hookrightarrow \mathfrak{g}$ 的选取（对这种 $i$，存在唯一的 Borel 子代数 $\mathfrak{b} \subset \mathfrak{g}$ 包含 $i(\mathfrak{b}_0)$）。但任意两个主嵌入 $\mathfrak{sl}_2 \hookrightarrow \mathfrak{g}$ 都由 $G = G_{ad}$ 的唯一元素共轭。于是可以等同对应于不同 $i$ 的各个 $V$，得到与 $\mathfrak{g}$ 典范相伴的向量空间（不是 $\mathfrak{g}$ 的子空间！）。
\item 设 $G$ 是对应于 $\mathfrak{g}$ 的伴随群，$B$ 是 $G$ 的 Borel 子群。命题 3.1.10 表明，对任意 $\mathfrak{g}$-oper $\mathfrak{F} = (\mathfrak{F}_B,\nabla)$，$\mathfrak{F}_B$ 同构于某个典范 $B$-丛 $\mathfrak{F}_B^0$，它不依赖于 $\mathfrak{F}$。\SourcePage{66}事实上，$\mathfrak{F}_B^0$ 是 2.6.5 中典范 $(\operatorname{Aut}^0 O)$-丛沿某个同态 $i_B\sigma\pi: \operatorname{Aut}^0 O \to B$ 的前推。这里 $\pi$ 是投影 $\operatorname{Aut}^0 O \to \operatorname{Aut}(O/\mathfrak{m}^3)$，$\mathfrak{m}$ 是 $O$ 的极大理想；$\sigma$ 是同构 $\operatorname{Aut}(O/\mathfrak{m}^3) \xrightarrow{\sim} B_0$，$B_0$ 是 $PSL_2$ 的 Borel 子群；$i_B: B_0 \to B$ 由主嵌入 $PSL_2 \to G$ 诱导（$\sigma$ 与 $i_B$ 在唯一共轭的意义下唯一）。
\end{enumerate}

\bdp{3.1.11} 设 $X$ 完备。则 $G$-oper 构成光滑代数栈，仍记作 $\operatorname{Op}_G(X)$（记号滥用）。若 $G$ 半单，这是 Deligne--Mumford 栈（见 3.1.4）；若 $G$ 伴随，则 $\operatorname{Op}_G(X) = \operatorname{Op}_{\mathfrak{g}}(X)$ 是概形，经（43）同构于仿射空间 $\underline{\operatorname{Op}}_{\mathfrak{g}}(X)$。

\emph{注.}
\begin{enumerate}[label=(\roman*),leftmargin=2.2em]
\item 若 $X$ 非完备，则 $\operatorname{Op}_{\mathfrak{g}}(X)$ 是归纳概形。
\item 若 $X$ 完备、连通，亏格 $g > 1$，则 $\dim\operatorname{Op}_{\mathfrak{g}}(X) = (g-1)\cdot\dim\mathfrak{g}$。事实上，由命题 3.1.10，$\dim\operatorname{Op}_{\mathfrak{g}}(X) = \dim\underline{\operatorname{Op}}_{\mathfrak{g}}(X) = \dim\Gamma(X,V_{\omega_X})$；Hitchin 的一个简单计算（见 2.2.4 的注 4）表明，当 $g > 1$ 时 $\dim\Gamma(X,V_{\omega_X}) = (g-1)\cdot\dim\mathfrak{g}$。实际上，3.1.13 将表明 $\Gamma(X,V_{\omega_X}) = \operatorname{Hitch}_{{}^{L}\mathfrak{g}}(X)$，所以直接使用 2.2.4(ii) 中提到的 Hitchin 公式即可：
\[
\dim\operatorname{Hitch}_{{}^{L}\mathfrak{g}}(X) = (g-1)\cdot\dim{}^{L}\mathfrak{g} = (g-1)\cdot\dim\mathfrak{g}.
\]
\item 设 $X$ 如注 ii，$G$ 是对应于 $\mathfrak{g}$ 的伴随群。有显然的态射 $i: \operatorname{Op}_{\mathfrak{g}}(X) \to \operatorname{LocSys}_G$，其中 $\operatorname{LocSys}_G$ 是 $X$ 上 $G$-局部系统的栈。可以证明：不能约化到非平凡抛物子群 $P \subset G$ 且没有非平凡自同构的 $G$-局部系统，构成开子栈 $U \subset \operatorname{LocSys}_G$。它实际上是光滑簇；$U$ 带典范辛结构。由 3.1.5，\SourcePage{67}$i(\operatorname{Op}_{\mathfrak{g}}(X)) \subset U$，且 $i$ 是集合论意义下的嵌入。事实上 $i$ 是闭嵌入，$i(\operatorname{Op}_{\mathfrak{g}}(X))$ 是 $U$ 的 Lagrange 子簇。此外，$i(\operatorname{Op}_{\mathfrak{g}}(X)) = \pi^{-1}(S)$，其中 $\pi: \operatorname{LocSys}_G \to \operatorname{Bun}_G$ 对应于忘掉联络，$S \subset \operatorname{Bun}_G$ 是同构于 $\mathfrak{F}_G^0$ 的 $G$-丛构成的局部闭子栈；$\mathfrak{F}_G^0$ 是对应于 3.1.10 注 iii 中引入的 $B$-丛 $\mathfrak{F}_B^0$ 的 $G$-丛（于是 $S$ 是幂幺群 $\operatorname{Aut}\mathfrak{F}_G^0$ 的分类栈）。
\end{enumerate}

\bdp{3.1.12} 记 $A_{\mathfrak{g}}(X)$ 为 $\operatorname{Op}_{\mathfrak{g}}(X)$ 的坐标环。我们将构造 $A_{\mathfrak{g}}(X)$ 上的典范滤链，以及分次代数的典范同构
\begin{equation}\tag{44}
\sigma_{A(X)}: \operatorname{gr} A_{\mathfrak{g}}(X) \xrightarrow{\sim} \mathfrak{z}^{cl}_{{}^{L}\mathfrak{g}}(X),
\end{equation}
其中 ${}^{L}\mathfrak{g}$ 表示 $\mathfrak{g}$ 的朗兰兹对偶，（44）的右端已在 2.2.2 中定义。我们给出两种等价的构造：3.1.13 中的构造直截了当，用到同构（43）；3.1.14 中的构造更自然。

\bdp{3.1.13} 利用 3.1.8，把 $A_{\mathfrak{g}}(X)$ 与 $\underline{\operatorname{Op}}_{\mathfrak{g}}(X)$ 的坐标环等同。记 $A^{cl}_{\mathfrak{g}}(X)$ 为向量空间 $\Gamma(X,V_{\omega_X})$ 的坐标环，它对应于仿射空间 $\underline{\operatorname{Op}}_{\mathfrak{g}}(X)$。考虑 $A^{cl}_{\mathfrak{g}}(X)$ 上与 $V$ 上 $\mathbb{G}_m$-作用 $a$ 所诱导者相反的作用（见 3.1.7）；$A^{cl}_{\mathfrak{g}}(X)$ 上相应的分次为正。它诱导 $A_{\mathfrak{g}}(X)$ 上的典范环滤链，以及典范同构 $\operatorname{gr} A_{\mathfrak{g}}(X) \xrightarrow{\sim} A^{cl}_{\mathfrak{g}}(X)$。

于是为定义（44），只剩构造分次同构 $A^{cl}_{\mathfrak{g}}(X) \xrightarrow{\sim} \mathfrak{z}^{cl}_{{}^{L}\mathfrak{g}}(X)$；这等价于构造概形的 $\mathbb{G}_m$-等变同构 $\Gamma(X,V_{\omega_X}) \xrightarrow{\sim} \operatorname{Hitch}_{{}^{L}\mathfrak{g}}(X)$。由 2.2.2，$\operatorname{Hitch}_{{}^{L}\mathfrak{g}}(X) := \Gamma(X,C_{\omega_X})$，$C := C_{{}^{L}\mathfrak{g}}$。因此只需构造概形的 $\mathbb{G}_m$-等变同构 $V_{\mathfrak{g}} \xrightarrow{\sim} C_{{}^{L}\mathfrak{g}}$。（$V_{\mathfrak{g}}$ 带 3.1.7 中的作用 $a$。）
% [Tutor Note: 原书 p.67 此处与本节末「见 3.1.7」所指的 G_m-作用 a，实际定义于 3.1.9，底本如此，保留原样。]

\SourcePage{68}

由 2.2.1，$C_{{}^{L}\mathfrak{g}} = \operatorname{Spec}(\operatorname{Sym}{}^{L}\mathfrak{g})^{{}^{L}G}$，其中 $G$ 是对应于 $\mathfrak{g}$ 的连通群。可以把 $(\operatorname{Sym}{}^{L}\mathfrak{g})^{{}^{L}G}$ 与 $(\operatorname{Sym}\mathfrak{g}^*)^G$ 等同，因为这两个分次代数都典范同构于 $(\operatorname{Sym}\mathfrak{h}^*)^W$，其中 $W$ 是 Weyl 群。故 $C_{{}^{L}\mathfrak{g}} = C'_{\mathfrak{g}}$，其中
\begin{equation}\tag{45}
C'_{\mathfrak{g}} = \operatorname{Spec}(\operatorname{Sym}\mathfrak{g}^*)^G,
\end{equation}
即 $C'_{\mathfrak{g}}$ 是 $\mathfrak{g}$ 关于 $G$ 伴随作用的仿射概形商。最后，由 Kostant 的工作 [Ko63] 中的定理 0.10，有典范同构 $V_{\mathfrak{g}} \xrightarrow{\sim} C'_{\mathfrak{g}}$，它把 $v \in V_{\mathfrak{g}}$ 映到 $v + i\left(\begin{pmatrix}0&0\\1&0\end{pmatrix}\right) \in \mathfrak{g}$ 在 $C'_{\mathfrak{g}}$ 中的像。该同构与 $\mathbb{G}_m$-作用交换。

\bdp{3.1.14} 描述 $A_{\mathfrak{g}}(X)$ 上的典范滤链与同构（44），还有更自然的方式。

有标准做法把滤 $\mathbb{C}$-代数与分次平坦 $\mathbb{C}[\hbar]$-代数等同（这里 $\deg\hbar = 1$）。具体说，带递增滤链 $\{A_i\}$ 的代数 $A$ 对应分次 $\mathbb{C}[\hbar]$-代数 $A^{\sim} = \bigoplus A_i$，乘以 $\hbar$ 即嵌入 $A_i \hookrightarrow A_{i+1}$。注意 $A = A^{\sim}/(\hbar-1)A^{\sim}$，$\operatorname{gr} A = A^{\sim}/\hbar A^{\sim}$。取谱可知，$\operatorname{Spec} A^{\sim}$ 是直线 $\mathbb{A}^1 = \operatorname{Spec}\mathbb{C}[\hbar]$ 上的平坦仿射概形；$A^{\sim}$ 上的分次等同于 $\operatorname{Spec} A^{\sim}$ 上与 $\mathbb{A}^1$ 的位似作用相容的 $\mathbb{G}_m$-作用。我们要构造概形 $\operatorname{Spec} A_{\mathfrak{g}}(X)^{\sim}$。

设 $\mathfrak{F}$ 是 $X$ 上的 $G$-torsor。记 $\mathcal{E}_{\mathfrak{F}}$ 为 $\mathfrak{F}$ 的无穷小对称构成的李代数胚；有典范正合列
\[
0 \to \mathfrak{g}_{\mathfrak{F}} \to \mathcal{E}_{\mathfrak{F}} \xrightarrow{\pi} \Theta_X \to 0.
\]

回顾：对 $\hbar \in \mathbb{C}$，$\mathfrak{F}$ 上的 $\hbar$-联络是 $\mathcal{O}_X$-线性映射 $\nabla_{\hbar}: \Theta_X \to \mathcal{E}_{\mathfrak{F}}$，满足 $\pi\nabla_{\hbar} = \hbar\operatorname{id}_{\Theta_X}$（通常的联络对应 $\hbar = 1$）。如 3.1.3 那样定义 $G$-$\hbar$-oper，只需把联络 $\nabla$ 换成 $\hbar$-联络 $\nabla_{\hbar}$。关于 $G$-oper 的上述结果都可移用到 $G$-$\hbar$-oper。特别地，$\mathfrak{g}$-$\hbar$-oper（即伴随群的 $\hbar$-oper）构成仿射概形\SourcePage{69}$\operatorname{Op}_{\mathfrak{g},\hbar}(X)$。

对 $\lambda \in \mathbb{C}^*$，有概形同构
\begin{equation}\tag{46}
\operatorname{Op}_{\mathfrak{g},\hbar}(X) \xrightarrow{\sim} \operatorname{Op}_{\mathfrak{g},\lambda\hbar}(X),
\end{equation}
由 $(\mathfrak{F}_B,\nabla_{\hbar}) \mapsto (\mathfrak{F}_B,\lambda\nabla_{\hbar})$ 定义。$\hbar$ 变化时，$\operatorname{Op}_{\mathfrak{g},\hbar}(X)$ 成为仿射 $\mathbb{C}[\hbar]$-概形 $\operatorname{Op}_{\mathfrak{g}}(X)^{\sim} = \operatorname{Spec} A_{\mathfrak{g}}(X)^{\sim}$ 的纤维。对 $\mathfrak{g}$-$\hbar$-oper 使用 3.1.9--3.1.10 的类比，可以证明 $A_{\mathfrak{g}}(X)^{\sim}$ 在 $\mathbb{C}[\hbar]$ 上平坦。态射（46）定义了 $\mathbb{G}_m$ 在 $\operatorname{Op}_{\mathfrak{g}}(X)^{\sim}$ 上的作用，即 $A_{\mathfrak{g}}(X)^{\sim}$ 的分次。$A_{\mathfrak{g}}(X) = A_{\mathfrak{g}}(X)^{\sim}/(\hbar-1)A_{\mathfrak{g}}(X)^{\sim}$ 上相应的滤链与 3.1.13 的滤链一致。

构造（44）等同于构造 $\operatorname{Op}_{\mathfrak{g},0}(X) = \operatorname{Spec}\operatorname{gr} A_{\mathfrak{g}}(X)$ 与 $\operatorname{Hitch}_{{}^{L}\mathfrak{g}}(X) = \operatorname{Spec}\mathfrak{z}^{cl}_{{}^{L}\mathfrak{g}}(X)$ 之间的 $\mathbb{G}_m$-等变同构。如 3.1.11 所述，$\operatorname{Hitch}_{{}^{L}\mathfrak{g}}(X) = \Gamma(X,C'_{\omega_X})$，其中 $C' = C'_{\mathfrak{g}}$ 由（45）定义。有层的典范映射
\begin{equation}\tag{47}
\operatorname{Op}_{\mathfrak{g},0} \longrightarrow C'_{\omega_X},
\end{equation}
它把 $(\mathfrak{F}_B,\nabla_0)$ 映到 $\nabla_0 \in \mathfrak{g}_{\mathfrak{F}}\otimes\omega_X$ 经投影 $\mathfrak{g} \to C'$ 的像。Kostant 的工作 [Ko63] 中的定理 0.10 与命题 19 表明（47）是双射，它诱导所要的同构 $\operatorname{Op}_{\mathfrak{g},0}(X) \xrightarrow{\sim} \Gamma(X,C'_{\omega_X})$。

\bdsec{3.2. 局部 oper 与 Feigin--Frenkel 同构}

\bdp{3.2.1} 把 $X$ 换成形式圆盘 $\operatorname{Spec} O$，$O \simeq \mathbb{C}[[t]]$。3.1 的构造与结果容易移用到这一情形。$\operatorname{Spec} O$ 上的 $\mathfrak{g}$-oper 构成概形 $\operatorname{Op}_{\mathfrak{g}}(O)$，它同构于可数多个变量的多项式环的谱。更确切地说，同构（43）把 $\operatorname{Op}_{\mathfrak{g}}(O)$ 与向量空间 $H^0(\operatorname{Spec} O,V_{\omega_O})$（$V := V_{\mathfrak{g}}$）对应的仿射空间等同。$\operatorname{Spec} O$ 上的 $G$-oper 构成代数栈 $\operatorname{Op}_G(O)$，它同构于 $\operatorname{Op}_{\mathfrak{g}}(O) \times B(Z)$，其中 $B(Z)$ 是中心 $Z \subset G$ 的分类栈，$\mathfrak{g} := \operatorname{Lie}(G/Z)$（该同构不完全是典范的；$\operatorname{Op}_G(O)$ 的典范描述见式（58））。

\SourcePage{70}

与整体情形一样（见 3.1.12--3.1.14），$\operatorname{Op}_{\mathfrak{g}}(O)$ 的坐标环 $A_{\mathfrak{g}}(O)$ 带典范滤链，并有典范同构
\begin{equation}\tag{48}
\sigma_A: \operatorname{gr} A_{\mathfrak{g}}(O) \xrightarrow{\sim} \mathfrak{z}^{cl}_{{}^{L}\mathfrak{g}}(O)
\end{equation}
（见（44））。注意 $\operatorname{Aut} O$ 以显然的方式作用在上述所有对象上。因此 $A_{\mathfrak{g}}(O)$ 是滤 $\operatorname{Aut} O$-代数，$\sigma_A$ 是分次 $\operatorname{Aut} O$-代数的同构。

\bdp{3.2.2} \emph{定理.}（[FF92]）存在滤 $\operatorname{Aut} O$-代数的典范同构
\begin{equation}\tag{49}
\varphi_O: A_{\mathfrak{g}}(O) \xrightarrow{\sim} \mathfrak{z}_{{}^{L}\mathfrak{g}}(O),
\end{equation}
使得 $\sigma_{\mathfrak{z}}\operatorname{gr}\varphi_O = \sigma_A$，其中 $\sigma_{\mathfrak{z}}: \operatorname{gr}\mathfrak{z}_{{}^{L}\mathfrak{g}}(O) \to \mathfrak{z}^{cl}_{{}^{L}\mathfrak{g}}(O)$ 是象征映射。\hfill$\Box$

\emph{注.}
\begin{enumerate}[label=(\roman*),leftmargin=2.2em]
\item 该同构由某些附加的相容性唯一确定；见 3.6.7。
\item Feigin 与 Frenkel 原来的构造是表示论的，对我们来说完全神秘。另一种几何构造见 ???；两种构造的比较见 ???。
\item 对 $\mathfrak{g} = \mathfrak{sl}_2$，（49）有简单的显式描述，本质上属于 Sugawara；见 ???。
\end{enumerate}
% [Tutor Note: 原书 p.70 上述三处交叉引用印作「???」，是 LaTeX 未解析的引用；底本如此，保留原样。]

\bdsec{3.3. 整体版本}

\bdp{3.3.1} 回到整体情形，即 $X$ 是完备曲线。我们将构造代数 $A_{\mathfrak{g}}(X)$ 与 $\mathfrak{z}_{{}^{L}\mathfrak{g}}(X)$ 之间的典范同构（后者由 2.7.4 的公式（27）定义）。

% ---- ch03_u02 ----
% ch03_u02 —— 原书 p.71–p.80（3.3.1 末至 3.5 开头）
\SourcePage{71}

取 $x\in X$。整体 $\mathfrak{g}$-oper 到 $\operatorname{Spec}\mathcal{O}_{x}$ 的限制，定义了仿射概形之间的态射
\[
  \operatorname{Op}_{\mathfrak{g}}(X)\longrightarrow\operatorname{Op}_{\mathfrak{g}}(\mathcal{O}_{x}).
\]
这是闭嵌入，于是得到坐标环之间的满射
\begin{equation}\tag{50}
  \theta_{x}^{A} : A_{\mathfrak{g}}(\mathcal{O}_{x})\twoheadrightarrow A_{\mathfrak{g}}(X).
\end{equation}
$\theta_{x}^{A}$ 与典范滤链严格相容（要看出这一点，可用同构 (24)）。

\bdp{3.3.2} \emph{定理.} 存在唯一的滤代数同构
\begin{equation}\tag{51}
  \varphi_{X} : A_{\mathfrak{g}}(X)\xrightarrow{\ \sim\ }\mathfrak{z}_{{}^{L}\mathfrak{g}}(X),
\end{equation}
使得对任意 $x\in X$，图表
\[
\begin{tikzcd}
  A_{\mathfrak{g}}(\mathcal{O}_{x}) \arrow[r, "\theta_{x}^{A}"] \arrow[d, "\varphi_{\mathcal{O}_{x}}\,\wr"']
    & A_{\mathfrak{g}}(X) \arrow[d, "\wr\,\varphi_{X}"] \\
  \mathfrak{z}_{{}^{L}\mathfrak{g}}(\mathcal{O}_{x}) \arrow[r, "\theta_{x}^{\mathfrak{z}}"']
    & \mathfrak{z}_{{}^{L}\mathfrak{g}}(X)
\end{tikzcd}
\]
交换（这里 $\varphi_{\mathcal{O}_{x}}$ 是 $O=\mathcal{O}_{x}$ 时的同构 (49)）。有 $\sigma_{\mathfrak{z}(X)}\cdot\operatorname{gr}\varphi_{X}=\sigma_{A(X)}$，其中 $\sigma_{A(X)}$ 是同构 (44)，而 $\sigma_{\mathfrak{z}(X)}:\operatorname{gr}\mathfrak{z}(X)\to\mathfrak{z}^{cl}(X)$ 的定义见 2.7.4 末。

\emph{证明.} 由于 $\theta_{x}^{A}$ 与 $\theta_{x}^{\mathfrak{z}}$ 都是满射，且都与滤链严格相容，只需证明存在使图表交换的同构 $\varphi_{X}$。按 2.6.5，有 $\mathcal{D}_{X}$-代数 $\mathcal{A}_{\mathfrak{g}}:=A_{\mathfrak{g}}(\mathcal{O})_{X}$，其纤维为 $A_{\mathfrak{g}}(\mathcal{O}_{x})$。任一整体 oper $\mathfrak{F}\in\operatorname{Op}_{\mathfrak{g}}(X)$ 定义截面 $\gamma_{\mathfrak{F}}:X\to\operatorname{Spec}\mathcal{A}_{\mathfrak{g}}$，其中 $\gamma_{\mathfrak{F}}(x)$ 是 $\mathfrak{F}$ 到 $\operatorname{Spec}\mathcal{O}_{x}$ 的限制。截面 $\gamma_{\mathfrak{F}}$ 是水平的，由此得到 $\operatorname{Op}_{\mathfrak{g}}(X)$ 与 $\operatorname{Spec}\mathcal{A}_{\mathfrak{g}}$ 的水平截面概形之间的同构（认为这需要证明的读者

\SourcePage{72}
可以在 3.3.3 中找到它）。过渡到坐标环，得到典范同构
\begin{equation}\tag{52}
  A_{\mathfrak{g}}(X)\xrightarrow{\ \sim\ }H_{\nabla}(X,\mathcal{A}_{\mathfrak{g}})
\end{equation}
（$H_{\nabla}$ 的定义见 2.6.2）。另一方面，(49) 给出 $\mathcal{D}_{X}$-代数的同构
\[
  \varphi : \mathcal{A}_{\mathfrak{g}}\xrightarrow{\ \sim\ }\mathfrak{z}_{{}^{L}\mathfrak{g}},
\]
从而有同构
\begin{equation}\tag{53}
  H_{\nabla}(X,\mathcal{A}_{\mathfrak{g}})\xrightarrow{\ \sim\ }H_{\nabla}(X,\mathfrak{z}_{{}^{L}\mathfrak{g}})
  =\mathfrak{z}_{{}^{L}\mathfrak{g}}(X).
\end{equation}
现在 $\varphi_{X}$ 就是 (52) 与 (53) 的复合。\hfill$\Box$

\bdp{3.3.3} 本小节（读者完全可以跳过）证明：$\mathfrak{g}$-oper 可以等同于 $\operatorname{Spec}\mathcal{A}_{\mathfrak{g}}$ 的水平截面（3.3.2 用到了这一等同）。

记 $\mathfrak{g}^{+}$ 为所有满足下述条件的 $a\in\mathfrak{g}^{-1}$ 的集合：$a$ 在 $\mathfrak{g}^{\alpha}$ 中的像，对每个单负根 $\alpha$ 都非零（记号见 3.1.1）。$\mathfrak{g}^{+}$ 是仿射概形。考虑 $\operatorname{Aut}^{0}\mathcal{O}$ 通过标准特征标 $\operatorname{Aut}^{0}\mathcal{O}\to\operatorname{Aut}(t\mathcal{O}/t^{2}\mathcal{O})=\mathbb{G}_{m}$ 在 $\mathfrak{g}^{+}$ 上的作用。记 $B$ 为 $\mathfrak{g}$ 所对应的伴随群的 Borel 子群，并给 $B$ 配备 $\operatorname{Aut}^{0}\mathcal{O}$ 的平凡作用。应用 2.6.7 中的函子 $\mathcal{J}:\{\operatorname{Aut}^{0}\mathcal{O}\text{-概形}\}\to\{\operatorname{Aut}\mathcal{O}\text{-概形}\}$，得到：$\mathcal{J}B$ 是态射 $\operatorname{Spec}\mathcal{O}\to B$ 的概形，$\mathcal{J}\mathfrak{g}^{+}$ 是 $\operatorname{Spec}\mathcal{O}$ 上 $\mathfrak{g}^{+}$-值微分形式的概形。群 $\mathcal{J}B$ 通过规范变换作用在 $\mathcal{J}\mathfrak{g}^{+}$ 上，而 $\operatorname{Op}_{\mathfrak{g}}(\mathcal{O})$ 是商概形。$\mathcal{J}B$ 在 $\mathcal{J}\mathfrak{g}^{+}$ 上的作用与态射 $\mathcal{J}\mathfrak{g}^{+}\to\operatorname{Op}_{\mathfrak{g}}(\mathcal{O})$ 都是 $\operatorname{Aut}\mathcal{O}$-等变的。事实上，$\mathcal{J}\mathfrak{g}^{+}$ 是 $\operatorname{Op}_{\mathfrak{g}}(\mathcal{O})$ 上的 $\mathcal{J}B$-torsor（主齐性空间）。此外，选取 $\eta\in\omega_{\mathcal{O}}^{+}:=\omega_{\mathcal{O}}\setminus t\omega_{\mathcal{O}}$ 定义了它的一个截面 $S_{\eta}\subset\mathcal{J}\mathfrak{g}^{+}$：
\[
  S_{\eta}:=\eta\cdot i(f)+V\otimes\omega_{\mathcal{O}}
\]
（这里 $f:=\begin{pmatrix}0&0\\1&0\end{pmatrix}$，而 $i,V$ 已在 3.1.9 中定义）。$S_{\eta}$ 是截面这一事实，正是命题 3.1.10 的局部形式。截面 $S_{\eta}$ 定义了诱导 torsor $\mathfrak{g}^{+}\times\omega_{\mathcal{O}}^{+}\to\operatorname{Op}_{\mathfrak{g}}(\mathcal{O})\times\omega_{\mathcal{O}}^{+}$ 的 $\operatorname{Aut}\mathcal{O}$-等变截面
\[
  s : \operatorname{Op}_{\mathfrak{g}}(\mathcal{O})\times\omega_{\mathcal{O}}^{+}\longrightarrow\mathfrak{g}^{+}\times\omega_{\mathcal{O}}^{+}.
\]

\SourcePage{73}
现在考虑 $\mathcal{D}_{X}$-概形 $(\mathcal{J}\mathfrak{g}^{+})_{X}$、$(\mathcal{J}B)_{X}$ 与 $\operatorname{Op}_{\mathfrak{g}}(\mathcal{O})_{X}=\operatorname{Spec}\mathcal{A}_{\mathfrak{g}}$。显然 $(\mathcal{J}B)_{X}$ 是 $X$ 上的群 $\mathcal{D}_{X}$-概形，而概形 $(\mathcal{J}\mathfrak{g}^{+})_{X}$ 是 $\operatorname{Op}_{\mathfrak{g}}(\mathcal{O})_{X}$ 上的 $(\mathcal{J}B)_{X}$-torsor。事实上 $(\mathcal{J}B)_{X}=\mathcal{J}(B_{X})$，而 $(\mathcal{J}\mathfrak{g}^{+})_{X}$ 是 $X$ 上 $\mathfrak{g}^{+}$-值微分形式的 jets 概形。显然
\[
  \operatorname{Op}_{\mathfrak{g}}=\operatorname{Sect}(\mathfrak{g}_{X}^{+})/\operatorname{Sect}(B_{X})
  =\operatorname{Sect}^{\nabla}((\mathcal{J}\mathfrak{g}^{+})_{X})/\operatorname{Sect}^{\nabla}((\mathcal{J}B)_{X})
  \subset\operatorname{Sect}^{\nabla}(\operatorname{Op}_{\mathfrak{g}}(\mathcal{O})_{X}).
\]
这里 Sect 表示 $X$-概形的截面层，$\operatorname{Sect}^{\nabla}$ 表示 $\mathcal{D}_{X}$-概形的水平截面层。为证明 $\operatorname{Op}_{\mathfrak{g}}=\operatorname{Sect}^{\nabla}(\operatorname{Op}_{\mathfrak{g}}(\mathcal{O})_{X})$，只需证明 $\operatorname{Sect}^{\nabla}((\mathcal{J}\mathfrak{g}^{+})_{X})\to\operatorname{Sect}^{\nabla}(\operatorname{Op}_{\mathfrak{g}}(\mathcal{O})_{X})$ 是满射。为此，使用由 $s$ 诱导的 $\mathcal{D}_{X}$-概形态射 $\operatorname{Op}_{\mathfrak{g}}(\mathcal{O})_{X}\times(\omega_{\mathcal{O}}^{+})_{X}\to(\mathfrak{g}^{+})_{X}$，以及如下事实：$(\omega_{\mathcal{O}}^{+})_{X}$（即处处非零的微分形式的 jets 概形）在 $X$ 的每一点的邻域内都有水平截面。

于是我们把 $\operatorname{Op}_{\mathfrak{g}}(X)$ 等同于 $\operatorname{Op}_{\mathfrak{g}}(\mathcal{O})_{X}=\operatorname{Spec}\mathcal{A}_{\mathfrak{g}}$ 的水平截面的集合。同样地，可把概形 $\operatorname{Op}_{\mathfrak{g}}(X)$ 等同于 $\operatorname{Spec}\mathcal{A}_{\mathfrak{g}}$ 的水平截面概形。

\emph{注.} 我们使用 $s$ 只是为了简化证明 ???。
% [Tutor Note: 原书 p.73 此处引用印作「???」，底本如此，保留原样。]

\bdsec{3.4. $G$-opers 与 $\mathfrak{g}$-opers}
本小节假设 $G$ 半单（其实一般情形也可类似处理，见 3.4.2 末的注 iii）。对每个单负根 $\alpha$，固定非零元 $y_{\alpha}\in\mathfrak{g}^{\alpha}$。令 $G_{ad}:=G/Z$，$B_{ad}:=B/Z$，$H_{ad}:=H/Z$，其中 $Z$ 是 $G$ 的中心。

\bdp{3.4.1} 存在显然的投影 $\operatorname{Op}_{G}(X)\longrightarrow\operatorname{Op}_{\mathfrak{g}}(X):=\operatorname{Op}_{G_{ad}}(X)$。我们将构造一个截面 $\operatorname{Op}_{\mathfrak{g}}(X)\to\operatorname{Op}_{G}(X)$，它依赖于 $\omega_{X}$ 的平方根的选取，即依赖于带同构 $\mathcal{L}^{\otimes 2}\xrightarrow{\sim}\omega_{X}$ 的线丛 $\mathcal{L}$。设 $(\mathfrak{F}_{B_{ad}},\nabla)$ 是 $\mathfrak{g}$-oper。把它提升为 $G$-oper，等价于把 $\mathfrak{F}_{B_{ad}}$ 提升为 $B$-丛，这又等价于把 $\mathfrak{F}_{H_{ad}}$ 提升为 $H$-丛（这里 $\mathfrak{F}_{H_{ad}}$ 是 $\mathfrak{F}_{B_{ad}}$ 沿 $B_{ad}\to H_{ad}$ 的前推）。在 $\mathfrak{g}=\mathfrak{sl}_{2}$ 的特殊情形，我们在 3.1.8 中构造了典范同构 $\mathfrak{F}_{H_{ad}}\xrightarrow{\sim}\omega_{X}$；3.1.8 的构造用到了固定元 $f\in\mathfrak{sl}_{2}/\mathfrak{b}_{0}$。完全类似地，

\SourcePage{74}
一般情形下可构造典范同构 $\mathfrak{F}_{H_{ad}}\xrightarrow{\sim}\lambda\omega_{X}$，其中 $\lambda\omega_{X}$ 是 $\omega_{X}$ 沿同态 $\lambda:\mathbb{G}_{m}\to H_{ad}$ 的前推，而 $\lambda$ 使得对任意单正根 $\alpha$，$\lambda(t)$ 在 $\mathfrak{g}^{\alpha}$ 上乘以 $t$（构造用到 3.4 开头固定的那些元素 $y_{\alpha}$）。存在唯一的同态 $\lambda^{\#}:\mathbb{G}_{m}\to H$，使得
\begin{equation}\tag{54}
  \lambda^{\#}(t)\ \operatorname{mod}Z=\lambda(t)^{2}
\end{equation}
（事实上，$\lambda$ 对应余权 $\check{\rho}:=$ 基本余权之和，而 $2\check{\rho}$ 属于余根格）。我们把 $\mathfrak{F}_{H_{ad}}=\lambda\omega_{X}$ 提升为 $H$-丛 $\lambda^{\#}\mathcal{L}$，其中 $\mathcal{L}$ 是 $\omega_{X}$ 的平方根。

\bdp{3.4.2} 记 $\omega^{1/2}(X)$ 为 $\omega_{X}$ 的平方根的群胚。对固定的 $\mathcal{L}\in\omega^{1/2}(X)$，有等价
\begin{equation}\tag{55}
  \Phi_{\mathcal{L}} : \operatorname{Op}_{\mathfrak{g}}(X)\times Z\operatorname{tors}(X)\xrightarrow{\ \sim\ }\operatorname{Op}_{G}(X),
\end{equation}
其中 $Z\operatorname{tors}(X)$ 是 $X$ 上 $Z$-torsor 的群胚。$\Phi_{\mathcal{L}}(\mathfrak{F},\mathcal{E})$ 的定义如下：用 $\mathcal{L}$ 把 $\mathfrak{F}\in\operatorname{Op}_{\mathfrak{g}}(X)$ 提升为 $G$-oper（见 3.4.1），再用 $\mathcal{E}$ 扭这个 $G$-oper。$\Phi_{\mathcal{L}}$ 对 $\mathcal{L}$ 的依赖方式为
\[
  \Phi_{\mathcal{L}\otimes\mathcal{A}}(\mathfrak{F},\mathcal{E})=\Phi_{\mathcal{L}}(\mathfrak{F},\mathcal{E}\cdot\alpha\mathcal{A}).
\]
这里 $\mathcal{A}$ 是 $\mathcal{O}_{X}$ 的平方根，或者说，是 $X$ 上的 $\mu_{2}$-torsor；而 $\alpha\mathcal{A}$ 是 $\mu_{2}$-torsor $\mathcal{A}$ 沿态射
\begin{equation}\tag{56}
  \alpha : \mu_{2}\longrightarrow Z,\qquad \alpha:=\lambda^{\#}|_{\mu_{2}}
\end{equation}
的前推。回忆 $\lambda^{\#}$ 由 (54) 定义。

\emph{注.}
\begin{enumerate}[label=(\roman*),leftmargin=2.2em]
  \item 若把 $\operatorname{Op}_{\mathfrak{g}}(X)$ 看作概形，把 $\operatorname{Op}_{G}(X)$ 与 $Z\operatorname{tors}(X)$ 看作代数栈，则 (55) 成为代数栈的同构。
  \item $\alpha$ 是「那个」主同态 $\mathrm{SL}_{2}\to G$ 到中心 $\mu_{2}\subset\mathrm{SL}_{2}$ 的限制。

\SourcePage{75}
  \item 若 $G$ 约化但非半单，而 $\mathfrak{g}:=\operatorname{Lie}(G/Z)$，则定义依赖于 $\mathcal{L}\in\omega^{1/2}(X)$ 的截面 $\operatorname{Op}_{\mathfrak{g}}(X)\to\operatorname{Op}_{G}(X)$ 为复合 $\operatorname{Op}_{\mathfrak{g}}(X)\to\operatorname{Op}_{[G,G]}(X)\to\operatorname{Op}_{G}(X)$。如果把 $Z\operatorname{tors}(X)$ 换成 $Z^{\nabla}\operatorname{tors}(X)$，即 $X$ 上带联络的 $Z$-torsor 的群胚，则 3.4.2 的结果仍然成立。
\end{enumerate}

\bdp{3.4.3} 下面是 3.4.2 的一个更自然的表述。先引入群胚 $Z\operatorname{tors}_{\theta}(X)$（$\theta$ 应使读者想到 $\theta$-特征，即 $\omega_{X}$ 的平方根）。$Z\operatorname{tors}_{\theta}(\mathcal{L})$ 的对象是偶对 $(\mathcal{E},\mathcal{L})$，其中 $\mathcal{E}\in Z\operatorname{tors}(X)$，$\mathcal{L}\in\omega^{1/2}(X)$，但我们宁愿写 $\mathcal{E}\cdot\mathcal{L}$ 而不写 $(\mathcal{E},\mathcal{L})$。
% [Tutor Note: 原书 p.75 此处印作 Ztors_{\theta}(\mathcal{L})，按上下文应为 Ztors_{\theta}(X)，底本如此，保留原样。]
令
\[
  \operatorname{Mor}(\mathcal{E}_{1}\cdot\mathcal{L}_{1},\mathcal{E}_{2}\cdot\mathcal{L}_{2})
  :=\operatorname{Mor}(\mathcal{E}_{1},\mathcal{E}_{2}\cdot\alpha(\mathcal{L}_{2}/\mathcal{L}_{1})),
\]
其中 $\alpha(\mathcal{L}_{2}/\mathcal{L}_{1})$ 是 $\mu_{2}$-torsor $\mathcal{L}_{2}/\mathcal{L}_{1}:=\mathcal{L}_{2}\otimes\mathcal{L}_{1}^{\otimes(-1)}$ 沿同态 (56) 的前推。态射的复合按显然方式定义。于是可把 3.4.2 重述为典范等价
\begin{equation}\tag{57}
  \Phi : \operatorname{Op}_{\mathfrak{g}}(X)\times Z\operatorname{tors}_{\theta}(X)\xrightarrow{\ \sim\ }\operatorname{Op}_{G}(X),
\end{equation}
其中 $\Phi(\mathfrak{F},\mathcal{L}\cdot\mathcal{E}):=\Phi_{\mathcal{L}}(\mathfrak{F},\mathcal{E})$，而 $\Phi_{\mathcal{L}}$ 由 (55) 定义。

在 3.2.1 的局部情形有类似的典范等价
\begin{equation}\tag{58}
  \operatorname{Op}_{\mathfrak{g}}(\mathcal{O})\times Z\operatorname{tors}_{\theta}(\mathcal{O})\xrightarrow{\ \sim\ }\operatorname{Op}_{G}(\mathcal{O}),
\end{equation}
其中 $Z\operatorname{tors}_{\theta}(\mathcal{O})$ 按整体情形定义。当然，$Z\operatorname{tors}_{\theta}(\mathcal{O})$ 的所有对象彼此同构，且一个对象的自同构群是 $Z$。对 $Z\operatorname{tors}(\mathcal{O})$ 同样如此。若把 $\mathcal{O}$ 的自同构也考虑进来，$Z\operatorname{tors}_{\theta}(\mathcal{O})$ 与 $Z\operatorname{tors}(\mathcal{O})$ 的差别就清楚了（见 3.5.2）。

\bdp{3.4.4} 要描述 $Z\operatorname{tors}_{\theta}(\mathcal{O})$ 的「经济」版本，我们需要一些抽象废话。

设 $Z$ 是交换群。范畴 $C$ 上的 $Z$-结构是态射 $Z\to\operatorname{Aut}\mathrm{id}_{C}$。等价地，$C$ 上的 $Z$-结构是 $Z$ 在 $\operatorname{Mor}(c_{1},c_{2})$（$c_{1},c_{2}\in\operatorname{Ob}C$）上的作用，使得对任意态射 $c_{1}\xrightarrow{f}c_{2}\xrightarrow{g}c_{3}$ 与任意 $z\in Z$，有 $z(gf)=(zg)f=g(zf)$。$Z$-范畴是带

\SourcePage{76}
$Z$-结构的范畴。设 $C$ 与 $C'$ 是 $Z$-范畴。若对任意 $c_{1},c_{2}\in C$，映射 $\operatorname{Mor}(c_{1},c_{2})\to\operatorname{Mor}(F(c_{1}),F(c_{2}))$ 是 $Z$-等变的，则称函子 $F:C\to C'$ 为 $Z$-函子。若 $A\to Z$ 是交换群的态射，且 $C$ 是 $A$-范畴，则如下定义诱导的 $Z$-范畴 $C\otimes_{A}Z$：$\operatorname{Ob}(C\otimes_{A}Z)=\operatorname{Ob}C$，$(C\otimes_{A}Z)$-态射 $c_{1}\to c_{2}$ 的集合为
\[
  (\operatorname{Mor}_{C}(c_{1},c_{2})\times Z)/A
  =\{\text{由 }A\text{-集 }\operatorname{Mor}_{C}(c_{1},c_{2})\text{ 诱导的 }Z\text{-集}\},
\]
而 $C\otimes_{A}Z$ 中态射的复合按显然方式定义。有自然的 $A$-函子 $C\to C\otimes_{A}Z$；并且对任意 $Z$-范畴 $C'$，任一 $A$-函子 $C\to C'$ 有唯一分解 $C\to C\otimes_{A}Z\xrightarrow{F}C'$，其中 $F$ 是 $Z$-函子。

记 $\omega^{1/2}(\mathcal{O})$ 为 $\omega_{\mathcal{O}}$ 的平方根的群胚。这是 $\mu_{2}$-范畴。$Z\operatorname{tors}(\mathcal{O})$ 与 $Z\operatorname{tors}_{\theta}(\mathcal{O})$ 是 $Z$-范畴。典范 $\mu_{2}$-函子 $\omega^{1/2}(\mathcal{O})\to Z\operatorname{tors}_{\theta}(\mathcal{O})$ 诱导等价
\[
  \omega^{1/2}(\mathcal{O})\otimes_{\mu_{2}}Z\to Z\operatorname{tors}_{\theta}(\mathcal{O}).
\]

\bdp{3.4.5} 读者也许更喜欢 $Z\operatorname{tors}_{\theta}(X)$ 与 $Z\operatorname{tors}_{\theta}(\mathcal{O})$ 的如下「具体」版本。定义正合列
\begin{equation}\tag{59}
  0\to Z\to\widetilde{Z}\to\mathbb{G}_{m}\to 0
\end{equation}
为
\begin{equation}\tag{60}
  0\to\mu_{2}\to\mathbb{G}_{m}\xrightarrow{f}\mathbb{G}_{m}\to 0,\qquad f(x):=x^{2}
\end{equation}
沿态射 (56) 的前推。记 $\widetilde{Z}\operatorname{tors}_{\omega}(X)$ 为把 $\mathbb{G}_{m}$-torsor $\omega_{X}$ 提升为 $\widetilde{Z}$-torsor 所成的群胚（即 $\widetilde{Z}\operatorname{tors}_{\omega}(X)$ 的对象，是 $X$ 上的 $\widetilde{Z}$-torsor，连同相应的 $\mathbb{G}_{m}$-torsor 与 $\omega_{X}$ 之间的一个同构）。由 (60) 到 (59) 的态射诱导函子 $F:\omega^{1/2}(X)\to\widetilde{Z}\operatorname{tors}_{\omega}(X)$。由
\[
  \mathcal{E}\cdot\mathcal{L}\mapsto\mathcal{E}\cdot F(\mathcal{L}),\qquad
  \mathcal{E}\in Z\operatorname{tors}(X),\ \mathcal{L}\in\omega^{1/2}(X)
\]
定义的函子 $Z\operatorname{tors}_{\theta}(X)\to\widetilde{Z}\operatorname{tors}_{\omega}(X)$ 是等价。

完全类似地可定义 $\widetilde{Z}\operatorname{tors}_{\omega}(\mathcal{O})$ 与典范等价 $Z\operatorname{tors}_{\theta}(\mathcal{O})\xrightarrow{\sim}\widetilde{Z}\operatorname{tors}_{\omega}(\mathcal{O})$。

\SourcePage{77}
借助 $\widetilde{Z}\operatorname{tors}_{\omega}(X)$ 与 $\widetilde{Z}\operatorname{tors}_{\omega}(\mathcal{O})$，可以容易地理解等价 (57) 与 (58)。例如，构造等价
\[
  \operatorname{Op}_{\mathfrak{g}}(X)\times\widetilde{Z}\operatorname{tors}_{\omega}(X)\xrightarrow{\ \sim\ }\operatorname{Op}_{G}(X).
\]
考虑如下行正合的交换图表：
\begin{equation}\tag{61}
  \begin{tikzcd}
    0\arrow[r] & \mu_{2}\arrow[r]\arrow[d, "\alpha"'] & \mathbb{G}_{m}\arrow[r]\arrow[d, "\lambda^{\#}"'] & \mathbb{G}_{m}\arrow[r]\arrow[d, "\lambda"] & 0\\
    0\arrow[r] & Z\arrow[r] & H\arrow[r] & H_{ad}\arrow[r] & 0
  \end{tikzcd}
\end{equation}
这里上行是 (60)；下行与态射 $\lambda,\lambda^{\#}$ 在 3.4.1 中定义。按 3.4.1，$X$ 上的 $G$-oper 等同于 $X$ 上的 $\mathfrak{g}$-oper 加上把 $H_{ad}$-torsor $\lambda_{*}(\omega_{X})$ 提升为 $H$-torsor。这样的提升等同于 $\widetilde{Z}\operatorname{tors}_{\omega}(X)$ 的一个对象：考察交换图表
\begin{equation}\tag{62}
  \begin{tikzcd}
    0\arrow[r] & Z\arrow[r]\arrow[d, equal] & \widetilde{Z}\arrow[r]\arrow[d] & \mathbb{G}_{m}\arrow[r]\arrow[d, "\lambda"] & 0\\
    0\arrow[r] & Z\arrow[r] & H\arrow[r] & H_{ad}\arrow[r] & 0
  \end{tikzcd}
\end{equation}
的右端（笛卡尔）方块（(62) 的上行是 (59)，而 (62) 与 (61) 的下行相同）。

\bdp{3.4.6} $Z\operatorname{tors}(X)$ 是（严格交换的）Picard 范畴（见 [Del73] 中的定义 1.4.2），而 $Z\operatorname{tors}_{\theta}(X)$ 是 $Z\operatorname{tors}(X)$ 上的一个 torsor；确切地说，$Z\operatorname{tors}_{\theta}(X)$ 是由 $\mu_{2}\operatorname{tors}(X)$ 上的 torsor $\omega^{1/2}(X)$ 沿对应于 (56) 的 Picard 函子 $\mu_{2}\operatorname{tors}(X)\to Z\operatorname{tors}(X)$ 诱导而来。我们在 §4 中将使用这套语言，所以这里回忆一下定义。

\emph{Picard 范畴}是指 [De-Mi] 意义下的张量范畴 $\mathcal{A}$（即对称交换幺半范畴），使得 $\mathcal{A}$ 的所有态射可逆（即 $\mathcal{A}$ 是群胚），且 $\mathcal{A}$ 的所有对象可逆；也就是说，对每个 $a\in\operatorname{Ob}\mathcal{A}$，存在 $a'\in\operatorname{Ob}\mathcal{A}$，使 $a\cdot a'$ 是单位对象（我们用 $\cdot$ 记「张量积」函子 $\mathcal{A}\times\mathcal{A}\to\mathcal{A}$；在 [De-Mi] 与 [Del73] 中它分别记作 $\otimes$ 与 $+$）。严格交换是指

\SourcePage{78}
对每个 $a\in\operatorname{Ob}\mathcal{A}$，交换同构 $a\otimes a\xrightarrow{\sim}a\otimes a$ 是恒等。

幺半范畴 $\mathcal{A}$ 在范畴 $C$ 上的一个\emph{作用}是函子 $\mathcal{A}\times C\to C$（记作 $\cdot$），连同结合约束，即函子同构 $(a_{1}\cdot a_{2})\cdot c\xrightarrow{\sim}a_{1}\cdot(a_{2}\cdot c)$（$a_{i}\in\mathcal{A}$，$c\in C$），它满足一个五边形公理，该公理与 $\mathcal{A}$ 中结合约束的五边形公理类似（见 [Del73] 与 [De-Mi]）；我们还要求对应于 $\mathcal{A}$ 的单位对象的函子 $F:C\to C$ 全忠实（此时同构 $F^{2}\xrightarrow{\sim}F$ 给出典范同构 $F\xrightarrow{\sim}\mathrm{id}$）。这一定义可在 [Pa] 与 [Yet] 的 §3 中找到。\emph{$\mathcal{A}$-模}是带 $\mathcal{A}$ 的作用的范畴。若 $C$ 与 $\widetilde{C}$ 是 $\mathcal{A}$-模，则\emph{$\mathcal{A}$-模函子} $C\to\widetilde{C}$ 是带函子同构 $\Phi(a\cdot c)\xrightarrow{\sim}a\cdot\Phi(c)$ 的函子 $\Phi:C\to\widetilde{C}$，且满足自然的相容条件（构造同构 $\Phi((a_{1}\cdot a_{2})\cdot c)\xrightarrow{\sim}a_{1}\cdot(a_{2}\cdot\Phi(c))$ 的两种方式必须给出同一结果；见 [Pa]、[Yet]）。$\mathcal{A}$-模函子也称为 $\mathcal{A}$-模的\emph{态射}。

Picard 范畴 $\mathcal{A}$ 上的一个\emph{torsor} 是 $\mathcal{A}$-模，使得对某个 $c\in\operatorname{Ob}C$，函子 $a\mapsto a\cdot c$ 是 $\mathcal{A}$ 与 $C$ 之间的等价（此时对一切 $c\in\operatorname{Ob}C$ 都成立）。

设 $\mathcal{A}$ 与 $\mathcal{B}$ 是 Picard 范畴。\emph{Picard 函子} $\mathcal{A}\to\mathcal{B}$ 是带函子同构 $F(a_{1}\cdot a_{2})\xrightarrow{\sim}F(a_{1})\cdot F(a_{2})$ 的函子 $F:\mathcal{A}\to\mathcal{B}$，且该同构与交换约束、结合约束相容。于是 $F$ 把 $\mathcal{A}$ 的单位对象送到 $\mathcal{B}$ 的单位对象，即 $F$ 是 [De-Mi] 意义下的张量函子。在 [Del73] 中，Picard 函子称为加法函子。

设 $F:\mathcal{A}_{1}\to\mathcal{A}_{2}$ 是 Picard 函子，$C_{i}$ 是 $\mathcal{A}_{i}$ 上的 torsor（$i=1,2$）。则 $C_{2}$ 带有 $\mathcal{A}_{1}$ 的作用。在这种情形下，$\mathcal{A}_{1}$-模函子 $C_{1}\to C_{2}$ 称为 $F$-\emph{仿射}函子。

\emph{例.}
\begin{enumerate}[label=\arabic*),leftmargin=2.2em]
  \item 设 $A$ 是交换代数群。则 $A\operatorname{tors}(X)$ 有典范的 Picard 范畴结构。

\SourcePage{79}
  \item 交换代数群的态射 $A\to B$ 诱导 Picard 函子 $A\operatorname{tors}(X)\to B\operatorname{tors}(X)$。
  \item 3.4.2 中的群胚 $\omega^{1/2}(X)$ 是 Picard 范畴 $\mu_{2}\operatorname{tors}(X)$ 上的 torsor。3.4.3 中的群胚 $Z\operatorname{tors}_{\theta}(X)$ 是 $Z\operatorname{tors}(X)$ 上的 torsor。
\end{enumerate}

若 $F:\mathcal{A}\to\mathcal{B}$ 是 Picard 范畴之间的 Picard 函子，$C$ 是 $\mathcal{A}$ 上的 torsor，则可得 $\mathcal{B}$ 上的\emph{诱导} torsor $\mathcal{B}\cdot_{\mathcal{A}}C$。读者可以从下例复原 $\mathcal{B}\cdot_{\mathcal{A}}C$ 的定义（见 3.4.3）：若 $\mathcal{A}=\mu_{2}\operatorname{tors}(X)$，$\mathcal{B}=Z\operatorname{tors}(X)$，$F$ 来自 (56)，$C=\omega^{1/2}(X)$，则 $\mathcal{B}\cdot_{\mathcal{A}}C=Z\operatorname{tors}_{\theta}(X)$。$\mathcal{B}\cdot_{\mathcal{A}}C$ 的对象记作 $b\cdot c$（$b\in\operatorname{Ob}\mathcal{B}$，$c\in\operatorname{Ob}C$，见 3.4.3）。

感兴趣的读者可以表述 $\mathcal{B}\cdot_{\mathcal{A}}C$ 的泛性质。我们需要如下较弱的性质。给定带 $\mathcal{B}$ 作用的范畴 $\widetilde{C}$ 与 $\mathcal{A}$-模函子 $\Phi:C\to\widetilde{C}$，有自然的方式构造 $\mathcal{B}$-模函子 $\Psi:\mathcal{B}\cdot_{\mathcal{A}}C\to\widetilde{C}$：令 $\Psi(b\cdot c):=b\cdot\Phi(c)$，并在态射上按显然方式定义 $\Psi$（即映射 $\operatorname{Mor}(b_{1}\cdot c_{1},b_{2}\cdot c_{2})\to\operatorname{Mor}(b_{1}\cdot\Phi(c_{1}),b_{2}\cdot\Phi(c_{2}))$ 是复合
\[
  \operatorname{Mor}(b_{1}\cdot c_{1},b_{2}\cdot c_{2})
  =\operatorname{Mor}(b_{1},b_{2}\cdot c_{2}/c_{1})
  \to\operatorname{Mor}(b_{1}\cdot\Phi(c_{1}),b_{2}\cdot c_{2}/c_{1}\cdot\Phi(c_{1}))
  \xrightarrow{\sim}\operatorname{Mor}(b_{1}\cdot\Phi(c_{1}),b_{2}\cdot\Phi(c_{2}))).
\]
同构 $\Psi(b_{1}\cdot(b_{2}\cdot c))\xrightarrow{\sim}b_{1}\cdot\Psi(b_{2}\cdot c)$ 是显然的那一个。

我们在下述情形使用这一构造。设有 Picard 函子 $\ell:\mathcal{B}\to\widetilde{\mathcal{B}}$、$\widetilde{\mathcal{B}}$ 上的 torsor $\widetilde{C}$、以及 $\ell'$-仿射函子 $\Phi:C\to\widetilde{C}$，其中 $\ell'$ 是复合 $\mathcal{A}\xrightarrow{F}\mathcal{B}\xrightarrow{\ell}\widetilde{\mathcal{B}}$。则上述构造给出 $\ell$-仿射函子 $\mathcal{B}\cdot_{\mathcal{A}}C\to\widetilde{C}$。

\bdp{3.4.7} 设 $Z$ 是交换群，$Z\operatorname{tors}$ 是 $Z$-torsor 的 Picard 范畴（一点上的）。下述注记将在 4.4.9 中使用。

\emph{注.}
\begin{enumerate}[label=(\roman*),leftmargin=2.2em]
  \item 从 $Z\operatorname{tors}$ 到 Picard 范畴 $\mathcal{A}$ 的 Picard 函子，「等同于」一个态射 $Z\to\operatorname{Aut}\mathbf{1}_{\mathcal{A}}$，其中 $\mathbf{1}_{\mathcal{A}}$ 是 $\mathcal{A}$ 的单位对象。更确切地说，从 Picard 函子范畴

\SourcePage{80}
$Z\operatorname{tors}\to\mathcal{A}$ 到 $\operatorname{Hom}(Z,\operatorname{Aut}\mathbf{1}_{\mathcal{A}})$ 的自然函子是等价。这里集合 $\operatorname{Hom}(Z,\operatorname{Aut}\mathbf{1}_{\mathcal{A}})$ 视为离散范畴。
  \item 把「Picard」换成「幺半」后，前述注记仍然成立。
  \item $Z\operatorname{tors}$ 在范畴 $C$ 上的一个作用，「等同于」 $C$ 上的 $Z$-结构，即态射 $Z\to\operatorname{Aut}\mathrm{id}_{C}$（注意幺半范畴 $\mathcal{A}$ 在 $C$ 上的作用等同于幺半函子 $\mathcal{A}\to\operatorname{Funct}(C,C)$，再应用前一注记）。
  \item 设 $C_{1}$ 与 $C_{2}$ 是 $Z\operatorname{tors}$ 上的模。按前一注记，$C_{1}$ 与 $C_{2}$ 是 3.4.4 意义下的 $Z$-范畴。容易看出，$(Z\operatorname{tors})$-模函子 $C_{1}\to C_{2}$ 与 3.4.4 意义下的 $Z$-函子相同（即函子 $F:C_{1}\to C_{2}$ 至多有一个 $(Z\operatorname{tors})$-模函子结构，且这样的结构存在当且仅当 $F$ 是 $Z$-函子）。
  \item $Z\operatorname{tors}$ 上的 torsor 「等同于」一个与 $Z\operatorname{tors}$ $Z$-等价的 $Z$-范畴。（我们需要这一条吗？？？）
% [Tutor Note: 原书 p.80 括注「(do we need this???)」为作者存疑之语，底本如此，保留原样。]
\end{enumerate}

\bdsec{3.5. 局部 opers II}
对大多数李代数 $\mathfrak{g}$（例如 $\mathfrak{g}=\mathfrak{sl}_{n}$，$n>5$），定理 3.2.2 中提到的性质并不能\emph{唯一地}确定 Feigin–Frenkel 同构 (49)。原因是 $A_{\mathfrak{g}}(\mathcal{O})$ 有许多在 $\operatorname{gr}A_{\mathfrak{g}}(\mathcal{O})$ 上诱导恒等的 $\operatorname{Aut}\mathcal{O}$-等变自同构。这一点从 3.2.1 中 $\operatorname{Op}_{\mathfrak{g}}(\mathcal{O})=\operatorname{Spec}A_{\mathfrak{g}}(\mathcal{O})$ 的几何描述即可看出，也可从将在 3.5.6 中给出的 $A_{\mathfrak{g}}(\mathcal{O})$ 的描述看出（见 (65)–(68)）。3.5–3.6 的目标是表述唯一确定 Feigin–Frenkel 同构的性质 3.6.7。这一性质以及 3.6.11 将用于我们主定理 5.2.6 的证明。在 3.7 与 3.8 中我们说明如何从 [FF92] 提取性质 3.6.7 与 3.6.11。人们可以（或许应该）在 3.5–3.8 之前先读 §4 与（§5 的一大部分？）。我们当然建议读者，在 ?? 中用到 3.6.11 之前，先跳过 3.5.16–3.5.18 与 3.6.8–3.6.11。
% [Tutor Note: 原书 p.80 有作者留待补全的「?」与残缺引用「??」，底本如此，保留原样。]

% ---- ch03_u03 ----
% ch03_u03 —— 原书 p.81–p.90（3.5 中段至 3.6 开头）
\SourcePage{81}
3.5 与 3.6 两节的思路\footnote{受 [Phys] 启发。}是：给 $A_{\mathfrak{g}}(O)$ 及其对应物 $\mathfrak{z}_{{}^{L}\mathfrak{g}}(O)$ 附加某些结构，以「消去」它们的自同构。对 $A_{\mathfrak{g}}(O)$ 而言，所附加的是 3.5.11 中的李代数胚 $\mathfrak{a}_{\mathfrak{g}}$。它在 $\mathfrak{z}_{{}^{L}\mathfrak{g}}(O)$ 中的对应物将在 3.6.5 中引入。$\mathfrak{a}_{\mathfrak{g}}$ 的定义很简单：它是 $\operatorname{Op}_{\mathfrak{g}}(O)$ 上万有 $G$-丛 $\mathfrak{F}^{0}_{G}$ 的无穷小对称代数胚。$\mathfrak{F}^{0}_{G}$ 从而 $\mathfrak{a}_{\mathfrak{g}}$ 都带有 $\operatorname{Der}O$ 的作用。结果是，偶 $(A_{\mathfrak{g}}(O),\mathfrak{F}^{0}_{G})$ 没有非平凡的 $\operatorname{Der}O$-等变自同构（见 3.5.9）；对 $(A_{\mathfrak{g}}(O),\mathfrak{a}_{\mathfrak{g}})$ 而言，这一点「几乎」成立（见 3.5.13）。

\bdp{3.5.1} 我们有 $\operatorname{Spec}O$ 上由 3.2.1 中的概形 $\operatorname{Op}_{\mathfrak{g}}(O)=\operatorname{Spec}A_{\mathfrak{g}}(O)$ 参数化的通用 $\mathfrak{g}$-oper 族。固定一个一维自由 $O$-模 $\omega_{O}^{1/2}$，并配上同构 $\omega_{O}^{1/2}\otimes\omega_{O}^{1/2}\xrightarrow{\sim}\omega_{O}$（当然，$\omega_{O}^{1/2}$ 在同构意义下唯一）。于是上述通用族提升为 $G$-oper 族；见 3.4.1\footnote{说实话，我们还须对每个负单根 $\alpha$ 选取非零的 $y_{\alpha}\in\mathfrak{g}^{\alpha}$（见 3.4.1）。}。这样，我们得到 $\operatorname{Spec}(A_{\mathfrak{g}}(O)\widehat\otimes O)=\operatorname{Spec}A_{\mathfrak{g}}(O)[[t]]$ 上的 $B$-丛 $\mathfrak{F}_{B}$，以及相伴 $G$-丛 $\mathfrak{F}_{G}$ 上沿 $\operatorname{Spec}O$ 的联络 $\nabla$。

\bdp{3.5.2} 考虑群归纳概形 $\operatorname{Aut}_{2}O:=\operatorname{Aut}(O,\omega_{O}^{1/2})$。我们有典范正合列
\begin{equation}\tag{63}
0\to\mu_{2}\to\operatorname{Aut}_{2}O\to\operatorname{Aut}O\to0,
\end{equation}
且 $\operatorname{Aut}_{2}O$ 连通。正合列 (63) 与连通性可以作为 $\operatorname{Aut}_{2}O$ 的另一定义。以 $\operatorname{Aut}_{2}^{0}O$ 记 $\operatorname{Aut}^{0}O$ 在 $\operatorname{Aut}_{2}O$ 中的原像。

$\operatorname{Aut}O$ 作用于 $A_{\mathfrak{g}}(O)$ 与 $O$，从而作用于 $\operatorname{Spec}(A_{\mathfrak{g}}(O)\widehat\otimes O)$。该作用典范地提升为 $\operatorname{Aut}_{2}O$ 在 $(\mathfrak{F}_{B},\nabla)$ 上的作用。$\mu_{2}\subset\operatorname{Aut}_{2}O$ 经态射 (56) 作用于 $\mathfrak{F}_{B}$。

\SourcePage{82}
\bdp{3.5.3} \emph{引理.} 设 $L$ 是代数群，$A$ 是带 $\operatorname{Aut}O$ 作用的代数。考虑 $\operatorname{Aut}O$ 经其在 $A$ 与 $O$ 上的作用诱导的、在 $A\widehat\otimes O$ 上的作用。令 $i:\operatorname{Spec}A\hookrightarrow\operatorname{Spec}(A\widehat\otimes O)$ 为自然嵌入，$\pi:\operatorname{Spec}(A\widehat\otimes O)\to\operatorname{Spec}A$ 为投影。

\begin{enumerate}[label=\arabic*),leftmargin=2.2em]
\item $i^{*}$ 是 $\operatorname{Spec}(A\widehat\otimes O)$ 上 $\operatorname{Aut}_{2}O$-等变 $L$-丛的范畴与 $\operatorname{Spec}A$ 上 $\operatorname{Aut}_{2}^{0}O$-等变 $L$-丛的范畴之间的等价。
\item $\pi^{*}$ 是 $\operatorname{Spec}A$ 上 $\operatorname{Aut}_{2}O$-等变 $L$-丛的范畴与 $\operatorname{Spec}(A\widehat\otimes O)$ 上带沿 $\operatorname{Spec}O$ 的 $\operatorname{Aut}_{2}O$-不变联络的 $\operatorname{Aut}_{2}O$-等变 $L$-丛的范畴之间的等价。
\item 这些等价与遗忘函子（从 $\operatorname{Spec}A$ 上 $\operatorname{Aut}_{2}O$-等变丛到 $\operatorname{Spec}A$ 上 $\operatorname{Aut}_{2}^{0}O$-等变丛）及遗忘函子（从带联络的丛到丛）相容。\hfill$\Box$
\end{enumerate}

\bdp{3.5.4} 以 $\mathfrak{F}^{0}_{B}$ 与 $\mathfrak{F}^{0}_{G}$ 记 $\mathfrak{F}_{B}$ 与 $\mathfrak{F}_{G}$ 到 $\operatorname{Op}_{\mathfrak{g}}(O)=\operatorname{Spec}A_{\mathfrak{g}}(O)\subset\operatorname{Spec}A_{\mathfrak{g}}(O)\widehat\otimes O$ 的限制。$\mathfrak{F}^{0}_{B}$ 是 $\operatorname{Op}_{\mathfrak{g}}(O)$ 上的 $B$-丛，$\mathfrak{F}^{0}_{G}$ 是相应的 $G$-丛。$\mathfrak{F}^{0}_{B}$ 是 $\operatorname{Aut}_{2}^{0}O$-等变的；由 3.5.3，$\mathfrak{F}^{0}_{G}$ 是 $\operatorname{Aut}_{2}O$-等变的。由于 $\mathfrak{F}_{G}$ 上的联络 $\nabla$ 不保持 $\mathfrak{F}_{B}$，$\operatorname{Aut}_{2}O$ 在 $\mathfrak{F}^{0}_{G}$ 上的作用不保持 $\mathfrak{F}^{0}_{B}$。由 3.5.3，带 $\operatorname{Aut}_{2}O$ 作用的 $\mathfrak{F}^{0}_{G}$ 与 $B$-结构 $\mathfrak{F}^{0}_{B}\subset\mathfrak{F}^{0}_{G}$ 一起「记得」通用 oper $(\mathfrak{F}_{B},\nabla)$。

\bdp{3.5.5} 以 $F^{0}_{H}$ 记 $\operatorname{Op}_{\mathfrak{g}}(O)$ 上对应于 $\mathfrak{F}^{0}_{B}$ 的 $H$-丛。由于 $\operatorname{Op}_{\mathfrak{g}}(O)$ 是（无穷维）仿射空间，$\operatorname{Op}_{\mathfrak{g}}(O)$ 上任何 $H$-丛都平凡，且 $H$ 在其平凡化集合上的作用传递。特别地，这适用于 $\mathfrak{F}^{0}_{H}$；于是 $\mathfrak{F}_{H}$ 是 $\operatorname{Spec}\mathbb{C}$ 上某个 $H$-丛 $F_{H}$ 的拉回。由 3.4.1，$F_{H}$ 是 $\operatorname{Spec}\mathbb{C}$ 上的 $\mathbb{G}_{m}$-丛 $\omega_{O}^{1/2}/t\omega_{O}^{1/2}$ 沿 (54) 所定义态射 $\lambda^{\#}:\mathbb{G}_{m}\to H$ 的前推。特别地，$\operatorname{Aut}_{2}^{0}O$ 在 $F_{H}$ 上的作用来自复合
\[
\operatorname{Aut}_{2}^{0}O\to\operatorname{Aut}(\omega_{O}^{1/2}/t\omega_{O}^{1/2})=\mathbb{G}_{m}\xrightarrow{\lambda^{\#}}H.
\]

\SourcePage{83}
于是 $\operatorname{Der}^{0}O$ 在 $F_{H}$ 上的作用是「显然」作用（即保持 $F_{H}$ 的任一平凡化的作用）与态射 $\operatorname{Der}^{0}O\to\mathfrak{h}$ 之和，后者由 $f(t)\cdot t\frac{d}{dt}\mapsto f(0)\check\rho$ 定义。这里 $\check\rho$ 是基本余权之和。

\bdp{3.5.6} 这里给出 $A_{\mathfrak{g}}(O)$ 与 $\mathfrak{F}^{0}_{G}$ 的一个显式描述，其精神与 3.1.9–3.1.10 相同。设 $e,f\in\mathfrak{sl}_{2}$ 为 3.1.8 中的矩阵。固定主嵌入 $i:\mathfrak{sl}_{2}\hookrightarrow\mathfrak{g}$，使得 $i(e)\in\mathfrak{b}$。若选取嘉当子代数 $\mathfrak{h}\subset\mathfrak{b}$ 使 $i([e,f])\in\mathfrak{h}$，则 $i([e,f])$ 可与 $2\check\rho$ 等同。正如 3.1.9，置 $V:=\operatorname{Ker}\operatorname{ad}i(e)$。选取基 $e_{1},\dots,e_{r}\in V$，使得 $e_{1}=i(e)$，且所有 $e_{j}$ 都是 $\operatorname{ad}\check\rho$ 的特征向量。事实上，$[\check\rho,e_{j}]=(d_{j}-1)e_{j}$，其中 $d_{j}$ 是 $\mathfrak{g}$ 上「基本」不变多项式的次数（特别地 $d_{1}=2$）。平凡 $G$-丛上的联络
\begin{equation}\tag{64}
\nabla_{\frac{d}{dt}}=\frac{d}{dt}+i(f)+u_{1}(t)e_{1}+\cdots+u_{r}(t)e_{r}
\end{equation}
定义一个 $\mathfrak{g}$-oper；由 3.1.10，这在 $\operatorname{Spec}O$ 上的 $\mathfrak{g}$-oper（$O:=\mathbb{C}[[t]]$）与 $r$-元组 $(u_{1}(t),\dots,u_{r}(t))$、$u_{j}(t)\in\mathbb{C}[[t]]$ 之间给出双射。把 $u_{j}(t)$ 写成 $u_{j0}+u_{j1}t+\cdots$。于是 $A_{\mathfrak{g}}(O)$ 是 $u_{jk}$（$1\le j\le r$，$0\le k<\infty$）的多项式环。3.5.1 与 3.5.4 中的丛 $\mathfrak{F}_{B},\mathfrak{F}_{G},\mathfrak{F}^{0}_{B},\mathfrak{F}^{0}_{G}$ 都平凡；我们通过选取 opers 的典范形式 (64) 将它们平凡化。

为描述 $\operatorname{Der}O$ 在 $A_{\mathfrak{g}}(O)$ 与 $\mathfrak{F}^{0}_{G}$ 上的作用，引入标准记号 $L_{n}:=-t^{n+1}\frac{d}{dt}\in\operatorname{Der}\mathbb{C}((t))$（于是 $n\ge-1$ 时 $L_{n}\in\operatorname{Der}O$）。置 $u_{j}:=u_{j0}$。则
\begin{equation}\tag{65}
u_{jk}=(L_{-1})^{k}u_{j}/k!
\end{equation}
\begin{equation}\tag{66}
L_{0}u_{j}=d_{j}u_{j}
\end{equation}
\begin{equation}\tag{67}
L_{n}u_{j}=0\qquad\text{若 }n>0,\ j\ne1
\end{equation}
\begin{equation}\tag{68}
L_{n}u_{1}=0\qquad\text{若 }n>0,\ n\ne2;\qquad L_{2}u_{1}=-3.
\end{equation}
于是 $A_{\mathfrak{g}}(O)$ 是以 $u_{1},\dots,u_{r}$ 为生成元、以关系 (66)–(68) 为定义关系的交换 $(\operatorname{Der}O)$-代数。以 $L_{n}^{hor}$ 记 $\mathfrak{F}^{0}_{G}$ 上的向量场，它\SourcePage{84}%
来自 $\mathfrak{F}^{0}_{G}$ 的平凡化与 $\operatorname{Der}O$ 在 $A_{\mathfrak{g}}(O)$ 上的作用。

$L_{n}$ 在 $\mathfrak{F}^{0}_{G}$ 上的作用为 $L_{n}^{hor}+M_{n}$，其中 $M_{n}\in\mathfrak{g}\otimes A_{\mathfrak{g}}(O)$。可以证明
\begin{equation}\tag{69}
M_{0}=-\check\rho
\end{equation}
\begin{equation}\tag{70}
M_{1}=-i(e),\qquad M_{n}=0\quad(n>1)
\end{equation}
\[
M_{-1}=i(f)+u_{1}e_{1}+\cdots+u_{r}e_{r}.
\]
后文只会用到 (69)（恐怕至少还会用到 (70)）!???!）。
% [Tutor Note: 原书 p.84 此处作「... (I’m afraid we’ll use at least (70)) !???!).」，括号不配对，底本如此，保留原样。]

\emph{注.} 若 $n\ge0$，则 $M_{n}\in i(\mathfrak{b}_{0})\subset\mathfrak{b}\subset\mathfrak{b}\otimes A_{\mathfrak{g}}(O)$，其中 $\mathfrak{b}_{0}:=\operatorname{Lie}B_{0}$，而 $B_{0}\subset SL_{2}$ 是上三角矩阵群。这意味着，我们把 $\operatorname{Aut}_{2}^{0}O$-等变丛 $\mathfrak{F}^{0}_{B}$ 与 $\operatorname{Spec}\mathbb{C}$ 上某个 $\operatorname{Aut}_{2}^{0}O$-等变 $B$-丛的拉回等同起来，而后者来自某个态射 $\operatorname{Aut}_{2}^{0}O\to B_{0}\to B$（参看 3.1.10 的注 (iii)）。

\bdp{3.5.7} 设 $A$ 是带 $\operatorname{Aut}O$ 作用的代数。则 $\operatorname{Der}O$ 作用于 $A$，$L_{0}$ 在 $A$ 上的作用可对角化，且 $L_{0}:A\to A$ 的特征值为整数。假设 $L_{0}:A/\mathbb{C}\to A/\mathbb{C}$ 的特征值为正。则 $A=\mathbb{C}\oplus A_{+}$，其中 $A_{+}$ 是 $A$ 中对应于正特征值的所有特征子空间之和。$A_{+}$ 是 $A$ 的唯一 $L_{0}$-不变极大理想。$\operatorname{Spec}A$ 中对应的点记作 $0$。由于 $[L_{0},L_{n}]=-nL_{n}$，有 $L_{-1}A_{+}\subset A_{+}$。假设
\begin{equation}\tag{71}
L_{1}A_{+}\subset A_{+}.
\end{equation}
特别地，若 $L_{0}$ 在 $A/\mathbb{C}$ 上的特征值大于 $1$，则 (71) 成立；例如 $A=A_{\mathfrak{g}}(O)$ 时如此（见 (66) 与 (65)）。

假设 $G$ 是 $\mathfrak{g}$ 对应的伴随群。设 $\mathcal{E}$ 是 $\operatorname{Spec}A$ 上的 $\operatorname{Aut}O$-等变 $G$-丛。代数 $\mathbb{C}L_{-1}+\mathbb{C}L_{0}+\mathbb{C}L_{1}\simeq\mathfrak{sl}_{2}$ 稳定 $0\in\operatorname{Spec}A$，故它作用在 $\mathcal{E}$ 于 $0$ 处的纤维上。于是得到态射 $\sigma:\mathfrak{sl}_{2}\to\mathfrak{g}$，它在共轭意义下定义。

\SourcePage{85}
\emph{例.} 点 $0\in\operatorname{Spec}A_{\mathfrak{g}}(O))=\operatorname{Op}_{\mathfrak{g}}(O)$ 是 $\mathfrak{sl}_{2}$-oper 沿主嵌入 $\mathfrak{sl}_{2}\to\mathfrak{g}$ 的前推，该 $\mathfrak{sl}_{2}$-oper 对应于 Sturm--Liouville 算子 $(d/dt)^{2}$。若 $A=A_{\mathfrak{g}}(O)$ 且 $\mathcal{E}=\mathfrak{F}^{0}_{G}$，则 $\sigma$ 是主嵌入。
% [Tutor Note: 原书 p.85 此处「$\operatorname{Spec} A_{\mathfrak{g}}(O))$」多一个右括号，底本如此，保留原样。]

\bdp{3.5.8} \emph{命题.}
\begin{enumerate}[label=\arabic*),leftmargin=2.2em]
\item 以下条件等价：
\begin{enumerate}[label=\alph*),leftmargin=2.2em]
\item $\operatorname{Aut}O$-等变 $G$-丛 $\mathcal{E}$ 同构于 $\varphi^{*}\mathfrak{F}^{0}_{G}$，其中 $\varphi:\operatorname{Spec}A\to\operatorname{Op}_{\mathfrak{g}}(O)$ 是某个 $\operatorname{Aut}O$-等变态射；
\item $\mathcal{E}$ 上存在 $\operatorname{Aut}^{0}O$-不变的 $B$-结构，使得相应的 $\operatorname{Aut}^{0}O$-等变 $H$-丛同构于 3.5.5\footnote{由于 $G$ 是伴随群，$\operatorname{Aut}_{2}^{0}O$ 在 $F_{H}$ 上的作用经 $\operatorname{Aut}^{0}O$ 分解。} 所定义的、$\operatorname{Spec}\mathbb{C}$ 上 $\operatorname{Aut}^{0}O$-等变 $H$-丛 $F_{H}$ 的拉回；
\item $\sigma:\mathfrak{sl}_{2}\to\mathfrak{g}$ 是主嵌入。
\end{enumerate}
\item a) 中提到的态射 $\varphi$ 与同构 $\mathcal{E}\xrightarrow{\sim}\varphi^{*}\mathfrak{F}^{0}_{G}$ 都唯一。
\item b) 中提到的 $B$-结构唯一。
\end{enumerate}

\emph{证明.} 由 3.5.5，a) 蕴含 b)。为从 b) 推出 c)，只需看 $0\in\operatorname{Spec}A$ 上的情形。下面从 b) 推出 a)，并说明 2) 由 3) 推出。为此只需证明：若固定具有 b) 所述性质的 $B$-结构 $\mathcal{E}_{B}\subset\mathcal{E}$，则恰有一种方式构造 $\operatorname{Aut}O$-等变态射 $\varphi:\operatorname{Spec}A\to\operatorname{Op}_{\mathfrak{g}}(O)$ 与同构 $f:\mathcal{E}\xrightarrow{\sim}\varphi^{*}\mathfrak{F}^{0}_{G}$，使得 $f(\mathcal{E}_{B})=\varphi^{*}\mathfrak{F}^{0}_{B}$。由 3.5.3，$\mathcal{E}$ 与 $\mathcal{E}_{B}$ 给出 $\operatorname{Spec}(A\widehat\otimes O)$ 上的 $G$-丛 $\widetilde{\mathcal{E}}_{G}$，带 $B$-结构 $\widetilde{\mathcal{E}}_{B}\subset\widetilde{\mathcal{E}}_{G}$、沿 $\operatorname{Spec}O$ 的联络 $\nabla$，以及 $\operatorname{Aut}O$ 在 $(\widetilde{\mathcal{E}}_{G},\widetilde{\mathcal{E}}_{B},\nabla)$ 上的作用。现在 $\varphi$ 与 $f$ 的唯一性清楚；为证其存在，须证 $(\widetilde{\mathcal{E}}_{G},\widetilde{\mathcal{E}}_{B},\nabla)$ 是 oper 族，即须证 3.1.2 定义的 $c(\nabla)$ 满足定义 3.1.3 的条件 1 与 2。在我们情形，$c(\nabla)$ 是 $(\mathfrak{g}/\mathfrak{b})_{\widetilde{\mathcal{E}}}\otimes_{O}\omega_{O}$ 的 $\operatorname{Aut}O$-不变截面；只需对其到 $\operatorname{Spec}A\subset\operatorname{Spec}(A\widehat\otimes O)$ 的限制 $c_{0}(\nabla)$ 验证条件 1 与 2。

\SourcePage{86}
$c_{0}(\nabla)$ 是 $H^{0}(\operatorname{Spec}A,(\mathfrak{g}/\mathfrak{b})_{\mathcal{E}_{B}})\otimes\omega_{O}/t\omega_{O}$ 的 $\operatorname{Aut}^{0}O$-不变元。令 $\operatorname{gr}^{k}\mathfrak{g}$ 与 3.1.1 中含义相同。既然知道 $\mathcal{E}_{B}$ 对应的 $H$-丛，便看出存在 $\operatorname{Aut}^{0}O$-等变同构
\begin{equation}\tag{72}
H^{0}(\operatorname{Spec}A,\operatorname{gr}^{k}\mathfrak{g}_{\mathcal{E}_{B}})\otimes\omega_{O}/t\omega_{O}\xrightarrow{\sim}A\otimes(\omega_{O}/t\omega_{O})^{\otimes(k+1)}\otimes\operatorname{gr}^{k}\mathfrak{g}.
\end{equation}
由于 $L_{0}$ 在 $(\omega_{O}/t\omega_{O})^{\otimes(k+1)}$ 上按乘以 $-k-1$ 的方式作用，$A\otimes(\omega_{O}/t\omega_{O})^{\otimes(k+1)}$ 的 $\operatorname{Aut}^{0}O$-不变部分在 $k<-1$ 时为 $0$，在 $k=-1$ 时为 $\mathbb{C}$。因此
\[
c_{0}(\nabla)\in\operatorname{gr}^{-1}\mathfrak{g}\subset A\otimes\operatorname{gr}^{-1}\mathfrak{g}=H^{0}(\operatorname{Spec}A,\operatorname{gr}^{-1}\mathfrak{g}_{\mathcal{E}_{B}})\otimes\omega_{O}/t\omega_{O}.
\]
这样我们验证了 3.1.3 的条件 1；还需在 $\operatorname{Spec}A$ 的某一点上验证条件 2，例如在 $0\in\operatorname{Spec}A$ 上。以 $(\widetilde{\mathcal{E}}^{0}_{G},\widetilde{\mathcal{E}}^{0}_{B},\nabla)$ 记 $(\widetilde{\mathcal{E}}_{G},\widetilde{\mathcal{E}}_{B},\nabla)$ 到 $\{0\}\times\operatorname{Spec}O\subset\operatorname{Spec}(A\widehat\otimes O)$ 的限制。则 $\widetilde{\mathcal{E}}^{0}_{G}$ 是平凡 $G$-丛，$\nabla$ 是平凡联络，$\mathfrak{sl}_{2}$ 经 3.5.7 提到的态射 $\sigma:\mathfrak{sl}_{2}\to\mathfrak{g}$ 与嵌入 $\mathfrak{sl}_{2}=\mathbb{C}L_{-1}+\mathbb{C}L_{0}+\mathbb{C}L_{1}\hookrightarrow\operatorname{Der}O$ 作用在 $(\widetilde{\mathcal{E}}^{0}_{G},\nabla)$ 上，且 $\widetilde{\mathcal{E}}^{0}_{B}$ 关于 $\mathfrak{sl}_{2}$ 不变。由于 $\sigma$ 是主嵌入，$(\widetilde{\mathcal{E}}^{0}_{G},\widetilde{\mathcal{E}}^{0}_{B},\nabla)$ 是 $0\in\operatorname{Op}_{\mathfrak{g}}(O)$ 对应的 oper。

下面证 3)。置 $\mathfrak{a}=H^{0}(\operatorname{Spec}A,\mathfrak{g}_{\mathcal{E}})$，$\mathfrak{a}_{k}:=\{a\in\mathfrak{a}\mid L_{0}a=ka\}$。若固定 $\mathcal{E}$ 上的 $B$-结构，则 3.1.1 的滤链 $\mathfrak{g}^{k}$ 在 $\mathfrak{a}$ 上诱导滤链 $\mathfrak{a}^{k}$。若该 $B$-结构具有 b) 所述性质，则 $\mathfrak{a}^{k}$ 是 $\operatorname{Aut}^{0}O$-不变的，且 $\mathfrak{a}^{k}/\mathfrak{a}^{k+1}$ 作为 $\operatorname{Aut}^{0}O$-模同构于 $A\otimes(\omega_{O}/t\omega_{O})^{\otimes k}\otimes\operatorname{gr}^{k}\mathfrak{g}$（见 (72)）。因此 $L_{0}$ 在 $\mathfrak{a}^{k}/\mathfrak{a}^{k+1}$ 上的特征值 $\ge-k$，且 $A$-模 $\mathfrak{a}^{k}/\mathfrak{a}^{k+1}$ 由其对应于特征值 $-k$ 的 $L_{0}$-特征向量生成。故
\begin{equation}\tag{73}
\mathfrak{a}^{k}=\sum_{i\le-k}A\mathfrak{a}_{i}.
\end{equation}
$\mathcal{E}$ 上的 $B$-结构可由 Borel 子代数 $\mathfrak{a}^{0}\subset\mathfrak{a}$ 重构。

剩下从 c) 推出 b)。按 (73) 定义 $\mathfrak{a}^{k}$。由于 $\mathfrak{a}$ 是有限型的自由 $L_{0}$-分次 $A$-模，$\mathfrak{a}^{k}/\mathfrak{a}^{k+1}$ 也是。因此 $\mathfrak{a}^{k}$ 定义 $\mathfrak{g}_{\mathcal{E}}$ 的一个向量\SourcePage{87}%
子丛。若 $k=0$，该子丛是李子代数，故它定义截面 $s:\operatorname{Spec}A\to S_{\mathcal{E}}$，其中 $S$ 是 $\mathfrak{g}$ 的子代数概形。一个无穷小计算表明，态射 $G/B\to S$，$g\mapsto g\mathfrak{b}g^{-1}$，是开嵌入；又因 $G/B$ 射影，它也是闭嵌入。由 c)，$s(0)\in(G/B)_{\mathcal{E}}\subset S_{\mathcal{E}}$，故 $s(\operatorname{Spec}A)\subset(G/B)_{\mathcal{E}}$，于是 $s$ 定义 $\mathcal{E}$ 上的 $B$-结构。显然它是 $\operatorname{Aut}^{0}O$-不变的。$\operatorname{Spec}A$ 上相应的 $\operatorname{Aut}^{0}O$-等变 $H$-丛是 $\operatorname{Spec}\mathbb{C}$ 上某个 $\operatorname{Aut}^{0}O$-等变 $H$-丛 $F$ 的拉回（对 $\operatorname{Spec}A$ 上任意 $\operatorname{Aut}^{0}O$-等变 $H$-丛与任意环面 $H$ 都成立；事实上可设 $H=\mathbb{G}_{m}$，把 $\mathbb{G}_{m}$-丛解释为线丛，并利用有限生成分次投射 $A$-模自由这一事实）。为求 $F$，看 $0\in\operatorname{Spec}A$ 上的情形。

\emph{注.} 命题 3.5.8 的证明表明：若 c) 成立，则 $\mathcal{E}$ 上存在唯一的 $\operatorname{Aut}^{0}O$-不变 $B$-结构。

\bdp{3.5.9} \emph{推论.} 若 $G$ 是伴随群，则偶 $(\operatorname{Op}_{\mathfrak{g}}(O),\mathfrak{F}^{0}_{G})$ 没有非平凡的 $\operatorname{Aut}O$-等变自同构。
这是命题 3.5.8 的陈述 2) 取 $A=A_{\mathfrak{g}}(O)$ 的情形。

\bdp{3.5.10} 回忆：交换 $\mathbb{C}$-代数 $R$ 上的李代数胚，是李 $\mathbb{C}$-代数 $\mathfrak{a}$ 配以 $R$-模结构及映射 $\varphi:\mathfrak{a}\to\operatorname{Der}R$，使得 1) $\varphi$ 是李代数态射与 $R$-模态射；2) 对 $a_{1},a_{2}\in\mathfrak{a}$ 与 $f\in R$，有 $[a_{1},fa_{2}]=f[a_{1},a_{2}]+v(f)a_{2}$，其中 $v:=\varphi(a_{1})$。

\emph{注.}
(i) [Ma87] 与 [Ma96] 是李代数胚与李群胚的标准文献。另见 [We] 与 [BB93]。本文只需李代数胚的定义。
(ii) 李代数胚也称为 $(\mathbb{C},R)$-李代数（见 [R]），并有各种其他名称（见 [Ma96]）。

\bdp{3.5.11} 以 $\mathfrak{a}_{\mathfrak{g}}$ 记 $\mathfrak{F}^{0}_{G}$ 的（整体）无穷小对称空间。$\mathfrak{a}_{\mathfrak{g}}$ 的元素是偶，由 $\operatorname{Op}_{\mathfrak{g}}(O)=\operatorname{Spec}A_{\mathfrak{g}}(O)$ 上的向量场\SourcePage{88}%
（即 $A_{\mathfrak{g}}(O)$ 的导子）及其到主 $G$-丛 $\mathfrak{F}^{0}_{G}$ 上 $G$-不变向量场的提升组成。$\mathfrak{a}_{\mathfrak{g}}$ 是 $A_{\mathfrak{g}}(O)$ 上的李代数胚。我们有典范正合列
\[
0\to\mathfrak{g}_{\mathrm{univ}}\to\mathfrak{a}_{\mathfrak{g}}\to\operatorname{Der}A_{\mathfrak{g}}(O)\to0,
\]
其中 $\mathfrak{g}_{\mathrm{univ}}$ 是 $\mathfrak{g}$ 的 $\mathfrak{F}^{0}_{G}$-扭的整体截面空间。当然，把 $G$ 换成伴随群 $G_{ad}$ 时，$\mathfrak{a}_{\mathfrak{g}}$ 与 $\mathfrak{g}_{\mathrm{univ}}$ 不变。故 $\mathfrak{a}_{\mathfrak{g}}$ 与 $\mathfrak{g}_{\mathrm{univ}}$ 不依赖 $\omega_{O}^{1/2}$ 的选取。

$\operatorname{Der}O$ 在 $\mathfrak{F}^{0}_{G}$ 上的作用诱导李代数态射 $\operatorname{Der}O\to\mathfrak{a}_{\mathfrak{g}}$。特别地，$\operatorname{Der}O$ 作用于 $\mathfrak{a}_{\mathfrak{g}}$。

\bdp{3.5.12} \emph{引理.} $\mathfrak{a}_{\mathfrak{g}}$ 在 $\mathfrak{g}_{\mathrm{univ}}$ 上的伴随表示定义了 $\mathfrak{a}_{\mathfrak{g}}$ 与 $\mathfrak{g}_{\mathrm{univ}}$ 的无穷小对称代数胚之间的同构。

\bdp{3.5.13} \emph{命题.} 偶 $(A_{\mathfrak{g}}(O),\mathfrak{a}_{\mathfrak{g}})$ 的 $\operatorname{Der}O$-等变自同构群等于 $\operatorname{Aut}\Gamma$，其中 $\Gamma$ 是 $\mathfrak{g}$ 的 Dynkin 图。

\emph{证明.} 设 $G$ 是 $\mathfrak{g}$ 对应的伴随群。以 $L$ 记 $(A_{\mathfrak{g}}(O),\mathfrak{g}_{\mathrm{univ}})$ 的 $\operatorname{Der}O$-等变自同构群。由 3.5.12，须证 $L=\operatorname{Aut}\Gamma$。有显然态射 $i:\operatorname{Aut}\Gamma=\operatorname{Aut}(G,B)/B\to L$ 与 $\pi:L\to\operatorname{Aut}\Gamma$，使得 $\pi i=\mathrm{id}$。$\operatorname{Ker}\pi$ 是 $(\operatorname{Op}_{\mathfrak{g}}(O),\mathfrak{F}^{0}_{G})$ 的 $\operatorname{Der}O$-等变自同构群；故由 3.5.9，$\operatorname{Ker}\pi$ 平凡。\hfill$\Box$

\bdp{3.5.14} \emph{命题.} 偶 $(A_{\mathfrak{g}}(O),\mathfrak{a}_{\mathfrak{g}})$ 没有能诱导 $\operatorname{gr}A_{\mathfrak{g}}(O)$ 平凡自同构的非平凡 $\operatorname{Der}O$-等变自同构（$\operatorname{gr}$ 对应 3.2.1 的滤链）。

\emph{证明.} 设 $\Gamma$ 是 $\mathfrak{g}$ 的 Dynkin 图。由 3.5.13 与 (48)，须证 $\operatorname{Aut}\Gamma$ 在 2.7.1 的代数 $\mathfrak{z}^{cl}_{{}^{L}\mathfrak{g}}(O)$ 上的作用忠实。故只需证 $\operatorname{Aut}\Gamma$ 在 $W\backslash\mathfrak{h}$ 上的作用忠实（$W$ 表示 Weyl 群）。设 $C\subset\operatorname{Aut}\mathfrak{h}$ 是\SourcePage{89}%
根系的自同构群。存在 $a\in\mathfrak{h}$，其在 $C$ 中的稳定化子平凡。于是 $\operatorname{Aut}\Gamma=C/W$ 在 $W\backslash\mathfrak{h}$ 上的作用忠实。\hfill$\Box$

\bdp{3.5.15} 给 $\mathfrak{a}_{\mathfrak{g}}$ 配备最弱的平移不变拓扑，使得 $\mathfrak{F}^{0}_{G}$ 全空间上任一正则函数的稳定化子为开集（回忆 $\mathfrak{a}_{\mathfrak{g}}$ 作用于 $\mathfrak{F}^{0}_{G}$）。这是使 $\mathfrak{g}_{\mathrm{univ}}$ 每个元素的 $\mathfrak{a}_{\mathfrak{g}}$-中心化子为开集的最弱平移不变拓扑。于是该拓扑可由 $\mathfrak{a}_{\mathfrak{g}}$ 上的李代数胚结构重构。显然典范态射 $\operatorname{Der}O\to\mathfrak{a}_{\mathfrak{g}}$ 连续。

\bdp{3.5.16} 以 $\mathfrak{a}_{\mathfrak{b}}$ 记 $\mathfrak{F}^{0}_{B}$ 的（整体）无穷小对称李代数胚。以 $\mathfrak{b}_{\mathrm{univ}}$（相应地 $\mathfrak{n}_{\mathrm{univ}}$）记 $\mathfrak{b}$（相应地 $\mathfrak{n}$）的 $\mathfrak{F}^{0}_{B}$-扭的整体截面空间。有典范正合列
\[
0\to\mathfrak{b}_{\mathrm{univ}}\to\mathfrak{a}_{\mathfrak{b}}\to\operatorname{Der}A_{\mathfrak{g}}(O)\to0.
\]
$\mathfrak{a}_{\mathfrak{b}}$ 是 $\mathfrak{a}_{\mathfrak{g}}$ 的子代数胚；事实上 $\mathfrak{a}_{\mathfrak{b}}$ 是 $\mathfrak{b}_{\mathrm{univ}}\subset\mathfrak{a}_{\mathfrak{g}}$ 的正规化子。$\operatorname{Der}^{0}O$ 在 $\mathfrak{a}_{\mathfrak{g}}$ 中的像含于 $\mathfrak{a}_{\mathfrak{b}}$。

$\mathfrak{n}_{\mathrm{univ}}$ 是 $\mathfrak{a}_{\mathfrak{b}}$ 的理想，$\mathfrak{a}_{\mathfrak{b}}/\mathfrak{n}_{\mathrm{univ}}$ 是 $\mathfrak{F}^{0}_{H}$ 的（整体）无穷小对称代数胚。由于 $\mathfrak{F}^{0}_{H}$ 平凡且其平凡化「几乎」唯一（见 3.5.5），$\mathfrak{a}_{\mathfrak{b}}/\mathfrak{n}_{\mathrm{univ}}$ 典范同构于 $\operatorname{Der}A_{\mathfrak{g}}(O)$ 与 $A_{\mathfrak{g}}(O)\otimes\mathfrak{h}$ 的半直和。以 $\mathfrak{a}_{\mathfrak{n}}$ 记 $\operatorname{Der}A_{\mathfrak{g}}(O)\subset\mathfrak{a}_{\mathfrak{b}}/\mathfrak{n}_{\mathrm{univ}}$ 在 $\mathfrak{a}_{\mathfrak{b}}$ 中的原像。

\emph{注.} 由 3.5.5，复合 $\operatorname{Der}^{0}O\to\mathfrak{a}_{\mathfrak{b}}/\mathfrak{n}_{\mathrm{univ}}=\operatorname{Der}A_{\mathfrak{g}}(O)\oplus(A_{\mathfrak{g}}(O)\otimes\mathfrak{h})$ 的像含于 $\operatorname{Der}A_{\mathfrak{g}}(O)\oplus\mathfrak{h}$；它等于自然态射 $\operatorname{Der}^{0}O\to\operatorname{Der}A_{\mathfrak{g}}(O)$ 与态射 $\operatorname{Der}^{0}O\to\mathfrak{h}$（$L_{0}\mapsto-\check\rho$，$n>0$ 时 $L_{n}\mapsto0$）之和。

\bdp{3.5.17} 下面用 $L_{0}$ 在 $\mathfrak{a}_{\mathfrak{g}}$ 上的作用描述 $\mathfrak{a}_{\mathfrak{b}}$、$\mathfrak{b}_{\mathrm{univ}}$ 等。将使用如下记号。若 $\operatorname{Der}O$ 作用在拓扑向量空间 $V$ 上，使 $L_{0}:V\to V$ 的特征值为整数，则以 $V^{\le k}$ 记 $V$ 中含所有满足 $L_{0}v=nv$、$n\le k$ 的 $v$ 的最小闭子空间。置 $V^{<k}:=V^{\le k-1}$。若 $V$ 是某代数 $A$ 上的拓扑模，\SourcePage{90}%
$W$ 是 $V$ 的子空间，则以 $A\cdot W$ 记 $V$ 中含所有 $aw$（$a\in A$，$w\in W$）的最小闭子空间。

\bdp{3.5.18} \emph{命题.} i) 下列等式成立：
\begin{equation}\tag{74}
\mathfrak{b}_{\mathrm{univ}}=A_{\mathfrak{g}}(O)\cdot(\mathfrak{g}_{\mathrm{univ}})^{\le0}
\end{equation}
\begin{equation}\tag{75}
\mathfrak{n}_{\mathrm{univ}}=A_{\mathfrak{g}}(O)\cdot\mathfrak{g}_{\mathrm{univ}}^{<0}
\end{equation}
\begin{equation}\tag{76}
\mathfrak{a}_{\mathfrak{b}}=A_{\mathfrak{g}}(O)\cdot(\mathfrak{a}_{\mathfrak{g}})^{\le0}
\end{equation}
\begin{equation}\tag{77}
\mathfrak{a}_{\mathfrak{n}}=A_{\mathfrak{g}}(O)\cdot\mathfrak{a}_{\mathfrak{g}}^{<0}
\end{equation}

ii) 态射
\[
\begin{aligned}
(\mathfrak{a}_{\mathfrak{g}})^{\le0}&\longrightarrow
A_{\mathfrak{g}}(O)(\mathfrak{a}_{\mathfrak{g}})^{\le0}/A_{\mathfrak{g}}(O)\mathfrak{a}_{\mathfrak{g}}^{<0}\\
&=\mathfrak{a}_{\mathfrak{b}}/\mathfrak{a}_{\mathfrak{n}}=A_{\mathfrak{g}}(O)\otimes\mathfrak{h}
\end{aligned}
\]
的像等于 $\mathfrak{h}$，故有典范同构
\begin{equation}\tag{78}
(\mathfrak{a}_{\mathfrak{g}})^{\le0}\big/\bigl(A_{\mathfrak{g}}(O)\cdot\mathfrak{a}_{\mathfrak{g}}^{<0}\cap(\mathfrak{a}_{\mathfrak{g}})^{\le0}\bigr)\xrightarrow{\sim}\mathfrak{h}.
\end{equation}

\emph{证明.} i) (74)–(77) 由 (69) 推出。或者注意到 (74) 与 (75) 是 (73) 的特例，并如下证明 (76)。由 (74)，$A_{\mathfrak{g}}(O)\cdot(\mathfrak{a}_{\mathfrak{g}})^{\le0}\supset\mathfrak{b}_{\mathrm{univ}}$，且 $A_{\mathfrak{g}}(O)\cdot(\operatorname{Der}A_{\mathfrak{g}}(O))^{\le0}=\operatorname{Der}A_{\mathfrak{g}}(O)$，故 $A_{\mathfrak{g}}(O)\cdot(\mathfrak{a}_{\mathfrak{g}})^{\le0}\supset\mathfrak{a}_{\mathfrak{b}}$。又由 (74)，$(\mathfrak{a}_{\mathfrak{g}}/\mathfrak{a}_{\mathfrak{b}})^{\le0}=(\mathfrak{g}_{\mathrm{univ}}/\mathfrak{b}_{\mathrm{univ}})^{\le0}=(\mathfrak{g}_{\mathrm{univ}})^{\le0}/(\mathfrak{b}_{\mathrm{univ}})^{\le0}=0$，故 $A_{\mathfrak{g}}(O)\cdot(\mathfrak{a}_{\mathfrak{g}})^{\le0}\subset\mathfrak{a}_{\mathfrak{b}}$。
ii) $(\mathfrak{a}_{\mathfrak{g}})^{\le0}$ 在 $A_{\mathfrak{g}}(O)\otimes\mathfrak{h}$ 中的像等于 $(A_{\mathfrak{g}}(O)\otimes\mathfrak{h})^{\le0}=\mathfrak{h}$。\hfill$\Box$

\bdsec{3.6. Feigin--Frenkel 同构 II}
\bdp{3.6.1} 设 $A$ 是 $\mathbb{C}[h]$ 上的结合代数，且作为 $\mathbb{C}[h]$-模平坦。置 $A_{0}:=A/hA$。以 $\mathfrak{z}$ 记 $A_{0}$ 的中心。若 $\mathfrak{z}=A_{0}$，即 $A_{0}$ 交换，则 $\mathfrak{z}$ 带有标准泊松括号
\begin{equation}\tag{79}
\{z_{1},z_{2}\}:=[\tilde z_{1},\tilde z_{2}]/h\bmod h.
\end{equation}

% ---- ch03_u04 ----
% ch03_u04 —— 原书 p.91–p.100（3.6.1 末、3.6.2–3.6.19、3.7 开头至 3.7.5）
\SourcePage{91}

其中 $z_{1}, z_{2} \in \mathfrak{z}$，而 $\tilde{z}_{i}$ 是 $z_{i}$ 在 $A$ 中的原像。Hayashi 在 [Ha88] 中指出：即使去掉假设 $\mathfrak{z} = A_{0}$，（79）仍是 $\mathfrak{z}$ 上的良定义泊松括号。特别地，（79）的右端属于 $\mathfrak{z}$。

\emph{注记}
\begin{enumerate}[label=(\roman*),leftmargin=2.2em]
  \item 上述情形有典范的李代数态射 $\varphi: \mathfrak{z} \to \operatorname{Der} A_{0}/\operatorname{Int} A_{0}$，其中 $\operatorname{Int} A_{0}$ 是内导子构成的空间。$\varphi$ 由 $\varphi(z) = D_{z}$ 定义，其中 $D_{z}(a) := [\tilde{z}, \tilde{a}]/h \bmod h$，而 $\tilde{z}, \tilde{a} \in A$ 分别是 $z \in \mathfrak{z}$ 与 $a \in A_{0}$ 的原像。若 $z' \in \mathfrak{z}$，则 $D_{z}(z') = \{z, z'\}$。$\operatorname{Der} A_{0}/\operatorname{Int} A_{0}$ 是 $\mathfrak{z}$-模，且 $\varphi(z_{1}z_{2}) = z_{1}\varphi(z_{2}) + z_{2}\varphi(z_{1})$。于是 $\varphi$ 诱导 $\mathfrak{z}$-模态射 $\Phi: \Omega^{1}_{\mathfrak{z}} \to \operatorname{Der} A_{0}/\operatorname{Int} A_{0}$。事实上 $\Phi$ 是 $\mathfrak{z}$ 上李代数胚的态射（李代数胚的定义见 3.5.10）。$\operatorname{Der} A_{0}/\operatorname{Int} A_{0}$ 上的李代数胚结构按显然方式定义；$\Omega^{1}_{\mathfrak{z}}$ 上的结构则是泊松括号诱导的标准代数胚结构（参看 [We88]）。即对 $z, z' \in \mathfrak{z}$ 令 $[dz, dz'] := d\{z, z'\}$；而态射 $\Omega^{1}_{\mathfrak{z}} \to \operatorname{Der} \mathfrak{z}$ 把 $dz$ 映到 $\operatorname{grad} z$，$(\operatorname{grad} z)(z') := \{z, z'\}$。
  \item 若把 $\mathbb{C}[h]$ 换成 $\mathbb{C}[h]/(h^{3})$，上述构造仍有意义。
\end{enumerate}

\bdp{3.6.2} 现在设 $\mathfrak{g}$ 是半单李代数，$K := \mathbb{C}((t))$。记 $A$ 为 2.5.1 中李代数 $\widetilde{\mathfrak{g}\otimes K}$ 的完备化泛包络代数，即 $A := \varprojlim_{n}(U\widetilde{\mathfrak{g}\otimes K})/J_{n}$，其中 $J_{n} \subset U\widetilde{\mathfrak{g}\otimes K}$ 是由 $\mathfrak{g}\otimes t^{n}\mathbb{C}[[t]] \subset \mathfrak{g}\otimes K \subset \widetilde{\mathfrak{g}\otimes K}$ 生成的左理想，$n \ge 0$。考虑 $A$ 上由 $ha = \mathbf{1}\cdot a - a$（$a \in A$）定义的 $\mathbb{C}[h]$-代数结构，其中 $\mathbf{1} \in \mathbb{C} \subset \widetilde{\mathfrak{g}\otimes K} \subset A$。$A$ 在 $\mathbb{C}[h]$ 上平坦，而 $A/hA$ 是 2.5.2 与 2.9.4 中的完备化扭泛包络代数 $\overline{U}' = \overline{U}'(\mathfrak{g}\otimes K)$。于是（79）定义了 $\overline{U}'$ 的中心 $\mathfrak{z}$ 上的泊松括号。[Ha88] 引入了这个括号，故称之为 Hayashi 括号。

\bdp{3.6.3} 对开的李子代数 $\mathfrak{a} \subset \mathfrak{g}\otimes O$，记 $\mathcal{J}_{\mathfrak{a}}$（相应地 $\widetilde{\mathcal{J}}_{\mathfrak{a}}$）为 $\overline{U}'$（相应地 $A = \overline{U}\widetilde{\mathfrak{g}\otimes K}$）中由 $\mathfrak{a} \subset \mathfrak{g}\otimes O \subset \mathfrak{g}\otimes K$ 生成的左理想的闭包。显然 $\mathcal{J}_{\mathfrak{a}}$ 是 $\widetilde{\mathcal{J}}_{\mathfrak{a}}$ 在 $\overline{U}'$ 中的像。令 $I_{\mathfrak{a}} := \mathcal{J}_{\mathfrak{a}} \cap \mathfrak{z}$。我们给 $\mathfrak{z}$ 装上由 $\overline{U}'$ 诱导的拓扑。

\SourcePage{92}

理想 $\mathcal{J}_{\mathfrak{a}}$（相应地 $I_{\mathfrak{a}}$）构成 $\overline{U}'$（相应地 $\mathfrak{z}$）中零的邻域基。

\bdp{3.6.4} \emph{引理.}
\begin{enumerate}[label=(\roman*),leftmargin=2.2em]
  \item $\{I_{\mathfrak{a}}, I_{\mathfrak{a}}\} \subset I_{\mathfrak{a}}$。
  \item $\mathfrak{z}$ 上的 Hayashi 括号连续。
\end{enumerate}

\emph{证明.} 用到如下事实：$A/\widetilde{\mathcal{J}}_{\mathfrak{a}}$ 带 3.6.2 的 $\mathbb{C}[h]$-模结构时平坦。\hfill$\Box$

\bdp{3.6.5} 令 $I := I_{\mathfrak{g}\otimes O}$。典范态射 $\mathfrak{z} \to \mathfrak{z}_{\mathfrak{g}}(O)$ 是满射（见 2.9.3--2.9.5），其核等于 $I$。于是 $\mathfrak{z}_{\mathfrak{g}}(O) = \mathfrak{z}/I$。

记 $I^{2}$ 为 $\mathfrak{z}$ 中由形如 $ab$（$a, b \in I$）的元素生成的闭理想。则 $I/I^{2}$ 是 $\mathfrak{z}_{\mathfrak{g}}(O)$ 上的李代数胚：交换子 $I/I^{2} \times I/I^{2} \to I/I^{2}$ 与映射 $I/I^{2} \to \operatorname{Der}\mathfrak{z}_{\mathfrak{g}}(O)$ 都由 Hayashi 括号诱导。李代数 $\operatorname{Der} O$ 作用在 $I/I^{2}$ 与 $\mathfrak{z}_{\mathfrak{g}}(O)$ 上。这些作用连续：$I/I^{2}$ 带由 $\mathfrak{z}$ 诱导的拓扑，$\mathfrak{z}_{\mathfrak{g}}(O)$ 取离散拓扑。

\bdp{3.6.6} 我们来表述定理 3.2.2 的一个更精确的版本。已有代数 $\mathfrak{z}_{\mathfrak{g}}(O)$ 及其上的李代数胚 $I/I^{2}$。另一方面，记 ${}^{L}\mathfrak{g}$ 为朗兰兹对偶，考虑代数 $A_{{}^{L}\mathfrak{g}}(O)$（见 3.2.1）及其上的李代数胚 $\mathfrak{a}_{{}^{L}\mathfrak{g}}$（见 3.5.11）。$I/I^{2}$ 与 $\mathfrak{a}_{{}^{L}\mathfrak{g}}$ 都带拓扑（见 3.6.5 与 3.5.15）。李代数 $\operatorname{Der} O$ 作用在所有这些对象上。$\mathfrak{z}_{\mathfrak{g}}(O)$ 与 $A_{{}^{L}\mathfrak{g}}(O)$ 都带滤链（见 1.2.5 与 3.2.1）。此外有态射 $\sigma_{A}^{-1}\sigma_{\mathfrak{z}}: \operatorname{gr}\mathfrak{z}_{\mathfrak{g}}(O) \to \operatorname{gr} A_{{}^{L}\mathfrak{g}}(O)$，其中 $\sigma_{\mathfrak{z}}: \operatorname{gr}\mathfrak{z}_{\mathfrak{g}}(O) \to \mathfrak{z}^{cl}_{\mathfrak{g}}(O)$ 是象征映射，而 $\sigma_{A}$ 是把 $\mathfrak{g}$ 换成 ${}^{L}\mathfrak{g}$ 后的同构（48）。

\bdp{3.6.7} \emph{定理.} 存在滤 $\operatorname{Der} O$-代数的同构
\begin{equation}\tag{80}
\varphi_{O}: A_{{}^{L}\mathfrak{g}}(O) \xrightarrow{\sim} \mathfrak{z}_{\mathfrak{g}}(O)
\end{equation}

\SourcePage{93}
使得 $\operatorname{gr}\varphi_{O}^{-1} = \sigma_{A}^{-1}\sigma_{\mathfrak{z}}$；并且 $\varphi_{O}$ 可延拓为李代数胚的拓扑 $\operatorname{Der} O$-等变同构
\begin{equation}\tag{81}
\mathfrak{a}_{{}^{L}\mathfrak{g}} \xrightarrow{\sim} I/I^{2}.
\end{equation}

本定理可从 [FF92] 中提取（见 3.7.12--3.7.17）。

\emph{注.} 由 3.5.14，同构（80）与（81）唯一。3.6.11 中我们将给出同构（81）的一条附加性质。但首先须为代数胚 $I/I^{2}$ 定义（78）的类比。

\bdp{3.6.8} 沿用 3.5.17 的记号。

\emph{引理.} 令 $\mathcal{J}_{-} := (\overline{U}')^{\le 0} \cap \mathcal{J}_{\mathfrak{a}}$，其中 $\mathfrak{a} = t\mathfrak{g}[[t]]$，而 $\mathcal{J}_{\mathfrak{a}}$ 见 3.6.3。则 $\mathcal{J}_{-}$ 是 $(\overline{U}')^{\le 0}$ 的双边理想，并且
\begin{equation}\tag{82}
(\overline{U}')^{\le 0} = U\mathfrak{g} \oplus \mathcal{J}_{-}.
\end{equation}

\emph{证明.} （82）显然。由于 $\mathcal{J}_{-}$ 是左理想且 $[\mathfrak{g}, \mathcal{J}_{-}] \subset \mathcal{J}_{-}$，（82）表明 $\mathcal{J}_{-}$ 是双边理想。\hfill$\Box$

定义 $\pi: (\overline{U}')^{\le 0} \to U\mathfrak{g}$ 为满足 $\pi(\mathcal{J}_{-}) = 0$ 且 $\pi(a) = a$（$a \in U\mathfrak{g}$）的态射。

这里给出 $\pi$ 的等价定义。令 $\operatorname{Vac}'_{\mathfrak{a}} := \overline{U}'/\mathcal{J}_{\mathfrak{a}}$，$\mathfrak{a} = t\mathfrak{g}[[t]]$。则 $\operatorname{Vac}'_{\mathfrak{a}}$ 是左 $\overline{U}'$-模，也是右 $U\mathfrak{g}$-模。$L_{0}$ 在 $\operatorname{Vac}'_{\mathfrak{a}}$ 上的特征值非负，且 $\operatorname{Ker}(L_{0}: \operatorname{Vac}'_{\mathfrak{a}} \to \operatorname{Vac}'_{\mathfrak{a}}) = U\mathfrak{g}$。于是 $U\mathfrak{g} \subset \operatorname{Vac}'_{\mathfrak{a}}$ 在 $(\overline{U}')^{\le 0}$ 的左作用下不变。$(\overline{U}')^{\le 0}$ 的左作用与 $U\mathfrak{g}$ 的右作用交换，故它定义了态射 $(\overline{U}')^{\le 0} \to U\mathfrak{g}$。这就是 $\pi$。

\bdp{3.6.9} 记 $C$ 为 $U\mathfrak{g}$ 的中心。则
\[
  \pi(\mathfrak{z}^{\le 0}) \subset C, \qquad \pi(\mathfrak{z}\cdot\mathfrak{z}^{<0} \cap \mathfrak{z}^{\le 0}) = 0.
\]
设 $m \subset C$ 是 $U\mathfrak{g}$ 的单位表示对应的极大理想。回顾 $I := \operatorname{Ker}(\mathfrak{z} \to \mathfrak{z}_{\mathfrak{g}}(O))$。则 $\pi(I^{\le 0}) \subset m$。由于 $(I^{2})^{\le 0} \subset I^{\le 0}\cdot I^{\le 0} + (\mathfrak{z}\cdot\mathfrak{z}^{<0} \cap \mathfrak{z}^{\le 0})$，有 $\pi((I^{2})^{\le 0}) \subset m^{2}$。于是 $\pi$ 诱导 $\mathbb{C}$-线性映射 $d: (I/I^{2})^{\le 0} \to m/m^{2}$，使得 $\mathfrak{z}_{\mathfrak{g}}(O)\cdot(I/I^{2})^{<0} \cap (I/I^{2})^{\le 0} \subset \operatorname{Ker} d$。

\SourcePage{94}

\emph{练习.} 对 $z_{1}, z_{2} \in \mathfrak{z}^{\le 0}$ 有 $\pi(\{z_{1}, z_{2}\}) = 0$（于是 $d$ 是李代数态射）。

\bdp{3.6.10} 把 $C$ 与 $\mathfrak{h}^{*}$ 上 $W$-不变多项式的代数等同，其中 $W$ 是 Weyl 群。则 $m$ 由 $\mathfrak{h}^{*}$ 上在 $\rho$ 处为零的 $W$-不变多项式构成，其中 $\rho :=$ 基本权之和。由于 $\rho \in \mathfrak{h}^{*}$ 是正则的，可把 $m/m^{2}$ 与 $\mathfrak{h}$ 等同：把 $m$ 中的 $W$-不变多项式对应到它在 $\rho$ 处的微分。于是我们构造了映射
\begin{equation}\tag{83}
d: (I/I^{2})^{\le 0}/\bigl(\mathfrak{z}_{\mathfrak{g}}(O)\cdot(I/I^{2})^{<0} \cap (I/I^{2})^{\le 0}\bigr) \to \mathfrak{h}
\end{equation}

\bdp{3.6.11} \emph{定理.} 下图反交换：
\begin{equation}\tag{84}
\begin{array}{ccc}
(\mathfrak{a}_{{}^{L}\mathfrak{g}})^{\le 0}/\bigl(A_{{}^{L}\mathfrak{g}}(O)\cdot\mathfrak{a}_{{}^{L}\mathfrak{g}}^{<0} \cap (\mathfrak{a}_{{}^{L}\mathfrak{g}})^{\le 0}\bigr) & \xrightarrow{\ \sim\ } & \mathfrak{h}^{*}\\[2mm]
\big\downarrow\ \cong & & \big\uparrow\ \cong\\[2mm]
d: (I/I^{2})^{\le 0}/\bigl(\mathfrak{z}_{\mathfrak{g}}(O)\cdot(I/I^{2})^{<0} \cap (I/I^{2})^{\le 0}\bigr) & \longrightarrow & \mathfrak{h}
\end{array}
\end{equation}
这里上方箭头是把 $\mathfrak{g}$ 换成 ${}^{L}\mathfrak{g}$ 后的同构（78），左方箭头由（81）诱导，右方箭头来自数量积（18）。

本定理可从 [FF92] 中提取（见 3.8.15--3.8.22）。

\bdp{3.6.12} 「临界」数量积（18）在 3.6.11 中出现，原因并不严重。读者或许更喜欢下述观点。记 $B$ 为 $\mathfrak{g}$ 上不变双线性型的集合。对每个 $b \in B$，有完备化扭泛包络代数 $\overline{U}'_{b} = \overline{U}'_{b}(\mathfrak{g}\otimes K)$，它对应余圈 $(u, v) \mapsto \operatorname{Res} b(du, v)$，$u, v \in \mathfrak{g}\otimes K$（故 $\overline{U}' = \overline{U}'_{c}$，其中 $c$ 由（18）定义）。把 3.6.1 的一般构造施于依赖参数 $h$ 的代数族 $\overline{U}'_{c+hb}$，就得到与 $b \in B$ 对应的 $\mathfrak{z}$ 上的泊松括号 $\{\}_{b}$（3.6.2 的括号对应 $b = c$）。$I/I^{2}$ 上的李代数胚结构依赖于 $b$。于是对每个非退化 $b \in B$，3.6.7 与 3.6.11 都成立（注意（84）中两个竖直箭头都依赖于 $b$）。

\SourcePage{95}

\bdp{3.6.13} 事实上，3.6.6--3.6.7 中提到的 $\operatorname{Der} O$ 在 $I/I^{2}$ 上的作用，来自典范态射 $\operatorname{Der} O \to I$，它本质上属于 Sugawara。在 3.6.14--3.6.15 简要回顾 Sugawara 公式之后，我们将在 3.6.16 解释这一点。这些公式还给出 $\mathfrak{z}_{\mathfrak{g}}(O)$ 的元素；当 $\mathfrak{g} = \mathfrak{sl}_{2}$ 时，它们生成 $\mathfrak{z}_{\mathfrak{g}}(O)$。我们在 3.6.18 中回顾这一点。3.6.18 与 3.6.19 在后文中都用不到 (?)。
% [Tutor Note: 原书 p.95 末括注「(?)」为作者自记之疑问，底本如此，保留原样。]

\bdp{3.6.14} 本小节回顾一般的 Sugawara 公式，3.6.15 回顾它们在临界水平的推论。

设 $A$ 为 $\widetilde{\mathfrak{g}\otimes K}$ 的完备化泛包络代数。作为向量空间，$\widetilde{\mathfrak{g}\otimes K}$ 是 $\mathfrak{g}\otimes K$ 与 $\mathbb{C} = \mathbb{C}\cdot 1$ 的直和。Sugawara 元素 $\widetilde{\mathcal{L}}_{n} \in A$ 由下式定义：
\begin{equation}\tag{85}
\widetilde{\mathcal{L}}_{n} := \frac{1}{2}\sum_{r+l=n} g^{\lambda\mu} :e_{\lambda}^{(r)}e_{\mu}^{(l)}:
\end{equation}
其中 $\{e_{\lambda}\}$ 是 $\mathfrak{g}$ 的一组基，$e_{\lambda}^{(r)} := e_{\lambda}t^{r} \in \mathfrak{g}((t)) = \mathfrak{g}\otimes K \subset \widetilde{\mathfrak{g}\otimes K}$，$(g^{\lambda\mu})$ 是 Gram 矩阵 $(e_{\lambda}, e_{\mu})$ 关于「临界」数量积（18）的逆，并且
\begin{equation}\tag{86}
:e_{\lambda}^{(r)}e_{\mu}^{(l)}: \,=
\begin{cases}
e_{\lambda}^{(r)}e_{\mu}^{(l)} & \text{若 } r \le l,\\
e_{\mu}^{(l)}e_{\lambda}^{(r)} & \text{若 } r > l.
\end{cases}
\end{equation}
当然，（85）中对 $\lambda$ 与 $\mu$ 的求和是隐含的。显然无穷级数（85）收敛，且当 $n \to \infty$ 时 $\widetilde{\mathcal{L}}_{n} \to 0$。

\emph{注.} 若 $n \ne 0$，则（85）中的 $:e_{\lambda}^{(r)}e_{\mu}^{(l)}:$ 可换成 $e_{\lambda}^{(r)}e_{\mu}^{(l)}$。事实上，由于 $g^{\lambda\mu}$ 对称，除非 $r + l = 0$ 且 $r \ne 0$，总有 $g^{\lambda\mu}[e_{\lambda}^{(r)}, e_{\mu}^{(l)}] = 0$。

\SourcePage{96}

下述公式的证明可在文献中找到\footnote{读者应注意，Kac--Moody 代数的专家通常给 $\mathfrak{g}$ 配上把（18）除以负的对偶 Coxeter 数所得的数量积。}，例如见 [KR] 的 Lecture 10 与 [Kac90] 的 §12.8：
\begin{equation}\tag{87}
\operatorname{ad}\widetilde{\mathcal{L}}_{n} = hL_{n},
\end{equation}
\begin{equation}\tag{88}
L_{m}(\widetilde{\mathcal{L}}_{n}) = (m-n)\widetilde{\mathcal{L}}_{m+n} + \delta_{m,-n}\cdot\frac{m^{3}-m}{12}\cdot(\dim\mathfrak{g})\cdot\mathbf{1}.
\end{equation}
（87）中 $\operatorname{ad}\widetilde{\mathcal{L}}_{n}$ 是算子 $A \to A$；$L_{n} := -t^{n+1}\frac{d}{dt} \in \operatorname{Der} K$ 也看作算子 $A \to A$（李代数 $\operatorname{Der} K$ 按显然方式作用在 $A$ 上）；$h$ 的含义同 3.6.2，即 $h: A \to A$ 是乘以 $\mathbf{1} - 1$。

利用（87）可把（88）改写成 Virasoro 形式：
\begin{equation}\tag{89}
[\widetilde{\mathcal{L}}_{m}, \widetilde{\mathcal{L}}_{n}] = h\Bigl((m-n)\widetilde{\mathcal{L}}_{m+n} + \delta_{m,-n}\cdot\frac{m^{3}-m}{12}\cdot(\dim\mathfrak{g})\cdot\mathbf{1}\Bigr).
\end{equation}

\bdp{3.6.15} $\widetilde{\mathcal{L}}_{n}$ 在 $A/hA = \overline{U}'$ 中的像记作 $\mathcal{L}_{n}$。由（87），$\mathcal{L}_{n}$ 属于中心 $\mathfrak{z} \subset \overline{U}'$，并且
\begin{equation}\tag{90}
\{\mathcal{L}_{n}, z\} = L_{n}(z), \qquad z \in \mathfrak{z}
\end{equation}
其中 $\{\ \}$ 表示 $\mathfrak{z}$ 上的 Hayashi 括号。由（88）与（89）有
\begin{equation}\tag{91}
L_{m}(\mathcal{L}_{n}) = (m-n)\mathcal{L}_{m+n} + \delta_{m,-n}\cdot\frac{m^{3}-m}{12}\cdot\dim\mathfrak{g}
\end{equation}
\begin{equation}\tag{92}
\{\mathcal{L}_{m}, \mathcal{L}_{n}\} = (m-n)\mathcal{L}_{m+n} + \delta_{m,-n}\cdot\frac{m^{3}-m}{12}\cdot\dim\mathfrak{g}.
\end{equation}

\bdp{3.6.16} 若 $n \ge -1$，则 $\mathcal{L}_{n} \in I := \operatorname{Ker}(\mathfrak{z} \to \mathfrak{z}_{\mathfrak{g}}(O))$。事实上，看（85）即知 $\mathcal{L}_{n}$ 消灭 $\operatorname{Vac}'$ 中的真空向量。若 $m, n \ge -1$，则「Virasoro 项」$\delta_{m,-n}(m^{3}-m)$ 为零，故有连续李代数态射 $\operatorname{Der} O \to I$，由 $L_{n} \mapsto \mathcal{L}_{n}$（$n \ge -1$）定义。它诱导连续代数态射
\begin{equation}\tag{93}
\operatorname{Der} O \to I/I^{2}.
\end{equation}

\SourcePage{97}

\emph{注.} 由（90），（93）诱导的 $\operatorname{Der} O$ 在 $I/I^{2}$ 上的作用与 3.6.6--3.6.7 中考虑的作用一致。

\bdp{3.6.17} \emph{引理.} （93）与同构 $I/I^{2} \xrightarrow{\sim} \mathfrak{a}_{{}^{L}\mathfrak{g}}$（它是（81）之逆）的复合，等于 3.5.11 中的态射 $\operatorname{Der} O \to \mathfrak{a}_{{}^{L}\mathfrak{g}}$。

\emph{证明.} 这两个态射 $\operatorname{Der} O \to \mathfrak{a}_{{}^{L}\mathfrak{g}}$ 诱导 $\operatorname{Der} O$ 在 $\mathfrak{a}_{{}^{L}\mathfrak{g}}$ 上的同一个作用。故由 3.5.12 它们相等。\hfill$\Box$

\bdp{3.6.18} 记 $\overline{\mathcal{L}}_{n}$ 为 $\mathcal{L}_{n}$ 在 $\mathfrak{z}/I = \mathfrak{z}_{\mathfrak{g}}(O)$ 中的像。若 $n \ge -1$，则 $\overline{\mathcal{L}}_{n} = 0$。自然态射 $\mathbb{C}[\overline{\mathcal{L}}_{-2}, \overline{\mathcal{L}}_{-3}, \ldots] \to \mathfrak{z}_{\mathfrak{g}}(O)$ 是单射；当 $\mathfrak{g} = \mathfrak{sl}_{2}$ 时它是同构。为看出这一点，只需计算 $\overline{\mathcal{L}}_{n}$ 的主象征，并利用 2.4.1 对 $\mathfrak{z}^{cl}_{\mathfrak{g}}(O)$ 的描述。若把 $\mathfrak{z}^{cl}_{\mathfrak{g}}(O)$ 与 $\mathfrak{g}^{*}\otimes\omega_{O}$ 上 $G(O)$-不变多项式的空间等同（见 2.4.1），则 $\overline{\mathcal{L}}_{n}$ 的主象征是多项式 $\ell_{n}: \mathfrak{g}^{*}\otimes\omega_{O} \to \mathbb{C}$，由 $\ell_{n}(\eta) = \frac{1}{2}\operatorname{Res}(\eta, \eta)L_{n}$ 定义。这里 $(\eta, \eta) \in \omega_{O}^{\otimes 2}$，$L_{n} \in \omega_{K}^{\otimes(-1)}$，$(\eta, \eta)L_{n} \in \omega_{K}$，故留数有意义。显然映射 $\mathbb{C}[\ell_{-2}, \ell_{-3}, \ldots] \to \mathfrak{z}^{cl}_{\mathfrak{g}}(O)$ 是单射；当 $\mathfrak{g} = \mathfrak{sl}_{2}$ 时它是同构。

对 $\mathfrak{g} = \mathfrak{sl}_{2}$，Feigin--Frenkel 同构是唯一的 $\operatorname{Der} O$-等变同构 $A_{{}^{L}\mathfrak{g}}(O) \xrightarrow{\sim} \mathfrak{z}_{\mathfrak{g}}(O)$。$\operatorname{Spec} O$ 上的 $\mathfrak{sl}_{2}$-oper 可表示成联络 $\frac{d}{dt} + \begin{pmatrix} 0 & u \\ 1 & 0 \end{pmatrix}$，$u = u(t) = u_{0} + u_{1}t + \ldots$；也可表示成 Sturm--Liouville 算子 $(\frac{d}{dt})^{2} - u(t): \omega_{O}^{-1/2} \to \omega_{O}^{3/2}$。有 $A_{\mathfrak{sl}_{2}}(O) = \mathbb{C}[u_{0}, u_{1}, \ldots]$，而 Feigin--Frenkel 同构把 $u_{j}$ 映到 $-2\overline{\mathcal{L}}_{-2-j}$。

对任意半单 $\mathfrak{g}$，我们在 3.5.6 中给出了 $A_{{}^{L}\mathfrak{g}}(O)$ 作为带 $\operatorname{Der} O$ 作用的代数的描述；见（64）--（68）。利用 Feigin--Frenkel 同构的 $\operatorname{Der} O$-等变性可见：若 $\mathfrak{g}$ 单，则 $\overline{\mathcal{L}}_{-2-j} \in \mathfrak{z}_{\mathfrak{g}}(O)$ 对应 $cu_{1j} \in A_{{}^{L}\mathfrak{g}}(O)$，其中 $c = -(\dim\mathfrak{g})/6$ (???)。
% [Tutor Note: 原书 p.97 括注「(???)」为作者自记之疑问，且下标印作「1j」，底本如此，保留原样。]

\bdp{3.6.19} 考虑真空模 $\operatorname{Vac}_{\lambda} := \operatorname{Vac}_{A}/(h-\lambda)\operatorname{Vac}_{A}$，其中 $\operatorname{Vac}_{A}$ 是 $A$ 模去由 $\mathfrak{g}\otimes O$ 生成的闭左理想所得的商。2.9.3 中我们提到，当 $\lambda \ne 0$ 时 $\operatorname{End}_{A}\operatorname{Vac}_{\lambda} = \mathbb{C}$。下述证明

\SourcePage{98}
系由 E. Frenkel 告知我们。正如 2.9.3--2.9.5 所解释的，任何自同态 $f: \operatorname{Vac}_{\lambda} \to \operatorname{Vac}_{\lambda}$ 都来自 $A/(h-\lambda)A$ 的某个中心元素 $z$。事实上，当 $\lambda \ne 0$ 时 $A/(h-\lambda)A$ 的中心等于 $\mathbb{C}$；但我们不去证明这一点，转而注意 $[\widetilde{\mathcal{L}}_{0}, z] = 0$，从而 $L_{0}(z) = 0$（见（87））。于是 $[L_{0}, f] = 0$，其中 $L_{0}$ 看作 $\operatorname{Vac}_{\lambda}$ 中的算子。因此 $f$ 保持空间 $\operatorname{Ker}(L_{0}: \operatorname{Vac}_{\lambda} \to \operatorname{Vac}_{\lambda})$，该空间由真空向量生成。由于 $A$-模 $\operatorname{Vac}_{\lambda}$ 由这个空间生成，$f$ 是纯量算子。

\bdsec{3.7. 中心与 Gelfand--Dikii 括号}

\bdp{3.7.1} 令 $Y := \operatorname{Spec} O$，$Y' := \operatorname{Spec} K$，其中按惯例 $O = \mathbb{C}[[t]]$，$K = \mathbb{C}((t))$。设 $A$ 是（交换的）$\operatorname{Aut} O$-代数。则对任意光滑曲线 $X$ 可得 $\mathcal{D}_{X}$-代数 $A_{X}$（见 2.6.5）。$Y$ 与 $Y'$ 虽不是字面意义上的曲线，2.6.5 的构造对它们仍适用（只需作一点修改，见下）。于是得到 $\mathcal{D}_{Y}$-代数 $A_{Y}$ 与 $\mathcal{D}_{Y'}$-代数 $A_{Y'}$，后者是 $A_{Y}$ 到 $Y'$ 的限制。$A_{Y}$ 在原点 $0 \in Y$ 处的纤维等于 $A$。

下面解释一些细节。2.6.5 中 $A_{X}$ 的定义用到某个概形 $X^{\wedge}$。由于 $Y$ 不是字面意义上的曲线，$Y^{\wedge}$ 的定义应作如下修改。记 $\Delta_{n}$ 为对角线 $\Delta \subset \operatorname{Spec} O\widehat{\otimes}O$ 的第 $n$ 个无穷小邻域。态射 $\operatorname{Spec} O\widehat{\otimes}O \to \operatorname{Spec} O\otimes O = Y\times Y$ 诱导嵌入 $\Delta_{n} \hookrightarrow Y\times Y$（若 $n > 0$，则 $\Delta_{n}$ 比 $Y\times Y$ 中对角线 $\Delta$ 的第 $n$ 个无穷小邻域\emph{更小}）。现在，$Y^{\wedge}$ 的 $R$-点只应考虑满足下述性质的 $R$-态射 $\gamma: \operatorname{Spec} R\widehat{\otimes}O \to Y$：对任意 $n$ 存在 $N$，使得 $\gamma$ 诱导的态射 $\operatorname{Spec} O/t^{n}O \times \operatorname{Spec} O/t^{n}O \times \operatorname{Spec} R \to Y\times Y$ 经过 $\Delta_{N}$ 分解（此时可取 $N = 2n-2$）。

\bdp{3.7.2} 有时我们会用到截面
\begin{equation}\tag{94}
Y \to Y^{\wedge},
\end{equation}

\SourcePage{99}
它对应于由下式定义的态射 $\gamma: \operatorname{Spec} O\widehat{\otimes}O \to Y = \operatorname{Spec} O$：
\begin{equation}\tag{95}
\gamma^{*}(t) = t\otimes 1 + 1\otimes t.
\end{equation}
截面（94）给出同构
\begin{equation}\tag{96}
A_{Y} \xrightarrow{\sim} A\otimes O_{Y}.
\end{equation}
当然，（94）与（96）不是典范的：它们依赖于局部参数 $t \in O$ 的选取。

\bdp{3.7.3} 在 3.7.1 的情形中，考虑从 $\mathbb{C}$-代数到集合的函子 $F$，它使 $F(R)$ 为水平 $Y'$-态射 $\operatorname{Spec} R\widehat{\otimes}K \to \operatorname{Spec} A_{Y'}$ 的集合；等价地，为水平 $K$-态射 $H^{0}(Y', A_{Y'}) \to R\widehat{\otimes}K$ 的集合。$F$ 可由一个归纳仿射归纳概形 $S$ 表示（可称之为 $\operatorname{Spec} A_{Y'}$ 的水平截面归纳概形）。事实上，$F$ 是函子 $R \mapsto \operatorname{Hom}(V, R\widehat{\otimes}K)$ 的闭子函子，其中 $V = H^{0}(Y', A_{Y'})$，$\operatorname{Hom}$ 表示 $K$-线性映射的集合。

记 $A_{K}$ 为 $S$ 上正则函数环。这是完备拓扑代数（$A_{K}$ 中对应于 $S$ 的闭子概形的理想构成零的邻域基）。

$A_{K}$ 带群归纳概形 $\operatorname{Aut} K$ 的作用（$\operatorname{Aut} K$ 的 $R$-点是拓扑 $R$-代数 $R\widehat{\otimes}K$ 的自同构）。

$\operatorname{Spec} A_{Y}$ 的水平截面概形典范同构于 $\operatorname{Spec} A$（水平截面 $s: Y \to \operatorname{Spec} A_{Y}$ 对应到 $s(0) \in \operatorname{Spec} A$）。这是 $S = \operatorname{Spec} A_{K}$ 的闭子概形，于是得到典范满态射
\begin{equation}\tag{97}
A_{K} \to A.
\end{equation}
显然它是 $\operatorname{Aut} O$-等变的。

\emph{例.} 设 $A = \mathbb{C}[u_{0}, u_{1}, u_{2}, \ldots]$，$u_{k} = (L_{-1})^{k}u_{0}/k!$，$L_{0}u_{0} = du_{0}$，$d \in \mathbb{Z}$（按惯例 $L_{n} := -t^{n+1}\frac{d}{dt} \in \operatorname{Der} O$）。则 $\mathcal{D}_{Y}$-概形 $\operatorname{Spec} A_{Y}$ 与 $Y$ 上 $d$-微分的 jets 概形

\SourcePage{100}
之间存在显然的同构 $f$。显然 $\operatorname{Aut} O = \operatorname{Aut} Y$ 按函子性作用在这两个概形上。$f$ 关于由 $L_{0}$ 与 $L_{-1}$ 生成的 $\operatorname{Aut} O$ 群归纳概形等变。借助 $f$，我们把 $\operatorname{Spec} A_{Y'}$ 的水平截面与 $Y'$ 上的 $d$-微分等同，即与 $\omega_{Y'}^{\otimes d}$ 的截面等同。$Y'$ 上的 $d$-微分可写成 $\sum_{i}\tilde{u}_{i}t^{i}(dt)^{\otimes d}$，故 $A_{K} = \mathbb{C}[[\ldots, \tilde{u}_{-1}, \tilde{u}_{0}, \tilde{u}_{1}, \ldots]]$，其中
\begin{equation}\tag{98}
\mathbb{C}[[\ldots, \tilde{u}_{-1}, \tilde{u}_{0}, \tilde{u}_{1}, \ldots]] := \varprojlim_{n}\mathbb{C}[\ldots, \tilde{u}_{-1},, \tilde{u}_{0}, \tilde{u}_{1}, \ldots]/(u_{-n}, u_{-n-1}, \ldots).
\end{equation}
% [Tutor Note: 原书 p.100 式（98）中被取极限的环内印有重复逗号「\tilde{u}_{-1},, \tilde{u}_{0}」，且商理想的生成元印作 $u_{-n}$ 而非 $\tilde{u}_{-n}$；底本如此，保留原样。]
显然 $L_{0}\tilde{u}_{k} = (d+k)\tilde{u}_{k}$，$L_{-1}\tilde{u}_{k} = (k+1)\tilde{u}_{k+1}$，而态射（97）把 $\tilde{u}_{k}$ 映到 $u_{k}$（$k \ge 0$），映到 $0$（$k < 0$）。

\bdp{3.7.4} 记 $\mathfrak{z}_{\mathfrak{g}}(K)$ 为 3.7.3 中代数 $A_{K}$ 在 $A = \mathfrak{z}_{\mathfrak{g}}(O)$ 这一特殊情形的特例（$\mathfrak{z}_{\mathfrak{g}}(O)$ 的定义见 2.5.1 或 2.7.2）。我们将定义从 $\mathfrak{z}_{\mathfrak{g}}(K)$ 到完备化扭泛包络代数 $\overline{U}' = \overline{U}'(\mathfrak{g}\otimes K)$ 的中心 $\mathfrak{z}$ 的典范态射。为此把（34）改写成 $K$-线性映射 $\mathfrak{z}_{\mathfrak{g}}(O)\otimes_{O}K \to \mathfrak{z}\widehat{\otimes}K$。利用非典范同构 $\mathfrak{z}_{\mathfrak{g}}(O)_{Y} \xrightarrow{\sim} \mathfrak{z}_{\mathfrak{g}}(O)\otimes O_{Y}$（见（96）），得到映射
\begin{equation}\tag{99}
H^{0}(Y', \mathfrak{z}_{\mathfrak{g}}(O)_{Y'}) \to \mathfrak{z}\widehat{\otimes}K,
\end{equation}
容易看出它是典范的，即与局部参数 $t \in O$ 的选取无关。事实上，（34）是（99）的非典范版本；（34）依赖于 $t$ 的选取，因为（32）中出现 $\zeta + t$，而它正是由（95）定义的非典范截面 $Y' \to Y'^{\wedge}$。

\bdp{3.7.5} \emph{定理.}
\begin{enumerate}[label=(\roman*),leftmargin=2.2em]
  \item 映射（99）是 $K$-代数的水平态射。因此（99）定义了连续态射
  \begin{equation}\tag{100}
  \mathfrak{z}_{\mathfrak{g}}(K) \to \mathfrak{z}.
  \end{equation}
  \item 复合 $\mathfrak{z}_{\mathfrak{g}}(K) \to \mathfrak{z} \to \mathfrak{z}_{\mathfrak{g}}(O)$ 就是取 $A = \mathfrak{z}_{\mathfrak{g}}(O)$ 时的态射（97）。
  \item 态射（100）是 $\operatorname{Aut} K$-等变的。
\end{enumerate}

% ---- ch03_u05 ----
% ch03_u05 —— 原书 p.101-p.111（3.7.5 末句至 3.7.17 末；本单元不含节标题）
\SourcePage{101}

我们不证明这个定理。事实上，唯一不平凡的陈述是：（99）（等价地，（34））是环同态；证明见 ???。
% [Tutor Note: 原书 p.101 此处交叉引用印作「???」（引用未解析），且与下文之间缺空格；底本如此，保留原样。]

上述定理的自然途径，基于 VOA（即顶点算子代数）这一概念，或基于它的几何版本；后者在 [BD] 中以「手性代数」之名引入。\footnote{在 2.9.4 -- 2.9.5 中，我们用到了 VOA 理论（或手性代数理论）的一些想法。} 下一小节（读者可以跳过）概述手性代数的途径。

\bdp{3.7.6} 光滑曲线 $X$ 上的手性代数，是一个（左）$\mathcal{D}_X$-模 $\mathcal{A}$，带一个态射
\begin{equation}\tag{101}
j_*j^{!}(\mathcal{A}\boxtimes\mathcal{A})\to\Delta_*\mathcal{A},
\end{equation}
其中 $\Delta:X\hookrightarrow X\times X$ 是对角线，$j:(X\times X)\setminus\Delta(X)\hookrightarrow X$。
% [Tutor Note: 原书 p.101 此处把 $j$ 的靶写成 $X$；按定义应为 $(X\times X)\setminus\Delta(X)\to X\times X$，底本如此，保留原样。]

态射（101）应满足若干公理，这里不予陈述。若（101）把 $\mathcal{A}\boxtimes\mathcal{A}$ 映为 0，则称该手性代数是\emph{交换的}。此时（101）诱导一个态射 $\Delta_*(\mathcal{A}\otimes\mathcal{A})=j_*j^{!}(\mathcal{A}\boxtimes\mathcal{A})/\mathcal{A}\boxtimes\mathcal{A}\to\Delta_*\mathcal{A}$；换言之，诱导一个态射
\begin{equation}\tag{102}
\mathcal{A}\otimes\mathcal{A}\to\mathcal{A}.
\end{equation}

此时手性代数公理只要求：带运算（102）的 $\mathcal{A}$ 是交换结合含幺代数。于是，交换手性代数与 2.6 意义下的交换结合含幺 $\mathcal{D}_X$-代数相同。另一方面，由 2.6.5，对应于 $\operatorname{Aut}O$-模 $\mathit{Vac}'$ 的 $\mathcal{D}_X$-模 $\mathit{Vac}'_X$ 具有自然的手性代数结构（见下面的注）。由嵌入 $\mathfrak{z}_{\mathfrak{g}}(O)\to\mathit{Vac}'$ 诱导的映射 $\mathfrak{z}_{\mathfrak{g}}(O)_X\to\mathit{Vac}'_X$ 是手性代数态射。给定点 $x\in X$，可定义从手性代数到结合拓扑代数的函子 $\mathcal{A}\mapsto\mathcal{A}_{((x))}$。若对某个交换 $\operatorname{Aut}O$-代数 $A$ 有 $\mathcal{A}=A_X$，则 $\mathcal{A}_{((x))}$ 是 3.7.3 中的代数 $A_{K_x}$。若 $\mathcal{A}=\mathit{Vac}'_X$，则 $\mathcal{A}_{((x))}$ 是完备扭泛包络代数 $\overline{U}'=\overline{U}'(\mathfrak{g}\otimes K)$。于是

\SourcePage{102}

由函子性，得到态射 $\mathfrak{z}_{\mathfrak{g}}(K)=\mathfrak{z}_{\mathfrak{g}}(O)_K\to\overline{U}'$。其像含于 $\mathfrak{z}$ 中，因为 $\mathfrak{z}_{\mathfrak{g}}(O)_X$ 是手性代数 $\mathit{Vac}'_X$ 的\emph{中心}。

\emph{注.} 我们概述 $\mathit{Vac}'_X$ 上手性代数结构的定义。首先，对每个 $n$，在 $\operatorname{Sym}^nX$ 上构造一个 $\mathcal{D}$-模 $\mathit{Vac}'_{\operatorname{Sym}^nX}$（$n=1$ 时得到 $\mathit{Vac}'_X$）。$\mathit{Vac}'_{\operatorname{Sym}^nX}$ 在 $D\in\operatorname{Sym}^nX$ 处的纤维 $\mathit{Vac}'_D$ 可描述如下。把 $D$ 视为 $X$ 的 $n$ 阶闭子概形，记 $O_D$ 为 $X$ 沿 $D$ 的形式完备化上的函数环，并以 $\operatorname{Spec}K_D=(\operatorname{Spec}O_D)\setminus D$ 定义 $K_D$。与 $n=1$ 的情形一样，定义 $\mathfrak{g}\otimes K_D$ 的中心扩张 $\widetilde{\mathfrak{g}\otimes K_D}$。$\mathit{Vac}'_D$ 是对应于 Harish-Chandra 对 $(\widetilde{\mathfrak{g}\otimes K_D},G(O_D))$ 的扭真空模（见 1.2.5）。记 $\mathit{Vac}'_{X\times X}$ 为 $\mathit{Vac}'_{\operatorname{Sym}^2X}$ 到 $X\times X$ 的拉回。则
\begin{equation}\tag{103}
j^{!}\mathit{Vac}'_{X\times X}=j^{!}(\mathit{Vac}'_X\boxtimes\mathit{Vac}'_X),
\end{equation}
\begin{equation}\tag{104}
\Delta^{\dagger}\mathit{Vac}'_{X\times X}=\mathit{Vac}'_X
\end{equation}
其中 $j$ 与 $\Delta$ 的含义同（101），$\Delta^{\dagger}$ 表示朴素拉回，即 $\Delta^{\dagger}=H^1\Delta^{!}$。把（101）定义为复合
\[
j_*j^{!}\mathit{Vac}'_X\boxtimes\mathit{Vac}'_X=j_*j^{!}\mathit{Vac}'_{X\times X}\to j_*j^{!}\mathit{Vac}'_{X\times X}/\mathit{Vac}'_{X\times X}=\Delta_*\mathit{Vac}'_X
\]
最后一步等式来自（104）。

\bdp{3.7.7} \emph{定理.}
\begin{enumerate}[label=(\roman*),leftmargin=2.2em]
\item 态射（100）是拓扑同构。
\item $G(K)$ 在 $\mathfrak{z}$ 上的伴随作用是平凡的。
\end{enumerate}
证明将在 3.7.10 中给出。它基于 Feigin--Frenkel 定理，故下面两点是本质的：$\mathfrak{g}$ 半单；$\mathfrak{g}\otimes K$ 的中心扩张对应「临界」数量积（18）。对定理 3.7.5 而言，这两点并不本质。

我们还将证明以下结论。

\bdp{3.7.8} \emph{定理.} 2.9.8 中定义的映射 $\operatorname{gr}\mathfrak{z}\to\mathfrak{z}^{cl}$ 诱导拓扑同构
\[
\operatorname{gr}_i\mathfrak{z}\xrightarrow{\ \sim\ }\mathfrak{z}^{cl}_{(i)}:=\{\mathfrak{z}^{cl}\text{ 中 }i\text{ 次齐次多项式构成的空间}\}.
\]

\SourcePage{103}

\bdp{3.7.9} \emph{定理.} 记 $\mathcal{I}_n$ 为 $\overline{U}'$ 中由 $\mathfrak{g}\otimes t^nO$（$n\ge 0$）拓扑生成的闭左理想。则理想 $I_n:=\mathcal{I}_n\cap\mathfrak{z}\subset\mathfrak{z}$ 由空间 $\mathfrak{z}^m_i$（$m<i(1-n)$）拓扑生成；这里 $\mathfrak{z}^m_i:=\{z\in\mathfrak{z}_i\mid L_0z=mz\}$，$\mathfrak{z}_i$ 是 $\mathfrak{z}$ 的标准滤链，$L_0:=-t\frac{d}{dt}\in\operatorname{Der}O$。

\bdp{3.7.10} 我们来证明上述各定理。像（100）的元素是 $G(K)$-不变的（见 2.9.6 的注）。于是 3.7.7(ii) 由 3.7.7(i) 推出。下面证明 3.7.7(i)、3.7.8 与 3.7.9。

由 2.5.2，$\operatorname{gr}\mathfrak{z}_{\mathfrak{g}}(O)=\mathfrak{z}^{cl}_{\mathfrak{g}}(O)$。按 2.4.1，$\mathfrak{z}^{cl}_{\mathfrak{g}}(O)$ 可与 $\mathfrak{g}^*\otimes\omega_O$ 上 $G(O)$-不变多项式函数构成的环等同。选取 $\mathfrak{g}^*$ 上 $G$-不变多项式代数的齐次生成元 $p_1,\dots,p_r$，并置 $d_j:=\deg p_j$。如下定义 $v_{jk}\in\mathfrak{z}^{cl}_{\mathfrak{g}}(O)$（$1\le j\le r$，$0\le k<\infty$）：
\begin{equation}\tag{105}
p_j(\eta)=\sum_{k=0}^{\infty}v_{jk}(\eta)t^k(dt)^{d_j},\qquad \eta\in\mathfrak{g}^*\otimes\omega_O.
\end{equation}

按 2.4.1，代数 $\mathfrak{z}^{cl}_{\mathfrak{g}}(O)$ 由 $v_{jk}$ 自由生成。$\operatorname{Der}O$ 在 $\mathfrak{z}^{cl}_{\mathfrak{g}}(O)$ 上的作用容易描述。特别地 $v_{jk}=(L_{-1})^kv_{j0}/k!$，$L_0v_{j0}=d_jv_{j0}$。把 $v_{j0}\in\mathfrak{z}^{cl}_{\mathfrak{g}}(O)=\operatorname{gr}\mathfrak{z}_{\mathfrak{g}}(O)$ 提升为元素 $u_j\in\mathfrak{z}_{\mathfrak{g}}(O)$，使得 $L_0u_j=d_ju_j$。于是代数 $\mathfrak{z}_{\mathfrak{g}}(O)$ 由 $u_{jk}:=(L_{-1})^ku_j/k!$（$1\le j\le r$，$0\le k<\infty$）自由生成。与 3.7.3 末尾的例子一样，我们看到 $\mathfrak{z}_{\mathfrak{g}}(O)_K=\mathbb{C}[[\dots,\tilde{u}_{j,-1},\tilde{u}_{j0},\tilde{u}_{j1},\dots]]$，且 $L_0\tilde{u}_{jk}=(d_j+k)\tilde{u}_{jk}$。

记 $\bar{u}_{jk}$ 为 $\tilde{u}_{jk}$ 在 $\mathfrak{z}$ 中的像。由 2.9.8，$\bar{u}_{jk}\in\mathfrak{z}_{d_j}$；且 $\bar{u}_{jk}$ 在 $\mathfrak{z}^{cl}_{(d_j)}$ 中的像，是由
\begin{equation}\tag{106}
p_j(\eta)=\sum_k\tilde{v}_{jk}(\eta)t^k(dt)^{d_j},\qquad \eta\in\mathfrak{g}^*\otimes\omega_K
\end{equation}
定义的函数 $\tilde{v}_{jk}:\mathfrak{g}^*\otimes\omega_K\to\mathbb{C}$。
% [Tutor Note: 原书 p.103 对一般的 $k$ 仍写 $\bar{u}_{jk}\in\mathfrak{z}_{d_j}$；按 $L_0\bar{u}_{jk}=(d_j+k)\bar{u}_{jk}$，其滤链次数应为 $d_j+k$，底本如此，保留原样。]

我们有拓扑代数的同构
\begin{equation}\tag{107}
\mathfrak{z}^{cl}=\mathbb{C}[[\dots\tilde{v}_{j,-1},\tilde{v}_{j0},\tilde{v}_{j1},\dots]]
\end{equation}

\SourcePage{104}

因为
\begin{equation}\tag{108}
\parbox{0.76\linewidth}{\raggedright
$\mathfrak{g}^*\otimes t^{-n}\omega_O$ 上 $G(O)$-不变多项式函数构成的代数，由 $\tilde{v}_{jk}$（$k\ge -nd_j$）的限制自由生成；而当 $k<-nd_j$ 时，$\tilde{v}_{jk}$ 到 $\mathfrak{g}^*\otimes t^{-n}\omega_O$ 的限制等于 0。}
\end{equation}

（这一陈述可直接化归为 2.4.1 中考虑的 $n=0$ 情形。）定理 3.7.8 由（107）推出。

现在考虑态射 $f_n:\mathfrak{z}_{\mathfrak{g}}(O)_K\to\mathfrak{z}/I_n$，其中 $I_n$ 已在 3.7.9 中定义。我们将证明
\begin{equation}\tag{109}
\parbox{0.74\linewidth}{\raggedright
$f_n$ 是满射，且其核是由 $u_{jk}$（$k<d_j(1-n)$）拓扑生成的理想 $J_n$。}
\end{equation}
定理 3.7.7 与 3.7.9 由（109）推出。

为证明（109），考虑复合 $\bar{f}_n:\mathfrak{z}_{\mathfrak{g}}(O)_K\to\mathfrak{z}/I_n\hookrightarrow(\overline{U}'/\mathcal{I}_n)^{G(O)}$。给 $\overline{U}'/\mathcal{I}_n$ 配上由 $\overline{U}'$ 上标准滤链诱导的滤链。$L_0$ 在此滤链第 $i$ 项上的特征值 $\ge i(1-n)$。故 $\operatorname{Ker}\bar{f}_n\supset J_n$，其中 $J_n$ 已在（109）中定义。现在 $\operatorname{gr}(\overline{U}'/\mathcal{I}_n)^{G(O)}$ 含于 $(\operatorname{gr}\overline{U}'/\mathcal{I}_n)^{G(O)}$ 中，后者即 $\mathfrak{g}^*\otimes t^{-n}\omega_O$ 上 $G(O)$-不变多项式的代数。利用（108）容易看出：由 $\bar{f}_n$ 诱导的映射 $\mathfrak{z}_{\mathfrak{g}}(O)_K/J_n\to(\overline{U}'/\mathcal{I}_n)^{G(O)}$ 是同构。这蕴含（109）。我们还证明了
\begin{equation}\tag{110}
\parbox{0.55\linewidth}{\raggedright
映射 $\mathfrak{z}\to(\overline{U}'/\mathcal{I}_n)^{G(O)}$ 是满射}
\end{equation}
从而
\begin{equation}\tag{111}
\mathfrak{z}=(\overline{U}')^{G(O)}.
\end{equation}

\bdp{3.7.11} \emph{注}
\begin{enumerate}[label=(\roman*),leftmargin=2.2em]
\item 这里给出（111）的另一个证明。\footnote{它类似于下述事实的证明：$\mathfrak{g}\otimes K$ 的可积离散表示是平凡的。我们无法直接使用这一事实，因为 $\overline{U}'$ 不离散。} 设 $u\in(\overline{U}')^{G(O)}$。取 $h\in H(K)$，其中 $H\subset G$ 是固定的嘉当子群。则 $h^{-1}uh$ 关于某个 Borel 子群 $B_h\subset G$ 不变。
\SourcePage{105}
于是 $h^{-1}uh$ 是 $G$-不变的（只需对 $h^{-1}uh$ 在\emph{离散空间} $\overline{U}'/\mathcal{I}_n$ 中的像证明即可，其中 $\mathcal{I}_n$ 已在 3.7.9 中定义）。因此，对任意 $h\in H(K)$，$u$ 关于 $h\mathfrak{g}h^{-1}\subset\mathfrak{g}\otimes K$ 不变。李代数 $\mathfrak{g}\otimes K$ 由 $\mathfrak{g}\otimes O$ 与 $h\mathfrak{g}h^{-1}$（$h\in H(K)$）生成。故 $u\in\mathfrak{z}$。
\item 事实上，对任意开子代数 $\mathfrak{a}\subset\mathfrak{g}\otimes K$，都有
\begin{equation}\tag{112}
\mathfrak{z}=(\overline{U}')^{\mathfrak{a}}.
\end{equation}
确实，可按如下方式修改上述证明。把 $u$ 写成 $u_\chi$ 的（无限）和，其中 $\chi\in\mathfrak{h}^*:=(\operatorname{Lie}H)^*$，且对 $a\in\mathfrak{h}$ 有 $[a,u_\chi]=\chi(a)u_\chi$。然后取 $h\in H(K)$，使得 $h$ 在 $H(K)/H(O)=\{\text{余权格}\}$ 中的像，关于包含 $\mathfrak{h}$ 的 Borel 子代数 $\mathfrak{b}\subset\mathfrak{g}$ 是「非常占优」的，从而 $h^{-1}\mathfrak{a}h\supset[\mathfrak{b},\mathfrak{b}]$。我们看到：除非 $\chi$ 占优，$u_\chi=0$；且 $h^{-1}u_0h$ 是 $\mathfrak{g}$-不变的。用 $h^{-1}$ 替换 $h$，可见 $u=u_0$，等等。
\item 这里给出 3.7.7(ii) 的另一个证明。考虑 $\overline{U}'$ 的典范滤链 $\overline{U}'_k$。由（109）可知，空间 $\overline{U}'_k\cap\mathfrak{z}$（$k\in\mathbb{N}$）之并在 $\mathfrak{z}$ 中稠密。故只需证明：对每个 $k$，$G(K)$ 在 $\overline{U}'_k\cap\mathfrak{z}$ 上的作用是平凡的。$G(K)$ 在 $\mathfrak{z}^{cl}$ 上的作用是平凡的（见（107）、（106））。故 $G(K)$ 在 $\operatorname{gr}\mathfrak{z}$ 上的作用是平凡的。$\mathfrak{g}\otimes K$ 在 $\widetilde{\mathfrak{g}\otimes K}$ 上对应于由（19）定义的 $G(K)$ 作用的那个作用，就是伴随作用；而 $\mathfrak{g}\otimes K$ 在 $\mathfrak{z}$ 上的伴随作用是平凡的。故 $G(K)$ 在 $\mathfrak{z}$ 上的作用经由 $\pi_0(G(K))$ 分解。群 $\pi_0(G(K))$ 有限（见 4.5.4），证毕。
\end{enumerate}

\bdp{3.7.12} 我们将从 [FF92] 推出定理 3.6.7。记 $A_{{}^{L}\mathfrak{g}}(O)$ 为 $\operatorname{Op}_{{}^{L}\mathfrak{g}}(O)$ 的坐标环（即 $\operatorname{Spec}O$ 上 ${}^{L}\mathfrak{g}$-oper 构成的概形）。设 $\varphi_O:A_{{}^{L}\mathfrak{g}}(O)\xrightarrow{\ \sim\ }\mathfrak{z}_{\mathfrak{g}}(O)$ 是满足 3.2.2 中条件的同构。它诱导 $\operatorname{Aut}K$-等变同构 $\varphi_K:A_{{}^{L}\mathfrak{g}}(K)\xrightarrow{\ \sim\ }\mathfrak{z}_{\mathfrak{g}}(K)$，其中 $A_{{}^{L}\mathfrak{g}}(K)$ 是 3.7.3 中对应于 $A=A_{{}^{L}\mathfrak{g}}(O)$ 的代数 $A_K$。

\SourcePage{106}

回忆：$A_K$ 是 $\operatorname{Spec}A_{Y'}$ 的水平截面构成的归纳概形的坐标环，$Y':=\operatorname{Spec}K$。若 $A=A_{{}^{L}\mathfrak{g}}(O)$，则 $\operatorname{Spec}A_{Y'}$ 是 $Y'$ 上 ${}^{L}\mathfrak{g}$-oper 的 jets 概形，其水平截面是 $Y'$ 上的 ${}^{L}\mathfrak{g}$-oper（参看 3.3.3）。于是 $A_{{}^{L}\mathfrak{g}}(K)$ 是 $\operatorname{Op}_{{}^{L}\mathfrak{g}}(K)$ 的坐标环；后者定义为 $\operatorname{Spec}K$ 上 ${}^{L}\mathfrak{g}$-oper 构成的归纳概形。它是关于 \emph{Gelfand--Dikii 括号}的泊松代数（我们在 3.7.14 中回忆其定义）。Gelfand--Dikii 括号依赖于 ${}^{L}\mathfrak{g}$ 上非退化不变双线性形式的选取。我们把它取为 $\mathfrak{g}$ 上形式（18）的\emph{对偶}：即，它在 $\mathfrak{h}^*={}^{L}\mathfrak{h}\subset{}^{L}\mathfrak{g}$ 上的限制，对偶于（18）在 $\mathfrak{h}$ 上的限制。

由 3.7.5 与 3.7.7，有典范同构 $\mathfrak{z}_{\mathfrak{g}}(K)\xrightarrow{\ \sim\ }\mathfrak{z}$，故可把 $\varphi_K$ 视为 $\operatorname{Aut}K$-等变同构
\begin{equation}\tag{113}
A_{{}^{L}\mathfrak{g}}(K)\xrightarrow{\ \sim\ }\mathfrak{z}.
\end{equation}
$\mathfrak{z}$ 是关于 Hayashi 括号的泊松代数（见 3.6.2）。

\bdp{3.7.13} \emph{定理.} [FF92] 存在同构
\begin{equation}\tag{114}
\varphi_O:A_{{}^{L}\mathfrak{g}}(O)\xrightarrow{\ \sim\ }\mathfrak{z}_{\mathfrak{g}}(O)
\end{equation}
满足 3.2.2 的条件，且相应的同构（113）与泊松结构相容。

我们将证明（见 3.7.16）：具备定理中所述性质的同构（114）满足 3.6.7 的条件。故它是唯一的（见 3.6.7 的注）。

\emph{注.} 如 3.6.12 所解释，可对 $\mathfrak{g}$ 上任一不变双线性形式 $B$ 配上 $\mathfrak{z}$ 上的一个泊松括号（3.6.2 的括号对应形式（18））。若 $B$ 非退化，则可考虑 ${}^{L}\mathfrak{g}$ 上的对偶形式，以及 $A_{{}^{L}\mathfrak{g}}(K)$ 上相应的 Gelfand--Dikii 括号。对应于（114）的同构（113）与这些泊松括号相容。

\SourcePage{107}

\bdp{3.7.14} 我们回忆 [DS85] 中 Gelfand--Dikii 括号的定义。这是 $\operatorname{Op}_{\mathfrak{g}}(K)$ 上的泊松括号（即其坐标环 $A_{\mathfrak{g}}(K)$ 上的泊松括号）。它依赖于 $\mathfrak{g}$ 上非退化不变双线性形式 $(\,,\,)$ 的选取。

记 $\widetilde{\mathfrak{g}\otimes K}$ 为 $\mathfrak{g}\otimes K$ 对应于 $(\,,\,)$ 的 Kac--Moody 中心扩张。作为向量空间，$\widetilde{\mathfrak{g}\otimes K}=(\mathfrak{g}\otimes K)\oplus\mathbb{C}$；其换位子由 2-余循环 $\operatorname{Res}(du,v)$（$u,v\in\mathfrak{g}\otimes K$）定义。拓扑对偶空间 $(\widetilde{\mathfrak{g}\otimes K})^*$ 是归纳概形。$(\widetilde{\mathfrak{g}\otimes K})^*$ 上正则函数的代数，是关于 Kirillov 括号\footnote{如 [We83] 所解释，「Kirillov 括号」由 Sophus Lie 发明，后来又被包括 A.A. Kirillov 在内的若干人重新发现。}的泊松代数（即唯一的连续泊松括号，使得从 $\widetilde{\mathfrak{g}\otimes K}$ 到 $(\widetilde{\mathfrak{g}\otimes K})^*$ 上正则函数代数的自然映射是李代数态射）。故 $(\widetilde{\mathfrak{g}\otimes K})^*$ 是泊松「流形」。记 $(\widetilde{\mathfrak{g}\otimes K})^*_1$ 为连续线性泛函 $l:\widetilde{\mathfrak{g}\otimes K}\to\mathbb{C}$ 构成的空间，其中 $l$ 在中心 $\mathbb{C}\subset\widetilde{\mathfrak{g}\otimes K}$ 上的限制是恒等。$(\widetilde{\mathfrak{g}\otimes K})^*_1$ 是 $(\widetilde{\mathfrak{g}\otimes K})^*$ 的泊松子流形。

我们把 $(\widetilde{\mathfrak{g}\otimes K})^*_1$ 与 $\operatorname{Conn}$ 等同，后者定义为 $\operatorname{Spec}K$ 上平凡 $G$-丛的联络构成的归纳概形：对联络 $\nabla=d+\eta$（$\eta\in\mathfrak{g}\otimes\omega_K$），配以 $l\in(\widetilde{\mathfrak{g}\otimes K})^*_1$，使得对任意 $u\in\mathfrak{g}\otimes K\subset\widetilde{\mathfrak{g}\otimes K}$ 有 $l(u)=\operatorname{Res}(u,\eta)$。容易验证：$\mathfrak{g}\otimes K$ 在 $\operatorname{Conn}$ 上的规范作用，对应于 $\mathfrak{g}\otimes K$ 在 $(\widetilde{\mathfrak{g}\otimes K})^*_1$ 上的余伴随作用；而 $G(K)$ 在 $(\widetilde{\mathfrak{g}\otimes K})^*$ 上的余伴随作用\footnote{它是对偶于由（19）定义的 $G(K)$ 在 $\widetilde{\mathfrak{g}\otimes K}$ 上的伴随作用（当然，（19）中的 $c$ 应换成我们 $\mathfrak{g}$ 上的双线性形式）。}，定义为：其在 $(\widetilde{\mathfrak{g}\otimes K})^*_1$ 上的限制对应于 $G(K)$ 在 $\operatorname{Conn}$ 上的规范作用。$G(K)$ 在泊松「流形」$(\widetilde{\mathfrak{g}\otimes K})^*_1$ 上的作用按字面意义不是哈密顿的，即：不能定义矩映射 $(\widetilde{\mathfrak{g}\otimes K})^*_1\to(\mathfrak{g}\otimes K)^*$。然而可以定义矩映射 $(\widetilde{\mathfrak{g}\otimes K})^*_1\to(\widetilde{\mathfrak{g}\otimes K})^*$：它就是恒等映射。

\SourcePage{108}

要点在于：$\operatorname{Op}_{\mathfrak{g}}(K)$ 可由 $\operatorname{Conn}=(\widetilde{\mathfrak{g}\otimes K})^*_1$ 经哈密顿约化得到（这种对 $\operatorname{Op}_{\mathfrak{g}}(K)$ 的解释，自动定义 $A_{\mathfrak{g}}(K)$ 上的泊松括号）。固定 Borel 子群 $B\subset G_{ad}$。设 $N$ 是其幂幺根，$\mathfrak{n}:=\operatorname{Lie}N$。由于 Kac--Moody 余循环在 $\mathfrak{n}\otimes K$ 上的限制平凡，有显然的分裂 $\mathfrak{n}\otimes K\to\widetilde{\mathfrak{g}\otimes K}$。它是 $B(K)$-等变的，且此性质唯一刻画它。$N(K)$ 在 $\operatorname{Conn}$ 上的作用是哈密顿作用：矩映射 $\mu:\operatorname{Conn}=(\widetilde{\mathfrak{g}\otimes K})^*_1\to(\mathfrak{n}\otimes K)^*$ 由上述分裂诱导。设 $\operatorname{Char}^{\star}\subset(\mathfrak{n}\otimes K)^*$ 是\emph{非退化特征标}的集合，即满足下述条件的李代数态射 $l:\mathfrak{n}\otimes K\to\mathbb{C}$ 之集：对每个单根 $\alpha$，$l$ 在 $\mathfrak{g}^\alpha\otimes K$ 上的限制非零。对每个 $l\in\operatorname{Char}^{\star}$，$N(K)$ 在 $\mu^{-1}(l)$ 上的作用自由，且商 $N(K)\backslash\mu^{-1}(l)$ 可典范地等同于 $\operatorname{Op}_{\mathfrak{g}}(K)$：确实，$\mu^{-1}(l)$ 由满足下述条件的联络 $\nabla=d+\eta\in\operatorname{Conn}$ 构成：$\eta=\sum_{\alpha\in\Gamma}J_\alpha+q$，其中 $q\in\mathfrak{b}\otimes\omega_K$，$\Gamma$ 是单根集，$J_\alpha=J_\alpha(l)$ 是 $\mathfrak{g}^{-\alpha}\otimes\omega_K$ 中固定的非零元。于是，$\operatorname{Op}_{\mathfrak{g}}(K)$ 由 $\operatorname{Conn}$ 对 $N(K)$ 的作用作关于 $l$ 的哈密顿约化得到，由此得到 $\operatorname{Op}_{\mathfrak{g}}(K)$ 上的泊松括号。它称为\emph{Gelfand--Dikii 括号}。它不依赖于 $l$。确实，对 $l,l'\in\operatorname{Char}^{\star}$，考虑同构
\begin{equation}\tag{115}
N(K)\backslash\mu^{-1}(l)\xrightarrow{\ \sim\ }N(K)\backslash\mu^{-1}(l')
\end{equation}
它来自（115）两端与 $\operatorname{Op}_{\mathfrak{g}}(K)$ 的等同。$H(K)$ 在 $\operatorname{Conn}=(\widetilde{\mathfrak{g}\otimes K})^*_1$ 上的（余）伴随作用保持相关结构（即 $\operatorname{Conn}$ 上的泊松括号、$N(K)$ 在 $\operatorname{Conn}$ 上的作用，以及矩映射 $\mu:\operatorname{Conn}\to(\mathfrak{n}\otimes K)^*$）。存在唯一的 $h\in H(K)$ 把 $l$ 变为 $l'$，且（115）由该 $h$ 的作用诱导。故（115）是泊松映射。

\emph{注}
\begin{enumerate}[label=(\roman*),leftmargin=2.2em]
\item 若 $\mathfrak{g}$ 上的双线性形式 $(\,,\,)$ 乘以 $c\in\mathbb{C}^*$，则 $\operatorname{Op}_{\mathfrak{g}}(K)$ 上的泊松括号乘以 $c^{-1}$。
\SourcePage{109}
\item 上面定义的 Gelfand--Dikii 括号是「第二 Gelfand--Dikii 括号」。第一个括号的定义，以及它与 Gelfand--Dikii 原始工作（[GD76]、[GD78]）之关系的说明，可在 [DS85] 中找到（见同上书第 2.3、3.6、3.7、6.5、8 节）。
\end{enumerate}

\bdp{3.7.15} 设 $\mathfrak{F}\in\operatorname{Op}_{\mathfrak{g}}(K)$，即 $\mathfrak{F}=(\mathfrak{F}_B,\nabla)$，其中 $\mathfrak{F}_B$ 是 $\operatorname{Spec}K$ 上的 $B$-丛，$\nabla$ 是相应 $G$-丛上满足 3.1.3 条件的联络（这里 $G$ 是对应于 $\mathfrak{g}$ 的伴随群，$B\subset G$ 是 Borel 子群）。我们将描述切空间 $T_{\mathfrak{F}}\operatorname{Op}_{\mathfrak{g}}(K)$ 与余切空间 $T^*_{\mathfrak{F}}\operatorname{Op}_{\mathfrak{g}}(K)$。然后写出 $\{\varphi,\psi\}(\mathfrak{F})$（$\varphi,\psi\in A_{\mathfrak{g}}(K)$）的显式公式。

\emph{注.} 当然 $\mathfrak{F}_B$ 总是平凡的，故可把 $\mathfrak{F}$ 视为平凡 $G$-丛中的联络 $\nabla$（即 $\nabla=d+q$，$q\in\mathfrak{g}\otimes\omega_K$），模去关于 $B$ 的规范变换。

为描述 $T_{\mathfrak{F}}\operatorname{Op}_{\mathfrak{g}}(K)$，须研究 $\mathfrak{F}=(\mathfrak{F}_B,\nabla)$ 的无穷小形变。由于 $\mathfrak{F}_B$ 不能形变，所有这些形变都来自 $\nabla$ 的无穷小形变，其形如 $\nabla(\varepsilon)=\nabla+\varepsilon q$，$q\in H^0(\operatorname{Spec}K,\mathfrak{g}^{-1}_{\mathfrak{F}}\otimes\omega_K)$（$\mathfrak{g}^{-1}$ 的定义见 3.1.1；$\mathfrak{g}^{-1}_{\mathfrak{F}}:=\mathfrak{g}^{-1}_{\mathfrak{F}_B}$ 是 $\mathfrak{g}^{-1}$ 的 $\mathfrak{F}_B$-扭）。计入 $\mathfrak{F}_B$ 的无穷小自同构，得到：
\begin{equation}\tag{116}
T_{\mathfrak{F}}\operatorname{Op}_{\mathfrak{g}}(K)=H^0(\operatorname{Spec}K,\operatorname{Coker}(\nabla:\mathfrak{b}_{\mathfrak{F}}\to\mathfrak{g}^{-1}_{\mathfrak{F}}\otimes\omega_K)).
\end{equation}
这里有切空间更方便的描述：
\begin{equation}\tag{117}
T_{\mathfrak{F}}\operatorname{Op}_{\mathfrak{g}}(K)=\operatorname{Coker}(\nabla:\mathfrak{n}^K_{\mathfrak{F}}\to\mathfrak{b}^K_{\mathfrak{F}}\otimes\omega_K)
\end{equation}
其中 $\mathfrak{n}^K_{\mathfrak{F}}:=H^0(\operatorname{Spec}K,\mathfrak{n}_{\mathfrak{F}})$，$\mathfrak{b}^K_{\mathfrak{F}}:=H^0(\operatorname{Spec}K,\mathfrak{b}_{\mathfrak{F}})$（从（117）右端到（116）右端的自然映射是同构）。利用 $\mathfrak{g}$ 上的不变数量积 $(\,,\,)$，把 $\mathfrak{b}^*$、$\mathfrak{n}^*$ 与 $\mathfrak{g}/\mathfrak{n}$、$\mathfrak{g}/\mathfrak{b}$ 等同，得到余切空间的如下描述：
\begin{equation}\tag{118}
T^*_{\mathfrak{F}}\operatorname{Op}_{\mathfrak{g}}(K)=\{u\in\mathfrak{g}^K_{\mathfrak{F}}\mid\nabla(u)\in\mathfrak{b}^K_{\mathfrak{F}}\otimes\omega_K\}/\mathfrak{n}^K_{\mathfrak{F}}.
\end{equation}

\SourcePage{110}

Gelfand--Dikii 括号的显式公式如下：
\begin{equation}\tag{119}
\{\varphi,\psi\}(\mathfrak{F})=\operatorname{Res}(\nabla(d_{\mathfrak{F}}\varphi),d_{\mathfrak{F}}\psi),\qquad \varphi,\psi\in A_{\mathfrak{g}}(K).
\end{equation}
此公式中，微分 $d_{\mathfrak{F}}\varphi$ 与 $d_{\mathfrak{F}}\psi$ 视为（118）右端的元素。

\bdp{3.7.16} \emph{定理.}\footnote{受 [Phys] 启发。}
\begin{enumerate}[label=(\roman*),leftmargin=2.2em]
\item 置 $I:=\operatorname{Ker}(A_{\mathfrak{g}}(K)\to A_{\mathfrak{g}}(O))$。则 $\{I,I\}\subset I$，故 $I/I^2$ 是 $A_{\mathfrak{g}}(O)$ 上的李代数胚。
\item 存在李代数胚的 $\operatorname{Aut}O$-等变拓扑同构
\begin{equation}\tag{120}
I/I^2\xrightarrow{\ \sim\ }\mathfrak{a}_{\mathfrak{g}}
\end{equation}
（$\mathfrak{a}_{\mathfrak{g}}$ 的定义见 3.5.11、3.5.15）。
\end{enumerate}
（此定理中，$I^2$ 表示由 $ab$（$a\in I$，$b\in I$）生成的子空间的\emph{闭包}。）

定理 3.6.7 由 3.7.13 与 3.7.16 推出。

\emph{注.} 同构（120）是唯一的（见 3.5.13 或 3.5.14）。

\bdp{3.7.17} 我们来证明定理 3.7.16。沿用 3.7.15 的记号。取 $\mathfrak{F}\in\operatorname{Op}_{\mathfrak{g}}(O)$。与（117）类似，$T_{\mathfrak{F}}\operatorname{Op}_{\mathfrak{g}}(O)$ 有如下描述：
\begin{equation}\tag{121}
T_{\mathfrak{F}}\operatorname{Op}_{\mathfrak{g}}(O)=\operatorname{Coker}(\nabla:\mathfrak{n}^O_{\mathfrak{F}}\to\mathfrak{b}^O_{\mathfrak{F}}\otimes\omega_O)
\end{equation}
其中 $\mathfrak{n}^O_{\mathfrak{F}}:=H^0(\operatorname{Spec}O,\mathfrak{n}_{\mathfrak{F}})$。$I/I^2$ 在 $\mathfrak{F}$ 处的纤维是余法空间 $T^{\perp}_{\mathfrak{F}}\operatorname{Op}_{\mathfrak{g}}(O)\subset T^*_{\mathfrak{F}}\operatorname{Op}_{\mathfrak{g}}(K)$。按（121），它用（118）的语言有如下描述：
\begin{equation}\tag{122}
T^{\perp}_{\mathfrak{F}}\operatorname{Op}_{\mathfrak{g}}(O)=\{u\in\mathfrak{g}^O_{\mathfrak{F}}\mid\nabla(u)\in\mathfrak{b}^O_{\mathfrak{F}}\otimes\omega_O\}/\mathfrak{n}^O_{\mathfrak{F}}.
\end{equation}

现在 $\{I,I\}\subset I$ 是清楚的：若 $\varphi,\psi\in I$，$\mathfrak{F}\in\operatorname{Op}_{\mathfrak{g}}(O)$，则 $d_{\mathfrak{F}}\varphi$ 与 $d_{\mathfrak{F}}\psi$ 属于（122）的右端，从而（119）的右端等于 0。

\SourcePage{111}

映射
\begin{equation}\tag{123}
I/I^2\to\operatorname{Der}A_{\mathfrak{g}}(O),
\end{equation}
是 $I/I^2$ 上代数胚结构的一部分，由 $\varphi\mapsto\partial_\varphi$ 定义：$\partial_\varphi(\psi):=\{\varphi,\psi\}$，$\varphi\in I$，$\psi\in A_{\mathfrak{g}}(K)/I=A_{\mathfrak{g}}(O)$。于是按（119），由（123）诱导的映射
\begin{equation}\tag{124}
T^{\perp}_{\mathfrak{F}}\operatorname{Op}_{\mathfrak{g}}(O)\to T_{\mathfrak{F}}\operatorname{Op}_{\mathfrak{g}}(O)
\end{equation}
是算子
\begin{equation}\tag{125}
\nabla:\{u\in\mathfrak{g}^O_{\mathfrak{F}}\mid\nabla(u)\in\mathfrak{b}^O_{\mathfrak{F}}\otimes\omega_O\}/\mathfrak{n}^O_{\mathfrak{F}}\to(\mathfrak{b}^O_{\mathfrak{F}}\otimes\omega_O)/\nabla(\mathfrak{n}^O_{\mathfrak{F}}).
\end{equation}

$I/I^2$ 上的代数胚结构，在映射（124）的核 $\mathfrak{a}_{\mathfrak{F}}$ 上诱导李代数结构。另一方面，$\mathfrak{a}_{\mathfrak{F}}$ 是（125）的核，即 $\mathfrak{a}_{\mathfrak{F}}=\{u\in\mathfrak{g}^O_{\mathfrak{F}}\mid\nabla(u)=0\}/\{u\in\mathfrak{n}^O_{\mathfrak{F}}\mid\nabla(u)=0\}$。由于 $\{u\in\mathfrak{n}^O_{\mathfrak{F}}\mid\nabla(u)=0\}=0$，有
\begin{equation}\tag{126}
\mathfrak{a}_{\mathfrak{F}}=\{u\in\mathfrak{g}^O_{\mathfrak{F}}\mid\nabla(u)=0\}.
\end{equation}
（126）的右端是 $\mathfrak{g}^O_{\mathfrak{F}}$ 的李子代数。

\emph{引理.} 由 $I/I^2$ 上代数胚结构给出的 $\mathfrak{a}_{\mathfrak{F}}$ 上的李代数结构，与由（126）诱导者一致。

\emph{证明.} 只需证明：若 $\varphi_1,\varphi_2\in A_{\mathfrak{g}}(K)$ 且 $d_{\mathfrak{F}}\varphi_i\in\mathfrak{a}_{\mathfrak{F}}$，则
\begin{equation}\tag{127}
d_{\mathfrak{F}}\{\varphi_1,\varphi_2\}=[d_{\mathfrak{F}}\varphi_1,d_{\mathfrak{F}}\varphi_2]
\end{equation}
（在（127）的右端，$d_{\mathfrak{F}}\varphi_i$ 经（126）视为 $\mathfrak{g}^O_{\mathfrak{F}}$ 的元素）。考虑 $\mathfrak{F}$ 的形变 $\mathfrak{F}(\varepsilon)$，$\varepsilon^2=0$。把 $\mathfrak{F}$ 写成 $(\mathfrak{F}_B,\nabla)$。不失一般性，可设 $\mathfrak{F}(\varepsilon)=(\mathfrak{F}_B,\nabla+\varepsilon q)$，$q\in\mathfrak{b}^K_{\mathfrak{F}}\otimes\omega_K$。把 $d_{\mathfrak{F}(\varepsilon)}\varphi_i$ 写成 $d_{\mathfrak{F}}\varphi_i+\varepsilon\mu_i$。则
\[
\begin{gathered}
\{\varphi_1,\varphi_2\}(\mathfrak{F}(\varepsilon))=\operatorname{Res}((\nabla+\varepsilon\operatorname{ad}q)(d_{\mathfrak{F}}\varphi_1+\varepsilon\mu_1),d_{\mathfrak{F}}\varphi_2+\varepsilon\mu_2)=\\
\varepsilon\operatorname{Res}([q,d_{\mathfrak{F}}\varphi_1],d_{\mathfrak{F}}\varphi_2)=\varepsilon\operatorname{Res}(q,[d_{\mathfrak{F}}\varphi_1,d_{\mathfrak{F}}\varphi_2])
\end{gathered}
\]
（我们用到了 $\nabla(d_{\mathfrak{F}}\varphi_i)=0$）。（127）的等式随之成立。\hfill$\Box$

% ---- ch03_u06 ----
% ch03_u06 —— 原书 p.112–p.121（3.7 末的无编号段落、3.8 全部 3.8.1–3.8.23）
% p.112 首段无段号，接 p.111 末 3.7.17 证明之后；3.8.23 正文止于 p.121 末，p.122 起为第 4 章。
\SourcePage{112}

根据该引理，映射（124）的核 $\mathfrak{a}_{\mathfrak{F}}$ 作为李代数与 $(\mathfrak{g}_{\mathrm{univ}})_{\mathfrak{F}}$ 一致，后者是 3.5.11 中李代数 $\mathfrak{g}_{\mathrm{univ}}$ 在 $\mathfrak{F}$ 处的纤维。映射（124）=（125）是满射，因为 $\nabla:\mathfrak{g}^{O}_{\mathfrak{F}}\to\mathfrak{g}^{O}_{\mathfrak{F}}\otimes\omega_{O}$ 是满射。容易证明，（121）与（122）是同胚，且映射（124）是开映射。

类似地可以证明，态射（123）是满射且开的，它的核可典范地等同于带离散拓扑的 $\mathfrak{g}_{\mathrm{univ}}$（这一等同对每个 $\mathfrak{F}\in\operatorname{Op}_{\mathfrak{g}}(O)$ 诱导出上述同构 $\mathfrak{a}_{\mathfrak{F}}\xrightarrow{\sim}(\mathfrak{g}_{\mathrm{univ}})_{\mathfrak{F}}$）。引理 3.5.12 给出连续的李代数胚态射 $f:I/I^{2}\to\mathfrak{a}_{\mathfrak{g}}$，使得图表
\[
\begin{tikzcd}
  0 \arrow[r] & \mathfrak{g}_{\mathrm{univ}} \arrow[r] \arrow[d, "\mathrm{id}"'] & I/I^{2} \arrow[r] \arrow[d, "f"']
    & \operatorname{Der}A_{\mathfrak{g}}(O) \arrow[r] \arrow[d, "\mathrm{id}"] & 0 \\
  0 \arrow[r] & \mathfrak{g}_{\mathrm{univ}} \arrow[r] & \mathfrak{a}_{\mathfrak{g}} \arrow[r]
    & \operatorname{Der}A_{\mathfrak{g}}(O) \arrow[r] & 0
\end{tikzcd}
\]
交换。由于图表的两行在拓扑意义下正合，$f$ 是拓扑同构。显然 $f$ 是 $\operatorname{Aut}O$-等变的。

\bdsec{3.8. oper 的奇点}

\bdp{3.8.1} 设 $U$ 是曲线 $X$ 的开稠子集。我们要把归纳概形 $\operatorname{Op}_{\mathfrak{g}}(U)$ 表示成若干闭子概形 $\operatorname{Op}_{\mathfrak{g},D}(X)$ 的并，其中 $D$ 跑遍 $X$ 的有限子概形，且满足 $D\cap U=\emptyset$。

由 3.1.9，有典范同构 $\underline{\operatorname{Op}}_{\mathfrak{g}}(U)\xrightarrow{\sim}\operatorname{Op}_{\mathfrak{g}}(U)$。这里 $\underline{\operatorname{Op}}_{\mathfrak{g}}(U)$ 是由 $\Gamma(U,\omega_{X}^{\otimes 2})$-torsor $\operatorname{Op}_{\mathfrak{sl}_{2}}(U)$ 经某个嵌入 $\Gamma(U,\omega_{X}^{\otimes 2})\subset\Gamma(U,V_{\omega_{X}})$ 诱导出的 $\Gamma(U,V_{\omega_{X}})$-torsor。该嵌入与 $V=V_{\mathfrak{g}}$ 的定义见 3.1.9。提醒一句：$V$ 是带 $\mathbb{G}_{m}$-作用（即分次）的向量空间，$V_{\omega_{X}}$ 表示用 $\mathbb{G}_{m}$-torsor $\omega_{X}$ 扭 $V$ 所得的扭。有 $\dim V=r:=\operatorname{rank}\mathfrak{g}$，且 $V$ 各分次分量的次数等于 $\mathfrak{g}$ 上「基本」不变多项式的次数 $d_{1},\ldots,d_{r}$。

\SourcePage{113}

若 $D$ 是 $X$ 的有限子概形，则有典范嵌入 $V_{\omega_{X}}\hookrightarrow V_{\omega_{X}(D)}$。记 $\overline{\operatorname{Op}}_{\mathfrak{g},D}(X)$ 为由 $\Gamma(X,V_{\omega_{X}})$-torsor $\underline{\operatorname{Op}}_{\mathfrak{g}}(X)$ 诱导出的 $\Gamma(X,V_{\omega_{X}(D)})$-torsor。显然 $\overline{\operatorname{Op}}_{\mathfrak{g},D}(X)$ 是 $\underline{\operatorname{Op}}_{\mathfrak{g}}(X\setminus D)$ 的闭子概形。记 $\operatorname{Op}_{\mathfrak{g},D}(X)$ 为 $\overline{\operatorname{Op}}_{\mathfrak{g},D}(X)$ 在 $\operatorname{Op}_{\mathfrak{g}}(X\setminus D)$ 中的像。若 $D\subset D'$，则 $\operatorname{Op}_{\mathfrak{g},D}(X)\subset\operatorname{Op}_{\mathfrak{g},D'}(X)$。对任意开稠子集 $U\subset X$，有
\[
\operatorname{Op}_{\mathfrak{g}}(U)=\bigcup_{D\cap U=\emptyset}\operatorname{Op}_{\mathfrak{g},D}(X).
\]

3.8.23 中将给出 $\operatorname{Op}_{\mathfrak{g},D}(X)$ 的「内蕴」描述，它不使用同构 $\underline{\operatorname{Op}}_{\mathfrak{g}}\xrightarrow{\sim}\operatorname{Op}_{\mathfrak{g}}$。该描述的局部版本见 3.8.7--3.8.10。

\bdp{3.8.2} 现在可以给出 2.8.6 中问题的答案：
\begin{equation}\tag{128}
  N_{\Delta}(G)=\operatorname{Op}_{{}^{L}\mathfrak{g},\Delta}(X).
\end{equation}
$N_{\Delta}(G)$ 定义为归纳概形 $N'_{\Delta}(G)$ 的子概形，后者经 Feigin--Frenkel 同构典范地等同于 $\operatorname{Op}_{{}^{L}\mathfrak{g}}(X\setminus\Delta)$。（128）是 $\operatorname{Op}_{{}^{L}\mathfrak{g}}(X\setminus\Delta)$ 的子概形之间的等式。

我们不证明（128）。3.8.6 中将给出一个提示。

\bdp{3.8.3} 3.8.1 中 $\operatorname{Op}_{\mathfrak{g},D}(X)$ 的定义在如下局部情形有意义：
\[
X=\operatorname{Spec}O,\qquad O:=\mathbb{C}[[t]],\qquad D=\operatorname{Spec}O/t^{n}O.
\]
此时记 $\operatorname{Op}_{\mathfrak{g},n}(O)$ 以代替 $\operatorname{Op}_{\mathfrak{g},D}(X)$。$\operatorname{Op}_{\mathfrak{g},n}(O)$ 是归纳概形 $\operatorname{Op}_{\mathfrak{g}}(K)$ 的闭子概形。当然 $\operatorname{Op}_{\mathfrak{g},0}(O)=\operatorname{Op}_{\mathfrak{g}}(O)$，$\operatorname{Op}_{\mathfrak{g},n}(O)\subset\operatorname{Op}_{\mathfrak{g},n+1}(O)$，且 $\operatorname{Op}_{\mathfrak{g}}(K)$ 是诸 $\operatorname{Op}_{\mathfrak{g},n}(O)$ 的归纳极限。

由 3.7.12，$A_{\mathfrak{g}}(K)$ 是 $\operatorname{Op}_{\mathfrak{g}}(K)$ 上正则函数构成的代数。记 $I_{n}$ 为 $A_{\mathfrak{g}}(K)$ 中对应于 $\operatorname{Op}_{\mathfrak{g},n}(O)\subset\operatorname{Op}_{\mathfrak{g}}(K)$ 的理想。显然 $I_{n}\supset I_{n+1}$，且 $I_{0}$ 就是 3.7.16(i) 中的理想 $I$。诸理想 $I_{n}$ 构成 $A_{\mathfrak{g}}(K)$ 中 $0$ 的邻域基。

\bdp{3.8.4} 下面给出 $A_{\mathfrak{g}}(K)$ 与 $I_{n}$ 的显式描述。沿用 3.5.6 的记号，于是 $\operatorname{Spec}K$ 上的 $\mathfrak{g}$-oper 与满足 $u_{j}(t)\in\mathbb{C}((t))$ 的算子（64）一一对应。把 $u_{j}(t)$ 写成 $\sum_{k}\tilde{u}_{jk}t^{k}$。则 $A_{\mathfrak{g}}(K)=\mathbb{C}[[\ldots,\tilde{u}_{j,-1},\tilde{u}_{j0},\tilde{u}_{j1},\ldots]]$（用记号（98））。理想 $I_{n}$ 由诸 $\tilde{u}_{jk}$（$k<-d_{j}n$）拓扑生成。\SourcePage{114} 3.5.6 中的 $u_{jk}$ 是 $\tilde{u}_{jk}$ 在 $A_{\mathfrak{g}}(O)=A_{\mathfrak{g}}(K)/I$ 中的像。

容易描述 $\operatorname{Der}K$ 在 $A_{\mathfrak{g}}(K)$ 上的作用。特别地
\begin{equation}\tag{129}
  L_{0}\tilde{u}_{jk}=(d_{j}+k)\tilde{u}_{jk}.
\end{equation}

与整体情形一样（见 3.1.12--3.1.14），$\operatorname{Op}_{\mathfrak{g}}(K)$ 的坐标环 $A_{\mathfrak{g}}(K)$ 带典范滤链。其第 $i$ 项由 $\tilde{u}_{jk}$ 的「多项式」组成，这些多项式的加权次数 $\le i$；这里约定 $\tilde{u}_{jk}$ 的权为 $d_{j}$。

\bdp{3.8.5} \emph{命题.} 理想 $I_{n}\subset A_{\mathfrak{g}}(K)$ 由诸空间 $A_{i}^{m}$（$m<i(1-n)$）拓扑生成，其中 $A_{i}^{m}$ 是 $A_{\mathfrak{g}}(K)$ 的滤链第 $i$ 项中满足 $L_{0}a=ma$ 的元素 $a$ 的集合。\hfill$\Box$

同构 $A_{{}^{L}\mathfrak{g}}(K)\xrightarrow{\sim}\mathfrak{z}$（见（113））保持滤链，且是 $\operatorname{Aut}K$-等变的。于是命题 3.8.5 推出如下陈述。

\bdp{3.8.6} \emph{命题.} Feigin--Frenkel 同构 $A_{{}^{L}\mathfrak{g}}(K)\xrightarrow{\sim}\mathfrak{z}$ 把 $I_{n}\subset A_{{}^{L}\mathfrak{g}}(K)$ 映到 3.7.9 中的理想 $I_{n}$。

这是（128）的证明中用到的组成部分之一。

\bdp{3.8.7} 我们要用「自然」的语言（不使用同构（43））描述 $\operatorname{Op}_{\mathfrak{g},n}(O)$。记 $\mathfrak{g}^{+}$ 为 $\mathfrak{g}$ 的如下局部闭既约子概形：它由所有满足下述条件的 $a\in\mathfrak{g}$ 组成——对正根 $\alpha$，若 $\alpha$ 非单则 $a_{-\alpha}=0$，若 $\alpha$ 单则 $a_{-\alpha}\ne 0$（$a_{-\alpha}$ 是 $a$ 在根子空间 $\mathfrak{g}^{-\alpha}$ 中的分量）。于是对任意 $\mathbb{C}$-代数 $R$，集合 $\mathfrak{g}^{+}(R)$ 由 $a\in\mathfrak{g}\otimes R$ 组成，使得对每个非单的 $\alpha>0$ 有 $a_{-\alpha}=0$，且对每个单的 $\alpha$，$a_{-\alpha}$ 生成 $R$-模 $\mathfrak{g}^{-\alpha}\otimes R$。

回忆：$\operatorname{Spec}K$ 上的 $\mathfrak{g}$-oper 是算子 $\frac{d}{dt}+q(t)$（$q\in\mathfrak{g}^{+}(K)$）的 $B(K)$-共轭类。这里 $B$ 是 $\mathfrak{g}$ 所对应的伴随群 $G$ 的 Borel 子群。

\SourcePage{115}

\bdp{3.8.8} \emph{定义.} $\operatorname{Spec}O$ 上的 $(\le n)$-奇异 $\mathfrak{g}$-oper 是算子 $\frac{d}{dt}+t^{-n}q(t)$（$q\in\mathfrak{g}^{+}(O)$）的 $B(O)$-共轭类。

\emph{注.}
\begin{enumerate}[label=(\roman*),leftmargin=2.2em]
\item $B(O)$ 在算子集 $\{\frac{d}{dt}+t^{-n}q(t)|q\in\mathfrak{g}^{+}(O)\}$ 上的作用是自由的。事实上，$B(K)$ 在 $\{\frac{d}{dt}+q(t)|q\in\mathfrak{g}^{+}(K)\}$ 上的作用是自由的（见 3.1.4）。
\item $n=0$ 时得到 $\operatorname{Spec}O$ 上 $\mathfrak{g}$-oper 的通常概念。
\end{enumerate}

\bdp{3.8.9} \emph{命题.} 映射「$\operatorname{Spec}O$ 上的 $(\le n)$-奇异 $\mathfrak{g}$-oper」$\to\operatorname{Op}_{\mathfrak{g}}(K)$ 是单射。其像等于 $\operatorname{Op}_{\mathfrak{g},n}(O)$。

\emph{证明.} 沿用 3.5.6 的记号。对每个 $v_{1},\ldots,v_{r}\in\mathbb{C}[[t]]$，算子
\begin{equation}\tag{130}
  \frac{d}{dt}+t^{-n}(i(f)+v_{1}(t)e_{1}+\cdots+v_{r}(t)e_{r})
\end{equation}
定义了 $\operatorname{Spec}O$ 上一个 $(\le n)$-奇异 $\mathfrak{g}$-oper。容易证明，这在算子（130）与 $\operatorname{Spec}O$ 上的 $(\le n)$-奇异 $\mathfrak{g}$-oper 之间给出双射。现在用 $B(K)$-共轭把（130）化成「典范形式」（64）。用 $t^{-n\check\rho}$ 共轭（130），得到
\begin{equation}\tag{131}
  \frac{d}{dt}+i(f)+n\check\rho t^{-1}+t^{-nd_{1}}v_{1}(t)e_{1}+\cdots+t^{-nd_{r}}v_{r}(t)e_{r}.
\end{equation}
% [Tutor Note: 原书 p.115 把 $\exp(-ne_{1}/2t)$ 与 $n(n-2)/4t^{2}$ 印作线性分数，
%  字面亦可读作 $(ne_{1}/2)t$、$(n(n-2)/4)t^{2}$；由本段共轭计算（$n=1$ 时应得 $-1/4t^{2}$）
%  可知应为 $e_{1}/(2t)$、$n(n-2)/(4t^{2})$。底本如此，正文照录。]
为消去 $n\check\rho t^{-1}$，再用 $\exp(-ne_{1}/2t)$ 共轭（131），得到算子（64），其中
\[
\begin{aligned}
  u_{j}(t)&=t^{-nd_{j}}v_{j}(t) \qquad \text{对 } j>1,\\
  u_{1}(t)&=t^{-nd_{1}}v_{1}(t)+n(n-2)/4t^{2},\qquad d_{1}=2.
\end{aligned}
\]
显然 $v_{j}\in\mathbb{C}[[t]]$ 当且仅当 $u_{j}\in t^{-nd_{j}}\mathbb{C}[[t]]$。\hfill$\Box$

\bdp{3.8.10} 若把 $\operatorname{Op}_{\mathfrak{g},n}(O)$ 的点看作 $\operatorname{Spec}O$ 上的 $(\le n)$-奇异 $\mathfrak{g}$-oper，则典范嵌入 $\operatorname{Op}_{\mathfrak{g},n}(O)\hookrightarrow\operatorname{Op}_{\mathfrak{g},n+1}(O)$ 把 $\frac{d}{dt}+t^{-n}q(t)$（$q\in\mathfrak{g}^{+}(O)$）的 $B(O)$-共轭类映到 $t^{\check\rho}(\frac{d}{dt}+t^{-n}q(t))t^{-\check\rho}$ 的 $B(O)$-共轭类（这是良定义的，因为 $t^{\check\rho}B(O)t^{-\check\rho}\subset B(O)$）。

\SourcePage{116}

\bdp{3.8.11} 记 $\operatorname{Inv}(\mathfrak{g})$ 为 $\mathfrak{g}$ 上 $G$-不变多项式构成的代数。有典范态射 $\mathfrak{g}\to\operatorname{Spec}\operatorname{Inv}(\mathfrak{g})=W\backslash\mathfrak{h}$，其中 $W$ 是 Weyl 群。

设有 $\operatorname{Spec}O$ 上的 $(\le 1)$-奇异 $\mathfrak{g}$-oper，即算子 $\frac{d}{dt}+t^{-1}q(t)$（$q\in\mathfrak{g}^{+}(O)$）的 $B(O)$-共轭类。$q(0)\in\mathfrak{g}$ 在 $\operatorname{Spec}\operatorname{Inv}(\mathfrak{g})$ 中的像称为该 oper 的\emph{留数}。于是我们定义了\emph{留数映射}
\begin{equation}\tag{132}
  \operatorname{Res}:\operatorname{Op}_{\mathfrak{g},1}(O)\to\operatorname{Spec}\operatorname{Inv}(\mathfrak{g})=W\backslash\mathfrak{h}.
\end{equation}
它是满射。因此它诱导嵌入
\begin{equation}\tag{133}
  \operatorname{Inv}(\mathfrak{g})\hookrightarrow A_{\mathfrak{g}}(K)/I_{1}
\end{equation}
（回忆：$A_{\mathfrak{g}}(K)/I_{1}$ 是 $\operatorname{Op}_{\mathfrak{g},1}(O)$ 的坐标环；见 3.8.3）。

\bdp{3.8.12} \emph{命题.} $\operatorname{Res}(\operatorname{Op}_{\mathfrak{g}}(O))\subset W\backslash\mathfrak{h}$ 由单点组成，该点是 $-\check\rho\in\mathfrak{h}$ 的像。

\emph{注.} 我们宁愿忘记 $-\check\rho$ 与 $\check\rho$ 在 $W\backslash\mathfrak{h}$ 中有相同的像。

\emph{证明.} 需要计算 3.8.10 所述映射 $\operatorname{Op}_{\mathfrak{g}}(O)\to\operatorname{Op}_{\mathfrak{g},1}(O)$ 与映射（132）的复合。若 $q(t)\in\mathfrak{g}^{+}(O)$，则 $t^{\check\rho}(\frac{d}{dt}+q(t))t^{-\check\rho}=\frac{d}{dt}+\frac{a-\check\rho}{t}+\{\text{正则项}\}$，其中 $a$ 属于对应于单负根的根子空间之和。现在 $a-\check\rho$ 与 $-\check\rho$ 在 $W\backslash\mathfrak{h}$ 中有相同的像。\hfill$\Box$

\bdp{3.8.13} \emph{命题.} 设 $f\in A_{\mathfrak{g}}(K)/I_{1}$，即 $f$ 是 $\operatorname{Op}_{\mathfrak{g},1}(O)$ 上的正则函数。则以下条件等价：
\begin{enumerate}[label=(\roman*),leftmargin=2.2em]
\item $f\in\operatorname{Inv}(\mathfrak{g})$，其中 $\operatorname{Inv}(\mathfrak{g})$ 按（133）与其像等同；
\item $f$ 是 $\operatorname{Aut}^{0}O$-不变的；
\item $L_{0}f=0$。
\end{enumerate}

\emph{证明.} 显然 (i)$\Rightarrow$(ii)$\Rightarrow$(iii)。下面从 (ii) 推出 (i)。考虑 $\operatorname{Spec}O$ 上的一个 $(\le 1)$-奇异 $\mathfrak{g}$-oper。它是联络 $\frac{d}{dt}+t^{-1}q(t)$（$q\in\mathfrak{g}^{+}(O)$）的 $B(O)$-共轭类。若把 $t$ 换成 $\lambda t$，该联络换成 $\frac{d}{dt}+t^{-1}q(\lambda t)$。由于 $L_{0}f=0$，$f$ 在该联络上的值\SourcePage{117}不依赖于 $\lambda$，因而只依赖于 $q(0)\in\mathfrak{g}^{+}$（因为 $\lim_{\lambda\to 0}q(\lambda t)=q(0)$）。剩下的只需用如下事实：$\mathfrak{g}^{+}$ 上 $B$-不变的正则函数可延拓为 $\mathfrak{g}$ 上 $G$-不变的多项式（见 [Ko63] 的定理 0.10）。\hfill$\Box$

\bdp{3.8.14} \emph{注.} 由 3.8.4，代数 $A_{\mathfrak{g}}(K)/I_{1}$ 由 $\tilde{u}_{jk}$（$k\ge -d_{j}$）自由生成，其中 $\tilde{u}_{jk}\in A_{\mathfrak{g}}(K)/I_{1}$ 是 $\tilde{u}_{jk}\in A_{\mathfrak{g}}(K)$ 的像。由 3.8.13 与（129），$\operatorname{Inv}(\mathfrak{g})\subset A_{\mathfrak{g}}(K)/I_{1}$ 由 $v_{j}:=\tilde{u}_{j,-d_{j}}$ 生成。同构 $\operatorname{Spec}\mathbb{C}[v_{1},\ldots,v_{r}]\xrightarrow{\sim}\operatorname{Spec}\operatorname{Inv}(\mathfrak{g})$ 是复合 $\operatorname{Spec}\mathbb{C}[v_{1},\ldots,v_{r}]\to\mathfrak{g}\to\operatorname{Spec}\operatorname{Inv}(\mathfrak{g})$，其中第一个映射等于 $i(f)-\check\rho+v_{1}e_{1}+\cdots+v_{r}e_{r}$（用 3.5.6 的记号）。

\bdp{3.8.15} 我们将证明定理 3.6.11。在 3.8.16--3.8.17 中，将叙述 Feigin--Frenkel 同构（113）的一个性质。该性质把定理 3.6.11 归结为某个陈述（见 3.8.19），后者只涉及 oper 与 Gelfand--Dikii 括号。该陈述将在 3.8.20--3.8.22 中证明。

\bdp{3.8.16} 沿用 3.5.17 的记号。此外，若 $\operatorname{Der}O$ 作用在向量空间 $V$ 上，则令 $V^{0}:=\{v\in V|L_{0}v=0\}$。

如 3.6.9 所解释，3.6.8 中的映射 $\pi$ 诱导态射
\begin{equation}\tag{134}
  (\mathfrak{z}/\mathfrak{z}\cdot\mathfrak{z}^{<0})^{0}=(\mathfrak{z}/\mathfrak{z}\cdot\mathfrak{z}^{<0})^{\le 0}
  =\mathfrak{z}^{\le 0}/(\mathfrak{z}\cdot\mathfrak{z}^{<0}\cap\mathfrak{z}^{\le 0})\to C
\end{equation}
其中 $C$ 是 $U\mathfrak{g}$ 的中心。现在（113）诱导同构
\begin{equation}\tag{135}
  (\mathfrak{z}/\mathfrak{z}\cdot\mathfrak{z}^{<0})^{0}\xrightarrow{\sim}(A_{{}^{L}\mathfrak{g}}(K)/I_{1})^{0},
\end{equation}
因为由 3.8.5，$I_{1}=A_{{}^{L}\mathfrak{g}}(K)\cdot A_{{}^{L}\mathfrak{g}}(K)^{<0}$。由 3.8.13，（135）的右端等于 $\operatorname{Inv}({}^{L}\mathfrak{g})$。于是（134）与（135）给出态射
\begin{equation}\tag{136}
  \operatorname{Inv}({}^{L}\mathfrak{g})\to C.
\end{equation}

记 $\operatorname{Inv}(\mathfrak{h}^{*})$ 为 $\mathfrak{h}^{*}$ 上 $W$-不变多项式构成的代数。由于 ${}^{L}\mathfrak{h}=\mathfrak{h}^{*}$，有典范同构 $\operatorname{Inv}({}^{L}\mathfrak{g})\xrightarrow{\sim}\operatorname{Inv}(\mathfrak{h}^{*})$。此外还有\SourcePage{118} Harish-Chandra 同构 $C\xrightarrow{\sim}\operatorname{Inv}(\mathfrak{h}^{*})$。于是（136）可看作映射
\begin{equation}\tag{137}
  \operatorname{Inv}(\mathfrak{h}^{*})\to\operatorname{Inv}(\mathfrak{h}^{*}).
\end{equation}

\bdp{3.8.17} \emph{定理.}（E. Frenkel，私人通信）

态射（137）把 $f\in\operatorname{Inv}(\mathfrak{h}^{*})$ 映到 $f^{-}$，其中 $f^{-}(\lambda):=f(-\lambda)$，$\lambda\in\mathfrak{h}^{*}$。\hfill$\Box$

\bdp{3.8.18} 利用 3.8.17，可以把图表（84）神秘的左下角换成其 oper 版本。下面的图表（143）本质上就是这样得到的。我们来定义（143）的下方箭头，它是 3.6.9--3.6.10 中所构造映射（83）的 oper 版本。

由 3.8.5，
\begin{equation}\tag{138}
  I_{1}=A_{\mathfrak{g}}(K)\cdot A_{\mathfrak{g}}(K)^{<0}.
\end{equation}

由 3.8.13，有典范同构
\begin{equation}\tag{139}
  (A_{\mathfrak{g}}(K)/I_{1})^{0}\xrightarrow{\sim}\operatorname{Inv}(\mathfrak{g}).
\end{equation}

对 $h\in\mathfrak{h}$，记 $m_{h}$ 为 $\operatorname{Inv}(\mathfrak{g})$ 中由在 $h$ 处取零值的多项式构成的极大理想。令 $m:=m_{-\check\rho}$。由 3.8.12，同构（139）诱导
\begin{equation}\tag{140}
  (I/I_{1})^{0}\xrightarrow{\sim}m.
\end{equation}

现在得到
\begin{equation}\tag{141}
  (I/(I^{2}+I_{1}))^{0}\xrightarrow{\sim}m/m^{2}
\end{equation}
（为从（140）得到（141），我们用到
\[
(I^{2})^{0}\subset(I^{0})^{2}+I\cdot I^{<0}\subset(I^{0})^{2}+A_{\mathfrak{g}}(K)\cdot A_{\mathfrak{g}}(K)^{<0}=(I^{0})^{2}+I_{1};
\]
见（138））。

对正则的 $h\in\mathfrak{h}$，我们把 $m_{h}/m_{h}^{2}$ 与 $\mathfrak{h}^{*}$ 等同：把 $\mathfrak{h}$ 上的 $W$-不变多项式对应到它在 $h$ 处的微分。特别地，对 $m=m_{-\check\rho}$，\SourcePage{119}有 $m/m^{2}\xrightarrow{\sim}\mathfrak{h}^{*}$（顺便说一句，若把 $m$ 写成 $m_{\check\rho}$，就会得到另一个同构 $m/m^{2}\xrightarrow{\sim}\mathfrak{h}^{*}$）。

最后，利用（138）把（141）的左端改写成用 $I/I^{2}$ 表达的式子，得到同构
\begin{equation}\tag{142}
  (I/I^{2})^{\le 0}/(A_{\mathfrak{g}}(O)\cdot(I/I^{2})^{<0}\cap(I/I^{2})^{\le 0})\xrightarrow{\sim}\mathfrak{h}^{*}.
\end{equation}

\bdp{3.8.19} \emph{命题.} 图表
\begin{equation}\tag{143}
\begin{tikzcd}[row sep=1.5em]
  (\mathfrak{a}_{\mathfrak{g}})^{\le 0}/(A_{\mathfrak{g}}(O)\cdot\mathfrak{a}_{\mathfrak{g}}^{<0}\cap(\mathfrak{a}_{\mathfrak{g}})^{\le 0})
    \arrow[r, "\sim"] & \mathfrak{h} \arrow[d, "\wr"] \\
  (I/I^{2})^{\le 0}/(A_{\mathfrak{g}}(O)\cdot(I/I^{2})^{<0}\cap(I/I^{2})^{\le 0})
    \arrow[r, "\sim"] \arrow[u, "\wr"] & \mathfrak{h}^{*}
\end{tikzcd}
\end{equation}
交换。这里下方箭头是同构（142），上方箭头是同构（78），左边箭头由同构（120）诱导（后者来自 $A_{\mathfrak{g}}(K)$ 上的 Gelfand--Dikii 括号），右边箭头由定义 Gelfand--Dikii 括号时所用 $\mathfrak{g}$ 上的不变数量积诱导。

该命题将在 3.8.20--3.8.22 中证明。

定理 3.6.11 由 3.8.17 与 3.8.19 推出。（143）的交换性蕴含（84）的反交换性，因为如下图表反交换：
\[
\begin{tikzcd}
  m_{\check\rho}/(m_{\check\rho})^{2} \arrow[rr, "\sim"] \arrow[dr] & & m_{-\check\rho}/(m_{-\check\rho})^{2} \arrow[dl] \\
  & \mathfrak{h}^{*} &
\end{tikzcd}
\]
这里上方箭头由 3.8.17 中的映射 $f\mapsto f^{-}$ 诱导。

\bdp{3.8.20} 我们要叙述命题 3.8.19 的证明中用到的引理。考虑复合
\begin{equation}\tag{144}
  I/I^{2}\to I/(I^{2}+I_{1})\xrightarrow{\sim}\mathfrak{a}_{\mathfrak{g}}/A_{\mathfrak{g}}(O)\cdot\mathfrak{a}_{\mathfrak{g}}^{<0}
  =\mathfrak{a}_{\mathfrak{g}}/\mathfrak{a}_{\mathfrak{n}}=\mathfrak{g}_{\mathrm{univ}}/\mathfrak{n}_{\mathrm{univ}}.
\end{equation}
这里第二个箭头来自（120）与（138）；$\mathfrak{a}_{\mathfrak{n}}$ 与 $\mathfrak{n}_{\mathrm{univ}}$ 已在 3.5.16 中定义，$\mathfrak{a}_{\mathfrak{g}}$ 已在 3.5.11 中定义；等式 $\mathfrak{a}_{\mathfrak{n}}=A_{\mathfrak{g}}(O)\cdot\mathfrak{a}_{\mathfrak{g}}^{<0}$\SourcePage{120} 已在 3.5.18 中证明。$I/I^{2}$ 在 $\mathfrak{F}=(\mathfrak{F}_{B},\nabla)\in\operatorname{Op}_{\mathfrak{g}}(O)$ 上的纤维等于 $\{u\in\mathfrak{g}^{O}_{\mathfrak{F}}|\nabla(u)\in\mathfrak{b}^{O}_{\mathfrak{F}}\otimes\omega_{O}\}/\mathfrak{n}^{O}_{\mathfrak{F}}$（见（122）），而 $\mathfrak{g}_{\mathrm{univ}}/\mathfrak{n}_{\mathrm{univ}}$ 在 $\mathfrak{F}$ 上的纤维等于 $(\mathfrak{g}_{\mathfrak{F}}/\mathfrak{n}_{\mathfrak{F}})_{0}$，即 $\mathfrak{g}_{\mathfrak{F}}/\mathfrak{n}_{\mathfrak{F}}$ 在原点 $0\in\operatorname{Spec}O$ 处的纤维。考虑映射
\[
\varphi,\psi:\{u\in\mathfrak{g}^{O}_{\mathfrak{F}}|\nabla(u)\in\mathfrak{b}^{O}_{\mathfrak{F}}\otimes\omega_{O}\}/\mathfrak{n}^{O}_{\mathfrak{F}}
  \to(\mathfrak{g}_{\mathfrak{F}}/\mathfrak{n}_{\mathfrak{F}})_{0}
\]
其中 $\varphi$ 由（144）诱导，$\psi$ 是在 $0$ 处取值。

\bdp{3.8.21} \emph{引理.} $\varphi=\psi$。

\emph{证明.} 由 3.7.17，$\varphi$ 与 $\psi$ 在 $\mathfrak{a}_{\mathfrak{F}}:=\{u\in\mathfrak{g}^{O}_{\mathfrak{F}}|\nabla(u)=0\}$ 上的限制相等。因此只需证明 $\operatorname{Ker}\varphi\subset\operatorname{Ker}\psi$。显然 $\operatorname{Ker}\varphi=T^{\perp}_{\mathfrak{F}}\operatorname{Op}_{\mathfrak{g},1}(O)$，即 $\operatorname{Op}_{\mathfrak{g},1}(O)$ 在 $\mathfrak{F}$ 处的余法空间。对任意 $q\in\mathfrak{b}^{O}_{\mathfrak{F}}$，oper $\mathfrak{F}_{q}:=(\mathfrak{F}_{B},\nabla+q\cdot\frac{d}{dt})$ 是 $(\le 1)$-奇异的。于是 $\mathfrak{b}^{O}_{\mathfrak{F}}\otimes t^{-1}\omega_{O}$ 在（117）右端的像含于切空间 $T_{\mathfrak{F}}\operatorname{Op}_{\mathfrak{g},1}(O)$ 中。因此 $T^{\perp}_{\mathfrak{F}}\operatorname{Op}_{\mathfrak{g},1}(O)\subset\operatorname{Ker}\psi$。\hfill$\Box$

\bdp{3.8.22} 现在证明 3.8.19。由于（142）的左端等于（141）的左端，可以把 3.8.19 改述如下。

设 $f\in\operatorname{Inv}(\mathfrak{g})$，$f(-\check\rho)=0$。把 $f$ 看作 $A_{\mathfrak{g}}(K)/I_{1}$ 的元素（见（133））。由 3.8.12，$f\in I/I_{1}$。$f$ 在 $I/(I^{2}+I_{1})$ 中的像可看作元素 $\nu\in\mathfrak{g}_{\mathrm{univ}}/\mathfrak{n}_{\mathrm{univ}}$（见（144））。另一方面，设 $\lambda\in\mathfrak{h}^{*}$ 是 $f\in\operatorname{Inv}(\mathfrak{g})$ 在 $\mathfrak{h}$ 上的限制在 $-\check\rho$ 处的微分。为证明 3.8.19，必须证明 $\nu$ 等于 $\lambda$ 在复合
\[
\mathfrak{h}^{*}\xrightarrow{\sim}\mathfrak{h}\subset\mathfrak{h}\otimes A_{\mathfrak{g}}(O)
  =\mathfrak{h}_{\mathrm{univ}}/\mathfrak{n}_{\mathrm{univ}}\subset\mathfrak{g}_{\mathrm{univ}}/\mathfrak{n}_{\mathrm{univ}}
\]
下的像。

由 3.8.21，这等价于如下陈述：设 $\mathfrak{F}=(\mathfrak{F}_{B},\nabla)\in\operatorname{Op}_{\mathfrak{g}}(O)$，$q\in\mathfrak{b}^{O}_{\mathfrak{F}}$，$\mathfrak{F}_{\varepsilon q}:=(\mathfrak{F}_{B},\nabla+\varepsilon q\frac{d}{dt})$，则
\begin{equation}\tag{145}
  \frac{d}{d\varepsilon}f(\operatorname{Res}(\mathfrak{F}_{\varepsilon q}))|_{\varepsilon=0}=\lambda(q_{\mathfrak{h}}(0)),
\end{equation}
其中 $q_{\mathfrak{h}}(t)\in\mathfrak{h}[[t]]$ 是 $q$ 在 $\mathfrak{b}^{O}_{\mathfrak{F}}/\mathfrak{n}^{O}_{\mathfrak{F}}=\mathfrak{h}\otimes O$ 中的像。与 3.8.12 的证明一样可以证明，$\operatorname{Res}(\mathfrak{F}_{\varepsilon q})$ 等于 $-\check\rho+\varepsilon q_{\mathfrak{h}}(0)$ 在 $W\backslash\mathfrak{h}$ 中的像。于是（145）是显然的。

\SourcePage{121}

\bdp{3.8.23} 本小节（读者完全可以跳过）给出 3.8.1 中概形 $\operatorname{Op}_{\mathfrak{g},D}(X)$ 的「内蕴」描述。它由对 3.8.7--3.8.10 作直接的「整体化」得到。

记 $G$ 为 $\mathfrak{g}$ 所对应的伴随群。设我们处于 3.1.2 的情形。于是有 $X$ 上的 $B$-丛 $\mathfrak{F}_{B}$、诱导出的 $G$-丛 $\mathfrak{F}_{G}$，以及 $\mathfrak{g}_{\mathfrak{F}}\otimes\omega_{X}$-torsor $\operatorname{Conn}(\mathfrak{F}_{G})$。设 $D$ 是 $X$ 的有限子概形。记 $\operatorname{Conn}_{D}(\mathfrak{F}_{G})$ 为由 $\operatorname{Conn}(\mathfrak{F}_{G})$ 诱导出的 $\mathfrak{g}_{\mathfrak{F}}\otimes\omega_{X}(D)$-torsor；于是 $\operatorname{Conn}_{D}(\mathfrak{F}_{G})$ 的截面是带 $(\le D)$-奇点的联络。与 3.1.2 一样定义 $c:\operatorname{Conn}_{D}(\mathfrak{F}_{G})\to(\mathfrak{g}/\mathfrak{b})_{\mathfrak{F}}\otimes\omega_{X}(D)$。$X$ 上 $(\le D)$-奇异 $\mathfrak{g}$-oper 的概念如下定义：把定义 3.1.3 中的 $\operatorname{Conn}$ 换成 $\operatorname{Conn}_{D}$，$\omega_{X}$ 换成 $\omega_{X}(D)$。

若 $X$ 完备，则 $X$ 上的 $(\le D)$-奇异 $\mathfrak{g}$-oper 构成概形。与 3.8.9 一样可以证明，从该概形到 $\operatorname{Op}_{\mathfrak{g}}(X\setminus D)$ 的自然态射是闭嵌入，且其像等于 $\operatorname{Op}_{\mathfrak{g},D}(X)$。于是可把 $\operatorname{Op}_{\mathfrak{g},D}(X)$ 看作 $X$ 上 $(\le D)$-奇异 $\mathfrak{g}$-oper 的模概形。

若 $D\subset D'$，则 $\operatorname{Op}_{\mathfrak{g},D}(X)\subset\operatorname{Op}_{\mathfrak{g},D'}(X)$。所以我们应当有自然的办法，从 $(\le D)$-奇异 $\mathfrak{g}$-oper $(\mathfrak{F}_{B},\nabla)$ 构造 $(\le D')$-奇异 $\mathfrak{g}$-oper $(\mathfrak{F}'_{B},\nabla')$。当然 $(\mathfrak{F}'_{B},\nabla')$ 应配有同构 $\alpha:(\mathfrak{F}'_{B},\nabla')|_{X\setminus\Delta}\xrightarrow{\sim}(\mathfrak{F}_{B},\nabla)|_{X\setminus\Delta}$，其中 $\Delta\subset X$ 是把 $D$、$D'$、$\Delta$ 看作有效除子时满足 $D'=D+\Delta$ 的有限子概形。联络 $\nabla'$ 由 $\nabla$ 与 $\alpha$ 重构，而 $(\mathfrak{F}'_{B},\alpha)$ 由下述性质定义（参看 3.8.10）：若 $x\in\Delta$，$f$ 是 $\Delta$ 在 $x$ 处的局部方程，$s$ 是 $\mathfrak{F}_{B}$ 在 $x$ 处的局部截面，则存在 $\mathfrak{F}'_{B}$ 在 $x$ 处的局部截面 $s'$，使得 $\alpha(s')=\lambda(f)s$，其中 $\lambda:\mathbb{G}_{m}\to H$ 是对应于 $\check\rho$ 的态射。
